%#! latex main.tex
\chapter{Analytic aspects of $p$-adic periods.}\label{chap-ana-p-period}
\addtocontents{toc}{\protect\par\vskip2mm\hskip20mm(Analysis)\par\vskip5mm}
\minitoc
\newpage
\section{Analytic continuation; topological point of view.}\label{sec-1}

\begin{abst}
 Thanks to Berkovich's presentation of $p$-adic analytic geometry, it is
 possible to make sense of the familiar monodromy principle in the
 exotic world of $p$-adic manifolds. Its application is however more
 limited than in the classical case, because (1) sheaves of solutions of
 linear differential equations 
 are usually not locally constant, and (2) many spaces (for instance, 
 annuli) are simply connected.
\end{abst}


\subsection{The Monodromy principle.}
 
Let $S$ be a topological space, and $\mathcal{F}$ an abelian sheaf on $S$.
For an open set $U\subseteq S$ and a section $f\in\Gamma(U,\mathcal{F})$, the 
support of $f$ is the subset $\Supp(f)=\{u\in U\ |\ f_u\neq 0\}$, 
which is easily seen to be closed in $U$.

\begin{dfn}\label{dfn-unique-cont}
We say that the sheaf $\mathcal{F}$ satisfies the {\it principle of unique
continuation} if for any open set $U\subseteq S$ and any section 
$f\in\Gamma(U,\mathcal{F})$ the support $\Supp(f)$ is open in $U$.
\end{dfn}

Note that the principle of unique continuation is a local property.

\begin{lem}\label{lem-unique-cont}
If $\mathcal{F}$ satisfies the principle of unique continuation, then
any two sections $f,g\in\Gamma(U,\mathcal{F})$
on a connected open subset $U\subseteq S$ coincide if (and only if)
their germs $f_s$, $g_s$ at some point $s\in U$ coincide.
The converse also holds if $S$ is locally connected.
\end{lem} 

\begin{proof}
 The first assertion is clear. For the converse, let $u$ be a point adherent to the complement 
 of $\Supp(f)$. There exists a sufficiently small connected open
 neighborhood $V$ of $u$ contained in $U$. Since $f_s=0$ at some
 point $s\in V$, $f\equiv 0$ on $V$. Hence $u\not\in\Supp(f)$. This shows that $\Supp(f)$ is open.
\end{proof}

\begin{lem}\label{lem-separated}
 Let $p: F\rightarrow S$
 be the local homeomorphism canonically attached to
 $\mathcal{F} $: $\mathcal{F}(U)$ is the set of (continuous) sections of $p$ over $U$. 
 If $F$ is separated (\ie Hausdorff), then $\mathcal{F}$ satisfies the
 principle of unique continuation. The converse also holds if $S$ is
 separated.
\end{lem}

\begin{proof} If $F$ is separated, then for any two sections $f, g$ of
 $p$ over an open subset $U\subset S$, the set of $s\in U$ such that
 $f(s)=g(s)$ is closed.  Applying this to the zero-section $g=0$, we see
 that the support of $f$ is open.

 Conversely, let $x,y $ be two points of $F$. If $p(x)\neq p(y)$, there
 are disjoint open neighborhoods $V_x, V_y$ of $p(x)$ and $p(y)$
 respectively in the separated space $S$; then $p^{-1}V_x$ and
 $p^{-1}V_y$ are disjoint open neighborhoods of $x$ and $y$
 respectively. If $p(x)= p(y)$, there are open neighborhoods $U_x, U_y$
 of $x$ and $y$ respectively, an open subset $V\subset S$, and sections
 $f,g$ of $p$ over $V$, such that $f(V)=U_x,\;g(V)=U_y$. By the
 principle of unique continuation, $f\neq g$ defines an open subset
 $W\subset V$ containing $p(x)=p(y)$; then $f(W)$ and $g(W)$ are
 disjoint open neighborhoods of $x$ and $y$ respectively.
\end{proof}

\begin{dfn}
 We say that the sheaf $\mathcal{F}$ 
 satisfies the {\it monodromy principle} if it has the following property:

 let $\Gamma: [a,b]\times [0,1] \rightarrow S$ be any 
 continuous map with
 $\Gamma(\{a\}\times [0,1])=\{x_a\}, \Gamma(\{b\}\times
 [0,1])=\{x_b\}$. Let $f_{x_a}$ be an element of the stalk
 $\mathcal{F}_{x_a}$. Assume that for any
 $t\in [0,1]$, $f_{x_a}$ extends to a global section $f^{\Gamma_t}$ of
 $\Gamma_t^{\ast} \mathcal{F}$ on
 $[a,b]$. Then this extension is unique and
 $f^{\Gamma_t}(b)\in \mathcal{F}_{x_b}$ is independent of $t$.

 In this situation, the section $f^{\Gamma_t}$ of 
 $\Gamma_t^{\ast} \mathcal{F}$ is called the
 {\it continuation of $f_{x_a}$ along the path $\Gamma_t$}, and 
 $f^{\Gamma_t}(b)\in \mathcal{F}_{x_b}$ its {\it value at $x_b$}.
\end{dfn}

\begin{pro}
 If $\mathcal{F}$ satisfies the principle of
 unique continuation, then it satisfies the principle of monodromy. 
 The converse also holds if
 $S$ is locally arcwise connected.
\end{pro}

\begin{proof}
 Assume that $\mathcal{F}$ satisfies the principle of
 unique continuation. This guarantees the uniqueness of the 
 continuation $f^{\Gamma_t}$ of
 $f_{x_a}$ along the path $\Gamma_t$ for any $t$. Moreover, by a 
 special case of the proper base change theorem
 (applied to the first projection $[a,b]\times [0,1]\rightarrow 
 [a,b]$, \cf \cite[IV, 1.4]{cos}) $f^{\Gamma_t}$
 extends to a section of
 $\Gamma^{\ast}
 \mathcal{F}$ on a suitable subset of the form $[a,b]\times 
 (t-\epsilon,t+\epsilon)$. By unicity, these sections glue
 together to a global section of $\Gamma^{\ast} \mathcal{F}$. The restriction 
 of this sheaf to
 $\{b\}\times [0,1]$ is the constant sheaf with stalk
 $\mathcal{F}_{x_b}$. Therefore the value 
 $f^{\Gamma_t}(b)\in \mathcal{F}_{x_b}$ is independent of $t$.

 Conversely, let $U$ be an arcwise connected open subset of $S$, let
 $x_a, x_b$ be two points of
 $U$, and let $f$ be a section of
 $\mathcal{F}$ over $U$. Let $\gamma: [a,b]\rightarrow S$ be any path from 
 $x_a$ to $ x_b$. It
 follows from the unicity of continuation of $f_{x_a}$ along $\gamma$ 
 (requested in the monodromy
 principle) that if $f_{x_a}$ is
 $0$, so is
 $f_{x_b}$.
\end{proof}

\begin{exa}\label{lem-homotopy}
  (1) Any locally constant abelian sheaf $\mathcal{F}$ on a
 topological space $S$ satisfies the principle of unique 
 continuation, hence the principle of monodromy.
 
 In fact, for any $\Gamma: [a,b]\times [0,1]\rightarrow S$ as in
 \lemref{lem-separated}, the inverse image
 $\Gamma^{\ast}\mathcal{F}$ is locally constant, hence constant and 
 canonically isomorphic to the
 constant sheaf attached to
 $\mathcal{F}_{x_a}$. Therefore the extension $f^\gamma$ of any 
 $f_{x_a}\in \mathcal{F}_{x_a}$ along any
 path $\gamma$ exists, and the value at the other extremity 
 $x_b=\gamma(b)$ depends only on the
 \emph{homotopy class} of $\gamma$.

  (2) When $S$ is a complex manifold, the structure sheaf
 $\mathcal{O}_S$ satisfies the principle of unique continuation, hence
 the principle of monodromy.
\end{exa}


\begin{para}
 Let $S$ be a topological space, connected and locally arcwise 
 (or simply) connected,
 and $\mathscr{F}$ a locally constant abelian sheaf on $S$.
 We fix a point $s\in S$, and denote by $\pi_1(S,s)$ the fundamental
 group based at $s$.\footnote{to remove any ambiguity, let us say that we
 adopt Deligne's convention: the composition in $\pi_1(S,s)$
 is induced by the juxtaposition of loops in the reverse order. As 
 such, $\pi_1(S,s)$ acts \emph{on the right} on the
 pointed universal covering.}
 To any loop $\gamma\colon [0,1]\rightarrow S$ based at $s$, let us associate
 the so-called {\it monodromy} along $\gamma$, defined by the composite
 \begin{equation*}
  \mathscr{F}_s\stackrel{\sim}{\longrightarrow}(\gamma^{\ast}\mathscr{F})_0
   \stackrel{\sim}{\longrightarrow}\Gamma([0,1],\gamma^{\ast}\mathscr{F})
   \stackrel{\sim}{\longrightarrow}(\gamma^{\ast}\mathscr{F})_1
   \stackrel{\sim}{\longrightarrow}\mathscr{F}_s.
 \end{equation*}
 By \examref{lem-homotopy} (1), the monodromy
 % $\mathscr{F}_s\stackrel{\sim}{\rightarrow}\mathscr{F}_s$
 $\mathscr{F}_s\xrightarrow{\sim}\mathscr{F}_s$
 depends only on the class of $\gamma$ in
 $\pi_1(S,s)$,
 hence gives rise to a left $\Z[\pi_1(S,s)]$-module structure on 
 $\mathscr{F}_s$.
 This construction yields an equivalence of the categories
 \begin{equation*}
  \left\{\textrm{locally constant sheaves on} \; S\right\}
   \stackrel{\sim}{\longrightarrow}
   \left\{\textrm{left\ }\Z[\pi_1(S,s)]\textrm{-modules}\right\}\;:
 \end{equation*}
  giving a locally constant sheaf amounts to giving its value at $s$
 together with the monodromy action.
\end{para}

%
%
\subsection{Rigid geometry and the problem of unique continuation.}
\label{subsection-rigid-and-unique-cont}
\begin{para}
 For $p$ a prime number, let $\Q_p$ denote as usual the
 completion of $\Q$ for the $p$-adic absolute value
 $|\;\;|_p$ : $|p^n\frac{a}{b}|_p = p^{-n}$ if the rational integers $a,
 b$ are prime 
 to $p$. This ultrametric absolute value extends in a unique way to each
 finite extension of $\Q_p$. These finite extensions are locally
 compact and totally disconnected. They are all complete, but ``the''
 algebraic closure $\ovl{\Q}_p$ of $\Q_p$ itself is not complete. Its
 completion, denoted by $\C_p$, turns out to be algebraically
 closed, and plays the role of $\C$ in $p$-adic analysis.


 \noindent In the sequel, $\matheur{D}(a,r^+)$
 (\resp $\matheur{D}(a,r^-)$) stands for the disk --- archimedean or not 
 --- of radius $r$ centered at $a$ with (\resp without) circumference.

 It might be surprising at first that geometries can be built upon
 $p$-adic numbers, whose ``fractal'' nature makes them hardly amenable to
 intuition as a continuum.  Nevertheless, Bourbaki's presentation of
 analytic geometry \cite{vdea} treats the real, complex and $p$-adic
 cases on equal footing. This approach is based on a local definition of
 analytic functions as sums of convergent power series. Its major
 drawback is that these analytic functions fail to satisfy the principle
 of unique continuation, essentially because two ultrametric disks are
 either concentric or disjoint (like drops of mercury)

 %PICTURE1
 \begin{figure}[h]
 \begin{picture}(250,60)(0,0)
  \put(0,0){\includegraphics{fig1.eps}}
 \end{picture}
 \caption{}
 \label{fig1}
 \end{figure}

\end{para}

\begin{para}
 M. Krasner had the idea to remedy this by using a definition of
 analytic functions \emph{\`a la Runge}.
 He properly founded ultrametric analysis by introducing his \emph{analytic
 elements} defined as
 uniform limits of rational functions, a global notion which overcomes,
 to some extent, problems stemming from the disconnectedness of the $p$-adics.

 The next step is due to J.~Tate. In order to deal with more general spaces
 than just subsets of the line, he introduced and developed the
 so-called \emph{rigid analytic geometry} (as opposed to Bourbaki's
 ``wobbly'' analytic geometry), based on affinoid algebras (topological
 $K$-algebras isomorphic to quotients of rings of restricted formal
 power series, \ie whose coefficients tends to $0$) and a suitable
 Grothendieck topology \cite{ras}.

 In a connected rigid analytic variety $S$, one has the following 
 avatar of analytic continuation \cite[0.1.13]{crecrasp}:  assume that
 there is no admissible covering formed by two disjoint non-empty open
 subsets, and let $f$ be an analytic function 
 on $S$. If there is a connected open subset $U$ over which $f$ 
 vanishes, then $f=0$.
\end{para}

\begin{para}
 The first achievement of rigid geometry was Tate's representation of an
 elliptic curve with ``bad reduction at $p$'' as a rigid analytic quotient
 $\C_p^{\times}/q^{\Z}$ of the multiplicative group by the discrete
 group generated by $q$.  Here $q$ is given by the usual series
 $\frac{1}{j}+\ldots$ in the $j$-invariant of the elliptic 
 curve ($|j| >1$), interpreted $p$-adically.  The quotient is
 obtained by gluing two affinoid annuli of width $|q|^{-1/2}$
 along their boundary; the inverse image in $\C_p^{\times}$ of each of
 these annuli consists of countably many disjoint copies of it.  This is
 analogous to Jacobi's partial uniformization $\C^{\times}/q^{\Z}$
 of a complex elliptic curve.
\end{para}

\begin{para}
 This result was extended by D.~Mumford to
 curves of higher genus \cite{aacodcoclr}.  Here, the
 partial uniformization which serves as a complex model is the 
 Schottky uniformization of complex curves of genus
 $g>1$, which we briefly recall (\cf \cite{osats}).  Let
 $D$ be an open subset of ${\mathbb{P}}^1(\C)$, limited by $2g$ 
 disjoint circles $C_1,C'_1,\ldots,C_g,C'_g$.  One
 assumes the existence, for any $i=1\ldots, g,$ of an element 
 $\gamma_i\in PGL_2(\C)$, with two distinct fixed
 points, which sends $C_i$ to $C'_i$ and $D$ outside itself. Then the 
 subgroup $\Gamma$ generated by the
 $\gamma_i$'s is discrete and freely generated by them, the complement 
 $\Omega\subset {\mathbb{P}}^1(\C)$ of the
 topological closure of the set of fixed points of
 $\Gamma$ is an open dense subset ($=\bigcup_{\gamma \in \Gamma} 
 \gamma(D)$), and $\Omega/\Gamma$ is a projective
 smooth complex curve of genus $g$.

 When $\C$ is replaced by $\C_p$, the same construction applies. The
 rigid analytic quotient $\Omega/\Gamma$ exists and is called a Mumford
 curve. Among projective smooth curves of genus $g$ over $\C_p$, Mumford
 curves are characterized by the existence of a reduction over the
 residue field ${\ovl{\F}}_p$ such that every irreducible component is
 isomorphic to ${\mathbb{P}}^1$ and intersect the others at double
 points, \cf \cite{sgamc}.

 Let us mention that T. Ichikawa has proposed a unified theory of 
 the archimedean and non-archimedean
 Schottky-Mumford uniformizations, \cf \cite{sutorsamcoig}.
\end{para}

\subsection{Berkovich geometry and the principle of monodromy}

\begin{para}
 Rigid analytic spaces are endowed with a 
 Grothendieck generalized topology, and their
 structure sheaf is a sheaf with respect to this topology. Hence
 it cannot be said to satisfy the principle of
 unique continuation in the strict sense of \defref{dfn-unique-cont}.
 Moreover, there is no non-trivial path in such spaces.  Therefore,
 rigid geometry is not a suitable setting for discussing $p$-adic 
 analytic continuation in an intuitive way.

 In contrast, Berkovich's viewpoint on $p$-adic geometry does 
 not suffer from these drawbacks: Berkovich's
 analytic spaces are genuine locally ringed topological spaces, which 
 are even locally arcwise connected. We refer to
 \cite{spasalc} for a compact technical presentation.
\end{para}

\begin{para}
 The buildings blocks are the same: affinoid algebras $\matheur{A}$
 (called strictly affinoid algebras by Berkovich), \ie topological
 $K$-algebras isomorphic to quotients of rings of restricted power
 series. 
 However, instead of attaching to $\matheur{A}$ its maximal spectrum
 $\mathrm{Spm}(\matheur{A})$ as in rigid geometry, 
 Berkovich analytic geometry deals with the {\it affinoid space}
 $\matheur{M}(\matheur{A})$ of all bounded multiplicative seminorms
 on $\matheur{A}$.

 This ``spectrum'' contains ``more points'' than the maximal spectrum
 $\mathrm{Spm}(\matheur{A})$, namely something 
 like ``generic points'', which ``complete'' $\mathrm{Spm}(\matheur{A})$.
 There is a natural inclusion $\mathrm{Spm}(\matheur{A})\subseteq 
 \matheur{M}(\matheur{A})$ which identifies $\mathrm{Spm}(\matheur{A})$
 with the subspace of all seminorms $|\cdot|$ with
 $\matheur{A}/\Ker |\cdot| =\C_p$.
 This ``completion'', in fact, simplifies the topology; e.g.\ Berkovich
 analytic spaces are locally arcwise connected. 
\end{para}

\begin{para}
 {\itshape Berkovich's affinoids.\ }
 Let us be a little more precise about the definition of 
 {\it affinoid spaces} in Berkovich geometry.
 Let $\matheur{A}$ be an affinoid algebra over a complete subfield $K$ of 
 $\C_p$.

 \begin{enumerate}
  \item A point of $\matheur{M}(\matheur{A})$ is a bounded multiplicative
	seminorm on $\matheur{A}$.
  \item The topology of $\matheur{M}(\matheur{A})$ is the weakest one so
	that the mapping
	$\matheur{M}(\matheur{A})\ni\chi\mapsto\chi(f)\in \R_{\ge 0}$
	is continuous for any $f\in\matheur{A}$.
  \item The sheaf of rings $\OO_{\matheur{M}(\matheur{A})}$ is defined by 
	$\Gamma(U,\OO_{\matheur{M}(\matheur{A})})=\varprojlim\matheur{A}_V$, 
	where $V$ runs over finite unions $\bigcup V_i$ of affinoid
	domains contained in $U$, and
	$\matheur{A}_V= \Ker\bigl(\prod_i\matheur{A}_{V_i}
	\rightrightarrows
	\prod_{ij}\matheur{A}_{V_i\cap V_j}\bigr)$, $\matheur{A}_{V_i}$
	being the affinoid algebra attached to the affinoid domain $V_i$.
 \end{enumerate}

 The value of a ``function'' $f\in \matheur{A}$ at a point $\chi\in 
 \matheur{M}(\matheur{A})$ is its image in the field
 $\matheur{A}/\Ker\chi$. This field inherits the absolute value induced by 
 the seminorm $\chi$, and its completion is
 denoted by $\mathscr{H}(\chi)$.\index{000H@$\mathscr{H}(\;)$}
\end{para}

\begin{exa}\label{exa-berkovich-space}
Let $K\{t\}$ denote the ring of restricted power series in one
variable.
Let us assume, for simplicity, that $K$ is algebraically closed.
In rigid geometry, the maximal spectrum $\mathrm{Spm}(K\{t\})$ is
just the closed disk $\matheur{D}(0,1^{+})$ in $K$ of radius $1$ in the usual sense.
In $\matheur{M}(K\{t\})$, the following four kinds of points occur:
\begin{enumerate}
 \renewcommand{\labelenumi}{(\arabic{enumi})}
 \item classical points (\ie those coming from $\mathrm{Spm}(K\{t\})$):
       $x\in\matheur{D}(0,1^{+})$, $\chi_x(f)=|f(x)|_K$ for $f\in K\{t\}$,
 \item  generic points of disks: $\chi=\eta_{x,r}$ for $0<r\leq 1$ with 
       $r\in |K^{\times}|$, 
       $\chi(f)=|f|_{\matheur{D}(x,r^{+})}$ (the sup-norm), 
 \item the same, for $r\not\in |K^{\times}|$,
 \item generic points of infinite decreasing families
       $\{\matheur{D}_{\alpha}\}$
       of closed disks with radius $\leq 1$: 
       $\chi(f)=\mathrm{inf}|f|_{\matheur{D}_{\alpha}}$.
\end{enumerate}

\vspace{2ex}%PICTURE2
\begin{figure}[h]
 \setlength{\unitlength}{1pt}
 \begin{picture}(150,60)(0,0)
  %\put(20,10){\includegraphics{pic1.eps}}
  %\put(37,54){$\matheur{t}_{0|1}=\matheur{t}_{1|1}$}
  \put(25,10){\line(1,1){35}}
  \put(95,10){\line(-1,1){35}}
  \put(25,9){\makebox(0,0)[c]{$\bp$}}
  \put(95,9){\makebox(0,0)[c]{$\bp$}}
  \put(60,44){\makebox(0,0)[c]{$\bp$}}
  \put(45,29){\makebox(0,0)[c]{$\bp$}}
  \put(55,58){\makebox(0,0)[l]{$\eta_{x,|x-y|}$}}
  \put(17,0){$x$}
  \put(96,0){$y$}
  \put(20,35){$\matheur{\eta}_{x,r}$}
  %\put(80,35){$\matheur{t}_{1|s}$}
 \end{picture}
 \caption{Generic points and paths on a Berkovich space.}
 \label{pic1}
\end{figure}
\vspace{1ex}

Therefore, two arbitrary distinct points in $\matheur{M}(K\{t\})$ can be 
connected by a unique path.
For example, two points of type (1), $x$ and $y$, are connected by a
path consisting of points of type (2) and (3) associated to the disks
$\matheur{D}(x,r^{+})$ and $\matheur{D}(y,r^{+})$ for $0<r\leq |x-y|$.
The complement of a point of type (1) or (4) is connected.
The affine line $\mathbb{A}^1$ is the union of the affinoid spaces
associated to the algebras of power series convergent in 
$\matheur{D}(0,r^{+})$, and the projective line $\mathbb{P}^1$ is the
Alexandroff compactification of $\mathbb{A}^1$.
\end{exa}

\begin{rem}
 Adding generic points to rigid spaces 
 to come up with Berkovich spaces simplifies the topology in contrast to 
 what happens in algebraic geometry, when
 generic points are added to varieties to produce schemes.
 Actually Berkovich's generic points are closed, unlike Grothendieck's
 ones.
\end{rem}

\begin{para}\label{para:strict-analytic}
 The construction of (strictly) analytic $K$-spaces by gluing affinoid
 spaces together is a little delicate (one is not gluing open
 subspaces); we refer to \cite{ecfnas}.  These spaces are \emph{locally 
 compact, locally countable at infinity, and locally arcwise connected.} 

 There is a fully faithful functor between Berkovich's Hausdorff 
 (strictly) analytic $K$-spaces and rigid analytic varieties: at the
 level of underlying sets, this functor sends a space $S$ to the subset of
 ``classical points'', \ie points $s$ for which $[\mathscr{H}(s):K]<\infty$.

 This functor establishes \emph{an equivalence of categories 
 between the category of paracompact (strictly)
 analytic $K$-spaces and quasi-separated rigid spaces over $K$ having 
 an admissible affinoid covering of finite type}.

 This equivalence respects the notion of dimension
 (topological dimension in the case of Berkovich analytic spaces), as
 well as the standard properties of local rings such as: reduced,
 normal, smooth..., and the properties: finite, \'etale...  of
 morphisms. Also, a paracompact analytic space is connected (in the
 usual sense) if and only if the corresponding rigid space does not
 admit an admissible covering by two disjoint non-empty open subsets.

 
 Furthermore, there is a canonical functor ``analytification''
 \begin{equation*}
  \left\{
  \begin{matrix}
   \text{separated schemes}\\
   \text{locally of finite type}\\
   \text{over $K$}
  \end{matrix}
  \right\}
  \rightarrow
  \left\{
  \begin{matrix}
   \text{paracompact (strictly)}\\
   \text{analytic $K$-spaces}
  \end{matrix}
  \right\}
 \end{equation*}
 \noindent
 and a canonical functor ``generic fiber''
 \begin{equation*}
  \left\{
  \begin{matrix}
   \text{separated formal schemes locally}\\
   \text{finitely presented over $\mathcal{O}_K$}
  \end{matrix}
  \right\}
  \rightarrow
  \left\{
  \begin{matrix}
   \text{paracompact (strictly)}\\
   \text{analytic $K$-spaces}
  \end{matrix}
  \right\}.
 \end{equation*}
\end{para}

\begin{para}\label{para:p-adic-manifold}
 {\itshape $p$-adic manifolds.} We shall be mainly 
 concerned with paracompact (strictly) $K$-analytic spaces $S$ which
 satisfy the following assumption:
 \begin{quote}
  any $s\in S$ has a neighborhood $U(s)$ which is isomorphic 
  to an affinoid subdomain of some space
  $V_s$ which admits locally an \'etale morphism to the affine space 
  $\mathbb{A}^{\dim S}$.
 \end{quote}
 For convenience, we shall call such a space $S$ a {\it
 $K$-manifold}\index{Kmanifold@$K$-manifold}, and simply 
 a {\it $p$-adic manifold}\index{padicmanifold@$p$-adic manifold} if $K=\C_p$.

 \noindent It is an important result of V. Berkovich \cite{spasalc}
 that \emph{$p$-adic manifolds are locally contractible}.  Therefore
 they have universal coverings. 
\end{para}

\begin{pro}  
 The structure sheaf of any $p$-adic 
 manifold $S$ satisfies the principle of unique
 continuation, and (equivalently) the principle of monodromy.
\end{pro}
\begin{proof} 
 Since the principle of unique continuation is local, we may replace $S$
 by the neighborhood $U(s)$ of any 
 point $s$ given in advance. Thus we may assume that $S$ 
 itself is affinoid: $S=\matheur{M}(\matheur{A})$,
 and connected, and we have to show that the homomorphism 
 $\iota_s:\;\matheur{A}\rightarrow \matheur{A}_s\simeq
 \underset{s\in V}{\liminj}\matheur{A}_V$ is injective ($V$ runs over
 the affinoid neighborhoods of $s$).  This follows from the
 fact that the homomorphism of completion at 
 $s : \matheur{A}\rightarrow  \widehat{\matheur{A}}_s =
 \underset{n}{\limproj}\matheur{A}/\mathcal{I}_s^n\;$ is injective and factors
 through $\iota_s$.
\end{proof}  

In the case of curves (\ie $p$-adic manifolds of dimension $1$),
analytic continuation is particularly intuitive, because there is a basis
of open subsets $\mathcal{U}$ with finite boundary such that two arbitrary
points in $\mathcal{U}$ are connected by a unique geometric path lying in
$\mathcal{U}$.

%However, while every simply connected curve is isomorphic to one of
%$\matheur{D}(0,1^{-})$, $\mathbb{A}^1$, or $\mathbb{P}^1$ in the
%complex case, there exist a plenty of those ultrametric curves; for
%example, every curve having good reduction.

\subsection{Topological coverings and \'etale coverings.}
\label{sub-top-cov-et-cov}
\begin{para}\label{para-complex-vs-padic}
 Complex manifolds are locally contractible,
 and have universal coverings.  There is no need to distinguish between
 topological
 coverings and \'etale coverings (finite or infinite).

 In $p$-adic rigid geometry, the situation is more complicated. It is
 natural to call
 topological covering any morphism $f: Y \rightarrow X$ such that there is
 an admissible
 cover $(X_i)$ of $X$ and an admissible cover $(Y_{ij})$ of $f^{-1}(X_i)$
 with disjoint $Y_{ij}$
 isomorphic to $(X_i)$ via  $f$. Indeed, such topological coverings
 correspond to locally constant
 sheaves of sets on $X$. It is still true that \emph{topological coverings
 are \'etale, but the converse is wrong}, even if one restricts to
 finite surjective morphisms. Indeed, the Kummer covering
 $z \mapsto z^n$ of the punctured unit disk is an \'etale 
 covering, but not a topological covering, if $n>1$.
\end{para}

\begin{para}
 It is again more convenient to tackle these questions in the framework
 of Berkovich's geometry.  For instance, one sees immediately that a
 Kummer covering $z\mapsto z^n$ as above is not a topological covering
 because a classical point has $n$ preimages, while the ``generic
 point'' $\eta_{0,1}$ (corresponding to the sup-norm on the disk) is its
 own single preimage. 
 Topological coverings of a $p$-adic manifold $X$ are defined in the
 usual way; they correspond to locally constant sheaves of sets on $X$.
 They coincide \cite[3.3.4]{staagonf} with topological coverings of the
 rigid analytic variety associated to $X$.
\end{para}

\begin{para}
 There are only three one-dimensional simply connected complex manifolds up to
 isomorphism: the projective line ${\P}^1(\C)$, the affine line $\C$, and
 the disk $\matheur{D}$ ($\simeq \mathfrak{h}$, the complex upper half
 plane). In contrast, there are plenty of \emph{simply-connected}
 one-dimensional 
 $p$-adic manifolds, for instance: any annulus, the line deprived from a
 finite number of points, more generally any connected $p$-adic manifold
 homeomorphic to a \emph{subset of} $\;{\P}^1(\C_p)$, any smooth
 projective curve with ``good reduction''... (all these Berkovich spaces
 look like ``bushy trees'').
\end{para}

\begin{para}
 One defines the (discrete) topological fundamental group
 $\pi_1^{\mathrm{top}}(S,s)$ of
 a pointed $p$-adic manifold $(S,s)$ in the usual way. The general
 topological theory of
 coverings applies: $(S,s)$ is naturally isomorphic to the quotient of the
 pointed universal
 covering $({\til{S}},{\tilde{s}})$ by $\pi_1^{\mathrm{top}}(S,s)$, and
 $\pi_1^{\mathrm{top}}(S,s)$ classify the
 topological coverings of $(S,s)$.   In the one-dimensional case, the
 topological fundamental
 group is a discrete \emph{free} group (more precisely, it is isomorphic to
 the fundamental
 group of the dual graph of the so-called semistable reduction of $S$
 (\cite[5.3]{efgonas})). 

 For instance, if $S=\Omega/\Gamma$ is the
 uniformization of a Mumford curve
 (\cf~\ref{subsection-rigid-and-unique-cont}), $\Omega$ is the universal
 covering of $S$ and the Schottky group $\Gamma$ is isomorphic to the 
 topological fundamental group.

% \noindent
 On the other hand, one can define the (profinite) algebraic fundamental
 group $\pi_1^{\mathrm{alg}}(S,s)$ \`a la Grothendieck, classifying the finite
 \'etale coverings of $(S,s)$. In contrast to the complex situation, the
 natural map $\pi_1^{\mathrm{alg}}(S,s) \rightarrow
 \pi_1^{\mathrm{top}}(S,s)^{\wedge}$
 to the profinite completion of $\pi_1^{\mathrm{top}}(S,s)$ is
 \emph{generally not injective} (\eg for annuli).
\end{para}

\subsection{Connections with locally constant sheaves of solutions.}
\label{sub-connection}
\begin{para}
 Let us briefly recall the complex situation. Let $S$ be a complex
 connected manifold, $(\mathcal{M},\nabla)$ a vector bundle of rank $r$
 with integrable connection on $S$. The classical Cauchy theorem shows
 that for any $s \in S$, the solution space
 $(\mathcal{M}\otimes\mathcal{O}_{S,s})^{\nabla}$ at $s$ has dimension
 $r$. Analytic continuation along paths gives rise to a homomorphism
 $\pi_1^{\mathrm{top}}(S,{s}) \rightarrow
 \Aut_{\C}((\mathcal{M}\otimes\mathcal{O}_{S,s})^{\nabla})$ (the
 \emph{monodromy}). The sheaf of germs of solutions
 $\mathcal{M}^{\nabla}$ is locally constant on $S$: its pull-back over
 the universal covering $\til{S}$ of $S$ is constant. Conversely, any
 complex representation $V$ of $\pi_1^{\mathrm{top}}(S,{s})$ of dimension $r$
 gives rise naturally to a vector bundle $\mathcal{M}$ of rank $r$ with
 integrable connection $\nabla$:  $\mathcal{M}= (V \times
 \til{S})/\pi_1^{\mathrm{top}}(S,{s}), \; \nabla(V) = 0$.  This sets up an
 equivalence of categories:
 \begin{equation*}
  \left\{
  \begin{matrix}
   \text{finite dimensional}\\
   \text{representations of $\pi_1^{\mathrm{top}}(S,{s})$}
  \end{matrix}
  \right\}
  \simeq
  \left\{
  \begin{matrix}
   \text{$S$-vector bundles with}\\
   \text{integrable connection}
  \end{matrix}
  \right\}.
 \end{equation*}
\end{para}  
\begin{para}
 Let $(S,s)$ be now a pointed connected $p$-adic manifold. It is still
 true that any $\C_p$-linear representation $V$ of
 $\pi_1^{\mathrm{top}}(S,s)$ of dimension $r$ gives rise to a vector
 bundle $\mathcal{M}=\mathcal{M}_V$ a vector bundle of rank $r$ with
 integrable connection $\nabla =\nabla_V$ (same formula).
 The functor
 \begin{equation*}
  \left\{
  \begin{matrix}
   \text{finite dimensional}\\
   \text{representations of $\pi_1^{\mathrm{top}}(S,s)$}
  \end{matrix}
  \right\}
  \rightarrow
  \left\{
  \begin{matrix}
   \text{$S$-vector bundles with}\\
   \text{integrable connection}
  \end{matrix}
  \right\}
 \end{equation*}
 is still fully faithful, but \emph{no longer surjective}; its essential
 image consists of those connections whose sheaf of solutions is locally
 constant, \ie  becomes constant over $\til{S}$.  In fact, the
 classical ``\emph{Cauchy theorem}'' according to which the solution
 space $(\mathcal{M}\otimes\mathcal{O}_{S,s})^{\nabla}$ at $s$ has
 dimension $r$ holds for every classical point $s$ of $S$ --- which
 corresponds to a point of the associated rigid variety ---, but
 \emph{does not hold} for non-classical points $s$ of the Berkovich
 space $S$ in general.
\end{para}

\begin{para}\label{para-cauchy} 
 Any non-trivial connection over the projective line minus a few points
 gives an example when the $p$-adic analogue of the ``Cauchy theorem''
 does not hold: indeed, in this case, the topological fundamental group
 is trivial.  

 However, a $p$-adic variant of Cauchy's theorem in the neighborhood of a
 non-classical point $s$ may be restored as follows (Dwork's technique of
 generic points): it suffices to extend the scalars from $\C_p$
 to a complete algebraically closed extension of $\mathscr{H}(s)$.  This
 transforms $s$ into a classical point, and neighborhoods of $s$ after
 scalar extensions are ``smaller'' than before.
\end{para}

\begin{para}\label{para-cauchy-monodromy} 
 When ``Cauchy's theorem'' holds at every point of $S$, one can continue
 the local solutions along paths and get the monodromy representation
 just as in the complex situation. 

 This nice category of $p$-adic connections has not yet
 attracted much attention.
\end{para}

\begin{exa}\label{exa-tate-ell-curve}
 Let us consider the case when $S$ is a Tate elliptic curve:
 $S=\C_p^{\times}/q^{\Z}$, with $s =$ its origin. Then
 $\pi_1^{\mathrm{top}}(S,s)=q^{\Z}$, and the connections on $S$ which arise
 from representations of $q^{\Z}$ are those which become trivial over
 $\til{S}=\C_p^{\times}$.
 In this correspondence, the multivalued solutions $\vec y$ (\ie 
 horizontal sections) of such a connection are the solutions of the
 associated linear $q$-difference equation $\vec y(qt)=M(q) \vec y(t)$,
 where $t$ is the standard coordinate on $\C_p^{\times}$ and the
 matrix $M(q)$ is the image of $q$ in the monodromy representation.

 The simplest example is given by the obvious one-dimensional 
 representation $M(q)=q$. It corresponds to
 $\mathcal{M}=\mathcal{O}_S,\;\nabla(1)=\omega_{\mathrm{can}}$ (the canonical
 differential on $S$ induced by $dt/t$); here the $q$-difference
 equation is $t.dy=y.dt$ with obvious solution $y=t$. 

 Let us now look at the representation $M(q)=\sqrt{q}$.  The associated
 $q$-difference equation is $t.dy=\frac{1}{2}y.dt$. There is the 
 obvious solution $\sqrt t$, which leads to
 $\mathcal{M}=\mathcal{O}_S,\;\nabla(1)=\frac{1}{2}\omega_{\mathrm{can}}$. 
 Here we encounter an interesting paradox: $\sqrt{t}$ is not a multivalued
 analytic function on $S$ (\ie it is not an analytic function on the
 universal covering $\C_p^{\times}\;$.  In the complex situation, a
 similar paradox arose in the work of G. Birkhoff in his theory of
 $q$-difference equations, which was pointed out
 and analyzed by M.~van der Put and M.~Singer in the last chapter of
 their book \cite{gtode}.  The solution of the paradox is that the vector
 bundle $\mathcal{M}$ associated to the representation $q \rightarrow
 \sqrt{q}$ of $q^{\Z}$ (or to the $q$-difference equation
 $y(qt)=\sqrt{q}.t$) is in fact a non-trivial vector bundle of rank one,
 and the basic solution is not $\sqrt{t}$, but
 $\frac{\theta(t/\sqrt{q})}{\theta(t)}$, where
 $\theta(t)=\prod_{n>0}(1-q^nt)\prod_{n\leq 0}(1-q^n/t)$.

\end{exa}

\begin{para}\label{para-Simpson} 
 In the previous example, it is easily seen that rank-one vector bundles
 with connection on $S$ which arise from a representation of
 $\pi_1^{\mathrm{top}}(S,s)$ form a space of dimension one, while the
 space of all rank-one vector bundles with connection on $S$ has
 dimension two. 

 We next consider the case of a $p$-adic manifold $S$ which ``is'' a
 geometrically irreducible algebraic ${\C}_p$-curve. Let $\ovl{S}$ be
 its projective completion. It follows from the Van Kampen theorem,
 together with the fact that punctured disks are simply-connected, that
 $\pi_1^{\mathrm{top}}(S,s) \to \pi_1^{\mathrm{top}}(\ovl{S},s)$ is
 an isomorphism. It 
 follows that the vector bundles with connection attached to
 representations of $\pi_1^{\mathrm{top}}(S,s)$ automatically extend
 to $\ovl{S}$. Hence we may assume without loss of generality that $S$ is
 compact.

 Vector bundles with connection on $S$ are algebrizable, and one can
 use C.~Simpson's construction (\cite{morotfgoaspv1}) to define the {\it
 moduli space of connections} of rank $r$ over $S$, denoted by
 $M_{\mathrm{dR}}(S,r)$. On the other hand, we have seen that the topological
 fundamental group $\pi_1^{\mathrm{top}}(S,s)$ is free on $b_1(\Delta)$
generators, being isomorphic to the fundamental group of the dual graph
$\Delta$ of the semistable reduction of $S$. Simpson has also studied
the {\it moduli space of representations of dimension $r$} of such a
group. We denote it by $M_{\mathrm{B}}(S,r)$; in fact, it depends only on the
couple of integers $(b_1(\Delta),r)$.

 Let us assume that $S$ is of genus $g\geq 2$. Simpson shows that 
 $M_{\mathrm{dR}}(S,r)$ is algebraic irreducible of dimension
 $2(r^2(g-1)+1)$.   On the other hand, $M_{\mathrm{B}}(S,r)$ is an algebraic
 irreducible affine variety of dimension $(r^2(b_1(\Delta)-1)+1)$.  We
 note that this dimension is maximal when the Betti number $b_1(\Delta
 )$ takes its maximal value, namely $g$. This corresponds to the case 
 where $S$ is a Mumford curve (a curve with totally degenerate reduction).  

 It turns out that the functor which associates a vector
 bundle with connection to any finite-dimensional representation of the
 topological fundamental group induces an {\rm injective analytic map}
 of moduli spaces $\iota : M_{\mathrm{B}}(S,r)\to M_{\mathrm{dR}}(S,r)$.

[The map is induced by the functor $V \to (\mathcal{M}_V,{\nabla}_V)$,
and its injectivity follows from the faithfulness of this functor. The
difficulty in showing that $\iota$ is analytic is that
$M_{\mathrm{dR}}(S,r)$ is 
a priori a moduli space for algebraic connections, not for all analytic
connections; we shall not pursue here in this direction].

In the complex situation, the corresponding map $\iota$ would always be
an analytic isomorphism (Riemann-Hilbert-Simpson). In the $p$-adic
cases, we see that the connections which satisfy Cauchy's theorem at
every point (classical or not) form a stratum of maximal dimension half
of the dimension of the total moduli space.
\end{para}


\vspace{10pt}\par
We shall leave the closer analysis of this kind of $p$-adic 
differential equations with global monodromy until later chapters, where
we show how they arise in the context of period mappings. In the next 
section, we shall deal with a very different kind of differential 
equations, which play a distinguished role in $p$-adic analysis under 
the name of {\it unit-root F-crystals}.



\newpage
\section{Analytic continuation; algebraic approach.}

\begin{abst}
 An equivalent, but more algebraic, approach to analytic continuation
 consists in interpreting complex analytic multivalued functions as
 limits of algebraic functions.
 Since the concept of topological coverings and that of \'etale coverings
 do not coincide in the $p$-adic setting, the algebraic approach 
 leads in that case to a theory quite different 
 from that of the previous section.  It turns out to be well-suited to
 the function-theoretic study of so-called unit-root $F$-isocrystals, and induces
 us to revisit two fundamental notions due to Dwork: Frobenius structure 
 and overconvergence.
\end{abst}

%
%
\subsection{Limits of rational functions.}\label{subsec-limit-rational}

\begin{para}\label{para-complex-approx}
{\itshape Over $\C$.\ }
 Let us briefly survey the complex analytic situation.
 Let $S$ be a connected complex analytic curve, and $U$ an open set
 of $S$ (not necessarily connected) which is holomorphically convex; 
 \index{holomorphically convex@holomorphically convex}
 \ie $S\setminus U$ has no compact connected component.
 Then $\OO(S)$ is dense in $\OO(U)$ for the topology of uniform
 convergence on every compact set.
 We denote this situation by $\OO(U)=\widehat{\OO(S)^U}$
 (\cf \eg \cite[13]{qadsdr}).

 Moreover, if $S$ is the Riemann surface coming from an affine algebraic
 curve $S^{\mathrm{alg}}$, then we also have 
 $\OO(U)=\widehat{\OO(S^{\mathrm{alg}})^U}$, which generalizes the theorem
 of Runge on approximation by rational functions.
 Indeed, considering an embedding $S^{\mathrm{alg}}\hookrightarrow
 (\mathbb{A}^N)^{\mathrm{alg}}$, we may extend analytic functions on $S$ to 
 analytic functions on $(\mathbb{A}^N)^{\mathrm{alg}}$
 (due to the vanishing of the first cohomology group of the
 coherent ideal sheaf defining $S$).
 An analytic functions on $(\mathbb{A}^N)^{\mathrm{alg}}$ can be 
 approximated by polynomials, which we restrict to $S$.
\end{para}

\begin{para}\label{para-ultra-approx}
{\itshape Over the $p$-adics.\ } We have an analogous situation. Let $K$ 
be a complete subfield of $\C_p$, and let $S$ be an analytic curve 
coming from a smooth affine algebraic curve over $K$. We consider a 
closed immersion $S \hookrightarrow {\A}^N$. For any $r > 0$, $S_r := S
\cap \matheur{D}_{{\A}^N}(0,r^+)$ is an affinoid domain in $S$. 
The same argument of polynomial approximation used in 
\ref{para-complex-approx} shows that $\mathcal{O}(S_r)$ is the 
completion of $\mathcal{O}(S^{\mathrm{alg}})$. 

 For any compact $Z \subset S$ (\eg an affinoid domain), we define the
 topological ring $\mathcal{H}(Z)$\index{000H@$\mathcal{H}(\;)$}
 of \emph{analytic elements}
 \index{analytic elements@analytic elements} on $Z$ to
 be the completion of $\Gamma (Z, \mathcal{O}_S) = \liminj_{Z\subset
 U \text{: open}}\Gamma (U, \mathcal{O}_S) $ under the sup-norm.
 Whenever $Z$ is contained in a $S_r$, $\mathcal{H}(Z)$ is also the
 completion of $\mathcal{O}(S_r)_Z$, where the subscript $Z$ denotes the
 localization with respect to the set of elements which do not vanish on
 $Z$ (\cf\ \cite{staagonf} and \cite{rdldaecpecd} for a more precise
version of the Runge theorem).

 The holomorphic convexity condition on $Z$ is that for some $r$, $Z$ is
 the intersection of affinoid neighborhoods defined by inequalities of
 the form $|f_i| \leq 1$ in $S_r$. If this condition is satisfied,
 $\mathcal{H}(Z)$ is in fact the completion of $\mathcal{O}(S_r)$ itself.
\end{para}

\begin{para}
{\itshape The Krasner-Dwork viewpoint.\ }
For $S=\mathbb{A}^1$, and $K=\C_p$, we deduce
$\mathcal{H}(Z)=\widehat{\left(K[z]_Z\right)}$, which is nothing but the 
M.\ Krasner's original definition of analytic elements on $Z$.

Let $\widehat{K(z)}$ be the completion of $K(z)$ by means of Gauss
norm.
One can interpret its elements as the analytic elements on a generic
disk, and then specialize in the complement $Z\subseteq\mathbb{P}^1$ of
a finite union of disks $\matheur{D}(a_i,1^-)$, $(i=1,2,\ldots)$.
This viewpoint leads us to understand analytic continuation
as a \emph{specialization}. 
\end{para}

\subsection{Limits of algebraic functions; complex case.}

Let $S$ be a Riemann surface coming from a complex affine algebraic
curve $S^{\mathrm{alg}}$, smooth and connected.
Let $\til{S}$ be the universal covering of $S$
(here we tacitly fix a base point $s\in S$).
We endow $\OO(\til{S})$ with the topology of uniform convergence on
every compact; note that the group $\pi_1(S)$ acts (continuously) on
$\OO(\til{S})$.

On the other hand, let $\OO(S^{\mathrm{alg}})^{\mathrm{et}}$ be the
maximal unramified integral extension of $\OO(S^{\mathrm{alg}})$.
Elements in $\OO(S^{\mathrm{alg}})^{\mathrm{et}}$ can be regarded as
unramified algebraic functions on $S^{\mathrm{alg}}$; thus we have
$\OO(S^{\mathrm{alg}})^{\mathrm{et}}\subset\OO(\til{S})$.

\begin{pro}\label{pro-gabber}
$\OO(S^{\mathrm{alg}})^{\mathrm{et}}$ is dense in $\OO(\til{S})$.
\end{pro}

For example, when $S^{\mathrm{alg}}=\mathbb{P}^1\setminus\{0,\infty\}$,
the function $\log z\in\OO(\til{S})$ can be written 
as $\displaystyle \lim_{n\rightarrow \infty}n(z^{1/n}-1)$, 
uniformly on every compact
subset of $\til{S}\simeq\mathbb{A}^1$.

\noindent We can use this formula to compute $\log_{\gamma}1=\lim
n(\zeta_n-1)=2i\pi$, where $\gamma$ is the counter-clockwise loop around
$0$ with the base point $1$.
%PICTURE3
\begin{figure}[h]
 \begin{picture}(70,80)(-35,-40)
  \put(0,-1){\makebox(0,0)[c]{$\bp$}}
  \put(0,-8){\makebox(0,0)[c]{$0$}}
  \put(30,-1){\makebox(0,0)[c]{$\bp$}}
  \put(33,-8){\makebox(0,0)[c]{$1$}}
  \qbezier(0,20)(5,20)(15,15)
  \qbezier(30,0)(25,10)(15,15)
  \qbezier(0,-20)(5,-20)(15,-15)
  \qbezier(30,0)(25,-10)(15,-15)

  \qbezier(-20,16)(-19,18)(-18,20)
  \qbezier(-20,16)(-18,15)(-16,14)

  \qbezier(0,20)(-10,21)(-20,16)
  \qbezier(-30,0)(-30,10)(-20,16)
  \qbezier(0,-20)(-10,-21)(-20,-16)
  \qbezier(-30,0)(-30,-10)(-20,-16)
 \end{picture}
 \caption{}
\end{figure}

\begin{proof}[Proof of \ref{pro-gabber}] (after O. Gabber).
 Consider a countable covering of $\til{S}$
 ($\simeq \mathbb{P}^1(\C),\C$ or $\matheur{D}(0,1^-)$) by
 relatively compact contractible open subsets which are oriented
 manifolds $U_n$ (\eg disks), such that $\ovl{U}_n\subset U_{n+1}$. For
 any $n$, let us consider the set
 \begin{equation*}
 \mathcal{F}_n=\{\gamma \in \pi_1(S,s)\mid
 \gamma\neq 1\text{ and }\gamma\ovl{U}_n \cap \ovl{U}_n \neq \emptyset\}.
 \end{equation*}
 If this set were infinite, we could find a sequence of
 pairwise distinct elements $\gamma_m \in \pi_1(S,s)$ and a sequence of
 points $s_m\in \ovl{U}_n$ such that $\gamma_m s_m\in \ovl{U}_n$. By
 compacity of $\ovl{U}_n$, we might assume that the sequences $s_m$ and
 $\gamma_m s_m$ converge to points $s'$ and $s''$ respectively. Then the
 sequence $\gamma_m s'$ converges to $s''$, which contradicts the
 discreteness of the $\pi_1(S,s)$-orbits in $\til{S}$.  So
 $\mathscr{F}_n$ is finite.

 On the other hand, $S$ is topologically a surface of genus
 $g$ with $t\geq 1$ punctures, so $\pi_1(S)$ is generated by $2g+t$
 generators $\gamma_i$ subject to the relation
 $$[\gamma_1,\gamma_{g+1}]\cdots[\gamma_g,\gamma_{2g}]\gamma_{2g+1}\cdots
 \gamma_{2g+t}=1.$$ Thus $\pi_1(S)$ is free with $2g+t-1$
 generators.  Therefore it is residually finite (\cf\ \cite[p.\ 150, ex.\
 34]{a1}); so we can find a subgroup $\Gamma_n\subseteq\pi_1(S)$ of
 finite index which avoids $\mathscr{F}_n$.

 We immediately see that 
 $\bigl((\Gamma_n\setminus\{1\})\!\cdot\!\ovl{U}_n\bigr)\cap\ovl{U}_n =
 \emptyset$.
 The restriction to $\ovl{U}_n$ of the canonical projection $\til{S}
 \rightarrow S_n := \til{S}/\Gamma_n$ is thus an embedding.
 We identify $\ovl{U}_n$ with its image.

 The part of the exact sequence of cohomologies with compact support
 \begin{equation*}
 \cdots\longrightarrow\mathrm{H}^1(\ovl{U}_n,\Z)\longrightarrow
 \mathrm{H}^2_{\mathrm{c}}(S_n\setminus\ovl{U}_n,\Z)\longrightarrow
 \mathrm{H}^2_{\mathrm{c}}(S_n,\Z)\longrightarrow\cdots
 \end{equation*}
 which can also be written as
 \begin{equation*}
 \ldots \rightarrow
 \mathrm{H}_1(\ovl{U}_n,\Z)=0\rightarrow
 \mathrm{H}_0(S_n\setminus \ovl{U}_n,\Z)\rightarrow
 \mathrm{H}_0(S_n,\Z)=\Z\rightarrow \ldots
 \end{equation*}
 and from which we deduce that $S_n\setminus \ovl{U}_n$ is
 connected.  Because $\partial U_n\subset \ovl{U_{n+1}\setminus
 \ovl{U}_n}$, we have $\ovl{S_n\setminus \ovl{U}_n}=S_n\setminus
 U_n$, hence $S_n\setminus U_n$ is connected.  Since $S_n\setminus U_n$
 is not compact, $U_n$ is holomorphically convex in $S_n$.
 Riemann's existence theorem assures that $S_n$ is the 
 analytification of an affine algebraic curve $S^{\mathrm{alg}}_n$.
 Here we can apply \ref{para-complex-approx} to deduce
 $\OO(U_n)=\widehat{\OO(S^{\mathrm{alg}}_n)^{U_n}}$.
 Note that 
 $\OO(S^{\mathrm{alg}}_n)\subset\OO(S^{\mathrm{alg}})^{\mathrm{et}}$.
 Since every compact set of $\til{S}$ is contained in some $U_n$, we 
 conclude the desired equality 
 $\OO(\til{S})=\widehat{\OO(S^{\mathrm{alg}})^{\mathrm{et}}}$.
\end{proof}

\begin{para}
 It follows from \propref{pro-gabber} that if $\Omega/\Gamma$ is 
 a Schottky partial uniformization of $S$, then
 the intersection $\mathcal{O}(S^{\mathrm{alg}})^{\mathrm{et}}\cap
 \mathcal{O}(\Omega)$ is dense in $\mathcal{O}(\Omega)$.  A variant of this
 statement (with a similar proof) holds, in the $p$-adic case, for a 
 Mumford curve.  We leave it to the reader.
\end{para}

\subsection{Limits of algebraic functions; $p$-adic case.}
\label{subsection-remedy}
\begin{para}
 In the $p$-adic situation, we have seen that there are ``much more''
 \'etale coverings than topological coverings. For instance, the unit disk
 $\matheur{D}=\matheur{D}(0,1^+)$ is simply-connected, but admits many
 non-trivial finite \'etale coverings, \eg the Artin-Schreier covering
 $\matheur{D}\rightarrow \matheur{D}, \;y\mapsto z=y^p-y$.
 
\propref{pro-gabber} suggests to replace, in the $p$-adic case, the ring 
 $\mathcal{O}(\til{S})$, which is often too
 small, by some kind of completion of $\mathcal{O}(S)^{\mathrm{et}}= \varinjlim
 \mathcal{O}(S'),$ where $S'$ runs over the finite \'etale
 connected coverings of $S$.
 
 Let us now assume that $S$ is an affinoid curve with good 
 reduction (hence simply connected). Then there is a
 Gauss $p$-adic norm on $\mathcal{O}(S)$, which extends uniquely to 
 $\mathcal{O}(S)^{\mathrm{et}}$.
 The completion $\widehat{\mathcal{O}(S)^{\mathrm{et}}}$ is
 however pathological in several senses: for instance, 
 it is difficult to give a function-theoretic 
 meaning to its elements, and the continuous derivations of
 $\mathcal{O}(S)$ extends to $\mathcal{O}(S)^{\mathrm{et}}$ but not to
 $\widehat{\mathcal{O}(S)^{\mathrm{et}}}$ in a natural way.
\end{para}

\begin{para}\label{para:r}
 To bypass such problems, we are going to look for some convenient {\it subring} of
 $\widehat{\OO(S)^{\mathrm{et}}}$.
 For simplicity, let us limit ourselves to the following situation:
 \begin{itemize}
  \item $K = {\widehat{\Q}}^{\mathrm{ur}}_p \subset \C_p$
	is the completion of
	the maximal unramified algebraic extension of ${\Q}_p$,
  \item $\mathfrak{v} = {\widehat{\Z}}^{\mathrm{ur}}_p$ is its ring of
	integers (the Witt ring of ${\ovl{\F}}_p$),
  \item $S_0$ is the affine line over ${\ovl{\F}}_p$, deprived from
	finitely many points ${\ovl{\zeta}}_1\ldots,{\ovl{\zeta}}_{\nu}$;
	we choose a point $s_0$ on $S_0$,
  \item $\mathcal{R}$ is the $p$-adic completion of
	$\mathfrak{v}[z, \frac{1}{(z-\zeta_1)\ldots (z-\zeta_{\nu})}]$,
	where $\zeta_1\ldots \zeta_{\nu}$ are liftings of
	${\ovl{\zeta}}_1\ldots,{\ovl{\zeta}}_{\nu}$ in $\mathfrak{v}$; its
	residue ring is $\mathcal{O}(S_0)$,
  \item $S = \matheur{M}\bigl(\mathcal{R}[\frac{1}{p}]\bigr) =
	\matheur{D}(0,1^+) \setminus \bigcup 
	\matheur{D}(\zeta_i,1^-)$, the associated affinoid domain over $K$.
	There is a natural specialization map $\sp: S \rightarrow S_0$ from
	characteristic $0$ to characteristic $p$.
 \end{itemize}

 Let us consider the integral closure ${\mathcal{R}^{\mathrm{et}}}$ of
 $\mathcal{R}$ in ${\mathcal{O}(S)^{\mathrm{et}}}$ and its $p$-adic
 completion $\widehat{\mathcal{R}^{\mathrm{et}}} \subset
 \widehat{\OO(S)^{\mathrm{et}}}$.
 It is thus endowed with (a natural extension of) the $p$-adic
 valuation, and its residue ring is
 $\mathcal{O}(S_0)^{\mathrm{et}}$. Moreover, there is a natural
 structure of ``differential ring'' (better: a connection) on
 $\widehat{\mathcal{R}^{\mathrm{et}}}$ with $\mathfrak{v}$ as ring of
 constants. The ``remarkable equivalence of categories'' of
 A. Grothendieck \cite[18.1.2]{ega44} allows to identify
 $\Aut_{\mathrm{cont}}(\widehat{\mathcal{R}^{\mathrm{et}}}/\mathcal{R})$
 with the algebraic fundamental group $\pi_1^{\mathrm{alg}}(S_0,s_0)$.
\end{para}

\begin{para}\label{para-analytic-cont}
 {\itshape ``Analytic continuation'' as specialization.\ } Let
 $\matheur{D}(s_0,1^-) = \sp^{-1}\lbrace s_0\rbrace$ be the residue
 class of $s_0$ in $S$.  The morphism $\mathcal{R}\rightarrow
 \mathcal{O}(\matheur{D}(s_0,1^-))$ extends non-canonically to a
 continuous morphism $\widehat{\mathcal{R}^{\mathrm{et}}}
 [\frac{1}{p}]\rightarrow \mathcal{O}(\matheur{D}(s_0,1^-))$, determined
 \emph{only up to the action of } $\pi_1^{\mathrm{alg}}(S_0,s_0)$. This
 allows to interpret elements of
 $\widehat{\mathcal{R}^{\mathrm{et}}}[\frac{1}{p}]$ as certain
 multivalued locally analytic functions (in the ``wobbly'' sense). 

The main difference here with Krasner's analytic continuation
 (\cf~\subsecref{subsec-limit-rational}) is the ambiguity coming from
 $\pi_1^{\mathrm{alg}}(S_0,s_0)$.  This provides a kind of multivalued
 analytic continuation, which may be interpreted as analytic
 continuation in characteristic $0$ along an ``\'etale path'' in
 characteristic $p$ (figure \ref{pic4}). 
 
\begin{center}%PICTURE4
 \begin{figure}[h]
  \begin{picture}(200,100)(0,0)
   \put(0,100){\makebox(0,0)[l]{$\matheur{D}(s_0,1^{-})$}}
   \put(0,0){\includegraphics{fig4.eps}}
   \put(15,5){\makebox(0,0)[c]{$s_0$}}
   \put(136,5){\makebox(0,0)[c]{$s_0'$}}
   \put(25,70){\makebox(0,0)[c]{$\times$}}
   \put(135,70){\makebox(0,0)[c]{$\times$}}
   \put(36,66){\makebox(0,0)[c]{$\bbp$}}
   \put(148,66){\makebox(0,0)[c]{$\bbp$}}
  \end{picture}
  \caption{analytic continuation along an underground path}
  \label{pic4}
 \end{figure}
\end{center}

 This is very close to the Robba-Christol theory of ``algebraic
 elements'' \cite{feea}, in which the main player is the \emph{complete}
 subalgebra of $H^{\infty}(\matheur{D}_{\C_p}(0,1^-))$ (the algebra of bounded
 analytic functions) generated the algebraic functions analytic on
 $\matheur{D}_{\C_p}(0,1^-)$. Underlying this ``algebraic'' analytic
 continuation, 
 there is a combinatorics of \emph{automata} (whereas a combinatorics of
 graphs underlies the ``topological'' analytic continuation, as we have
 seen).\footnote{for these and further aspects of the Krasner, resp.
 Christol-Robba, analytic continuation, we refer to \cite{paeap} and to the mimeographed
 notes of the numerous conferences of the Groupe d'\'etude d'analyse
 ultram\'etrique, Paris, devoted to this subject (1973-1980).}
\end{para}

\begin{rem}
 The complex formula $\log z = \lim_n n\bigl(z^\frac{1}{n}-1\bigr)$ has the
 following $p$-adic analogue: $\log z = \lim_n p^{-n}(z^{p^n}-1)$ for $z
 \in \matheur{D}(1,1^-)$. However, the convergence is not uniform on
 $\matheur{D}(1,1^-)$, and the analytic function $\log z$ is not bounded
 on that disk.  Nevertheless $\log : \matheur{D}(1,1^-) \rightarrow
 {\A}^1$ defines an infinite Galois \'etale covering of the $p$-adic
 affine line (not at all a topological covering), with Galois group
 $\mu_{p^{\infty}}$, the $p$-primary torsion in $\C_p$.
\end{rem}

\subsection{$\widehat{\mathcal{R}^{\mathrm{et}}}$ and unit-root $F$-crystals.}
\label{sub-r-and-unit-root}
\begin{para}
 Let us assume for technical simplicity that $p\neq 2$. Let $\mathcal{M}$ be
 a free $\mathcal{R}$-module of finite rank $\mu$, endowed with a connection
 $\nabla: \mathcal{M} \rightarrow
 \Omega_\mathcal{R}\otimes_\mathcal{R}\mathcal{M}$, where
 $\Omega_\mathcal{R}$ is the rank-one module of continuous differentials
 (relative to $\mathfrak{v}$).

 In general, solutions make sense only very locally: typically, analytic
 solutions exist in disks of radius $p^\frac{-1}{p-1}$ (the radius of
 $p$-adic convergence of the exponential function, which is the basic
 example).  A very important criterion, due to B.~Dwork, for the
 convergence of analytic solutions in any open disk of radius $1$, is
 the existence of a so-called Frobenius structure.
\end{para}

\begin{para}\label{para:f-cris}
 {\itshape $F$-crystals.}  Let $\sigma$ be the Frobenius automorphism of
 $\mathfrak{v}$ lifting the $p$th-power map in $\ovl{\F}_p$. There
 are many $\sigma$-linear endomorphisms $\phi$ of $\mathcal{R}$ which reduce
 to the $p$th-power map of $\mathcal{O}(S_0)$ in characteristic
 $p$. For instance, we can take $\phi(z) = z^p$, so that
 $\phi\bigl((z-\zeta_1)\ldots (z-\zeta_{\nu})\bigr) =
 (z^p-\sigma(\zeta_1))\ldots
 (z^p-\sigma(\zeta_{\nu})) \equiv
 \bigl((z-\zeta_1)\ldots (z-\zeta_{\nu})\bigr)^p
 \pmod{p}$. One can show that any lifting $\phi$ extends uniquely to a 
 $\pi_1^{\mathrm{alg}}(S_0,s_0)$-equivariant endomorphism of
 $\widehat{\mathcal{R}^{\mathrm{et}}}$.

 A Frobenius structure on $(\mathcal{M},\nabla)$ is a rule $F$ which
 associates to every lifting $\phi$ a homomorphism $F(\phi):
 \phi^{\ast}\mathcal{M} \rightarrow \mathcal{M}$ (that is to say, a
 $\phi$-linear endomorphism of $\mathcal{M}$) such that:

\begin{enumerate}
 \item $F(\phi)\otimes \Q$ is an \emph{isomorphism},
 \item $F(\phi)$ is \emph{horizontal}, \ie compatible with
       the connections
       $\phi^{\ast}\nabla$ and $\nabla$ respectively,
 \item for any two liftings $\phi$ and $\phi^{\prime}$, the
       homomorphisms $F(\phi)$ and
       $F(\phi^{\prime})$ are related by $F(\phi^{\prime}) = F(\phi)\; \circ
       \;\chi(\phi^{\prime},\phi),$ where
       $\chi(\phi^{\prime},\phi): {\phi^{\prime}}^{\ast}\mathcal{M} \rightarrow
       \phi^{\ast}\mathcal{M} $ is
       the ``Taylor isomorphism'' given by the formulas
       \begin{align*}
	\chi(\phi^{\prime},\phi)({\phi^{\prime}}^{\ast}m) & =
	 \sum_{n\geq 0}{\phi}^{\ast}
	\biggl(\nabla \Bigl(\frac{d}{dz}\Bigr)^n(m)\biggr)
	 ({\phi^{\prime}}(z)-{\phi}(z))^n/n!\\
	& = \sum_{n\geq 0}{\phi}^{\ast} 
	 \biggl(\nabla \Bigl(z\frac{d}{dz}\Bigr)^n(m)\biggr)
	 \biggl(\log\frac{{\phi^{\prime}}(z)}{{\phi}(z)}\biggr)^n\big/n!.
       \end{align*}\label{item:3}
\end{enumerate} 

In virtue of (\ref{item:3}), a Frobenius structure amounts to the data of
the semi-linear horizontal endomorphism $F(\phi)$, for the standard
$\phi(z) = z^p$.

\noindent
The triple $(\mathcal{M},\nabla, F)$ is
called an $F$-{\it crystal}, \cf \cite{tdd}\footnote{the endomorphism
$F_{S_0}$ of $S_0$ given by the $p$th-power map, as well as $\sigma$, $\phi$
and $F$ are all called ``Frobenius'' in the usual jargon, without
causing too much confusion, it seems... }.  The very existence of a
horizontal isomorphism $\phi^{\ast}\mathcal{M}\otimes {\Q} \simeq
\mathcal{M} \otimes {\Q}$ implies that the radius of convergence
${\rho}$ of any local solution of $(\mathcal{M},\nabla)$ satisfies
$\min(1,{\rho})=\min(1,{\rho}^p)$, that is to say: ${\rho}\geq 1$ (since
${\rho}\neq 0$).
\end{para}

\begin{para}\label{para-unit-root}
 {\itshape Unit-root $F$-crystals.}  This is the case where $F(\phi)$
 (not only $F(\phi)\otimes {\Q}$) is an isomorphism (for one, or
 equivalently, for all $\phi$).  The name comes from their first
 appearance in Dwork's computation of the $p$-adic units among the
 reciprocal zeroes of the zeta-function of hypersurfaces in
 characteristic $p$.

 N. Katz has constructed a functor:
 \begin{equation*}
  \left(\begin{array}{l}
   \text{continuous ${\Z}_p$-representations}\\
   \text{of $\pi_1^{\mathrm{alg}}(S_0,s_0)$}
  \end{array}\right)
   \longrightarrow \text{(unit-root $F$-crystals)}
 \end{equation*}
 which associates to any free ${\Z}_p$-module $V$ of rank $r$
 with a continuous action of
 $\pi_1^{\mathrm{alg}}(S_0,s_0)$ a unit-root
 $F$-crystal $(\mathcal{U}_V,\nabla_V, F_V)$ over $\mathcal{R}$ of rank
 $r$.
 Let us recall its definition at the level of objects.
 For any $m\in \N$, we set $S_m=\Spec\;\mathcal{R}/p^m\mathcal{R}$.
 Starting from a representation $\rho$ of $\pi_1^{\mathrm{alg}}(S_0,s_0)$, let
 $G_n$ denote the image of $\rho$ in $GL(V/p^nV)$.  The homomorphism
 $\pi_1^{\mathrm{alg}}(S_0,s_0)\rightarrow G_n$ corresponds to an \'etale 
 covering $S_{n,0}\rightarrow S_0$, which has a
 unique \'etale lifting $\pi_{n,m}:\;S_{n,m}\rightarrow S_m$. The action 
 of $G_n$ on $S_{n,0}$ extends uniquely to
 $S_{n,m}$; on the other hand, the action of $\phi$ on $S_m$ extends 
 uniquely to $S_{n,m}$, and the $\phi$- and
 $G_n$-actions commute. The opposite action makes
 $\mathcal{O}_{S_{n,m}}$ into a right $G_n$-module.
 If we set $\mathcal{U}_n=
 \pi_{n,n_\ast}\mathcal{O}_{S_{n,n}}
 \otimes_{(\mathfrak{v}/p^n\mathfrak{v})[G_n]} V$, 
 we then have a compatible system of isomorphisms
 \begin{equation*}
  \Phi_n=\phi\otimes \id:\;\phi^{\ast}\mathcal{U}_n\rightarrow\mathcal{U}_n
 \end{equation*} 
 and a compatible system of connections
 \begin{equation*}
  \nabla_n=d\otimes \id:\;\mathcal{U}_n\rightarrow
 \pi_{n,n_\ast}\Omega_{S_{n,n}}\otimes_{\mathcal{R}/p^m\mathcal{R}}
 \mathcal{U}_n \simeq \Omega_\mathcal{R}\otimes_\mathcal{R} {\mathcal{U}}_n
 \end{equation*} 
 The unit-root crystal attached to $\rho$ is given 
 by $\mathcal{U}=\varprojlim_n \mathcal{U}_n,
 \nabla=\varprojlim_n \nabla_n,\Phi=\varprojlim_n \Phi_n$.

 The next statement summarizes results of Katz and R.~Crew \cite{fapr}.
\end{para}

\begin{pro}
 \begin{enumerate}
  \item The Katz functor is an equivalence of categories;
  \item any unit-root $F$-crystal $(\mathcal{U},\nabla, F)$ is
	solvable in $\widehat{\mathcal{R}^{\mathrm{et}}}$, \ie :
	$\mathcal{U}\otimes_\mathcal{R}
	\widehat{\mathcal{R}^{\mathrm{et}}}\simeq
	(\mathcal{U}\otimes_\mathcal{R}
	\widehat{\mathcal{R}^{\mathrm{et}}})^{\nabla}\otimes_{\mathfrak{v}}
	\widehat{\mathcal{R}^{\mathrm{et}}}$;
  \item the rule  $\;(\mathcal{U},\nabla, F) \mapsto
	(\mathcal{U}\otimes_\mathcal{R}
	\widehat{\mathcal{R}^{\mathrm{et}}})^{\nabla =0, F(\phi)=\id}$
	(with Galois action coming from that on
	$\widehat{\mathcal{R}^{\mathrm{et}}}$) provides an inverse of
	the Katz functor. 
  \end{enumerate}
\end{pro}
More precisely, we have the action of $\phi$ and $d/dz$ on 
$\widehat{\mathcal{R}^{\mathrm{et}}}$ commute with the action of
$\pi_1^{\mathrm{alg}}(S_0,s_0)$, and we have a canonical isomorphism
\begin{equation*}
 \mathcal{U}\otimes_\mathcal{R}\widehat{\mathcal{R}^{\mathrm{et}}}
  \simeq V\otimes_{\Z_p}\widehat{\mathcal{R}^{\mathrm{et}}}
\end{equation*}
compatible with $\phi, d/dz$ and $\pi_1^{\mathrm{alg}}(S_0,s_0)$
(diagonal action of $\phi, d/dz$ on the left hand side, diagonal action
of $\pi_1^{\mathrm{alg}}(S_0,s_0)$ on the right hand side), which allows
to reconstruct the representation $V$ from the unit-root $F$-crystal and
conversely. In fact, the connection as well as $V$ can be reconstructed
from $(\mathcal{U}, \Phi)$ alone.

This proposition may be compared with \ref{para-cauchy-monodromy}, though it
applies to a quite different type of $p$-adic connections. In
\ref{para-cauchy-monodromy} (for dimension 1), the analytic curve $S$ had
typically bad reduction and the main player was the \emph{discrete
fundamental group} $\pi_1(\Delta)$ of the dual graph of the semistable
reduction, together with the $\pi_1(\Delta)$-module
$\mathcal{O}(\til{S})$. Here the curve $S$ has good reduction
$S_0$ and the main player is the \emph{compact fundamental group}
$\pi_1^{\mathrm{alg}}(S_0)$, together with the
$\pi_1^{\mathrm{alg}}(S_0)$-module $\widehat{\mathcal{R}^{\mathrm{et}}}$.

Before presenting one of Dwork's classical unit-root $F$-crystals, let
us just mention that the above theory extends with little change to the
case when the base ring $\mathfrak{v}$ is a finite ramified extension of
${\widehat{\Z}}^{\mathrm{ur}}_p$: one has to fix an extension of $\sigma$ to
$\mathfrak{v}$, to replace ${\Z}_p$ by $\mathfrak{v}^{\sigma}$, and
to impose some mild nilpotence constraint on $\nabla$ if the
ramification index is $\geq p-1$ (also, it is customary to extend the
definition of Frobenius structure on replacing $\phi$ by some power).

\begin{exa}\label{exa-dwork-exp}
 {\itshape Dwork's exponential.} We denote by
 $\pi$ a fixed $(p-1)$th
 root of $-p$, and set $\mathfrak{v}={\widehat{\Z}}^{\mathrm{ur}}_p[\pi]$, with
 $\sigma(\pi) = \pi$. Let us
 consider the differential equation 
 \begin{equation*}\label{eq:ast}
  f^{\prime}(z)=-\pi f(z) \tag{$\ast$}
 \end{equation*}
 over $S=\matheur{D}(0,1^+)$, which has the analytic solution $f_a =
 e^{-\pi(z-a)}$ in any residue class $\matheur{D}(a,1^-)\subset
 \matheur{D}(0,1^+)$.  The change of variable $z\mapsto z^p$ leads to
 the differential equation

 \begin{equation*}\label{eq:astp}
   g'(z)=-p\pi z^{p-1}g(z) = {\pi}^p z^{p-1}g(z) \tag{$\ast_p$}
 \end{equation*} 

 \noindent
 Dwork's exponential function is 
 \begin{equation*}
  E_{\pi}(z) =
   e^{\pi(z-z^p)} \in {\Z}_p[\pi][[z]].
 \end{equation*}
 This is an invertible element of $\mathcal{R}$, the $\pi$-adic completion of
 $\mathfrak{v}[z]$. This function provides the unit-root Frobenius structure
 which relates (\ref{eq:ast}) and (\ref{eq:astp}).

 Let us reformulate this in the setting of \ref{para-unit-root}: the
 relevant $F$-crystal is $(\mathcal{R}, \nabla, F)$, with $\nabla(1) =
 \pi dz$, $F(\phi)(1) = E_{\pi}(z)$ for the standard $\phi$.  This is
 the unit-root $F$-crystal attached to the finite character of
 $\pi_1^{\mathrm{alg}}({\A}^1_{{\ovl{\F}}_p},0)$ corresponding to the
 Artin-Schreier covering $z=y^p-y.$ The solution $f_0$ of $(\ast)$
 belongs to $\widehat{\mathcal{R}^{\mathrm{et}}}$: indeed, $e^{-\pi z}=
 E_{\pi}(y) \in {\widehat{\mathfrak{v}[y]}} \subset
 \widehat{\mathcal{R}^{\mathrm{et}}}$;
 explicitly, $e^{-\pi z}=\lim_n(1-\pi p^n z)^{p^{-n}}$.

 Dwork has computed the value of his exponential for any $(p-1)$th
 root of unity $\zeta_{p-1}$ (\eg $\zeta_{p-1} =1$): this is

 \begin{equation*}
  E_{\pi}(\zeta_{p-1})=\zeta_p^{\zeta_{p-1}} 
 \end{equation*}

 where $\zeta_p$ is the unique $p$th root of unity $\equiv
 1+\pi\pmod{\pi}$ (\cf \cite[chap. 14]{cf1a2}).  In the spirit of
 \ref{para-analytic-cont}, this may be understood as follows:

 $\zeta_p^{\zeta_{p-1}}=$ \emph{the value at $(z=0)$ of the analytic
 continuation of $f_0$ along the ``wild loop''} corresponding to the
 path from $(y=0)$ to $(y=\zeta_{p-1})$ on the Artin-Schreier covering
 $z=y^p-y$ in characteristic $p$ (figure \ref{pic2}).  If one
 changes $\zeta_{p-1}$, this value $\zeta_p^{\zeta_{p-1}}$ is multiplied
 by some $p$th root of unity.

\begin{figure}[h]%PICTURE5
 \begin{picture}(200,130)(0,0)
%  \put(0,0){\includegraphics[height=4cm,keepaspectratio,clip]{fig5.eps}}
  \put(0,0){\includegraphics[height=4cm,clip]{fig5.eps}}
  \put(40,102){\makebox(0,0)[c]{$\scriptstyle{\zeta_{p-1}}$}}
  \put(48,99.5){$\bbp$}
  \put(43,84){\makebox(0,0)[c]{$\scriptstyle{0}$}}
  \put(48,81){$\bbp$}
  \put(40,55){\makebox(0,0)[r]{$\scriptstyle{y^p-y=x}$}}
  \put(43,23){\makebox(0,0)[c]{$\scriptstyle{0}$}}
  \put(48,20){$\bbp$}
  \put(180,90){\makebox(0,0)[c]{$y$}}
%  \put(180,77){\makebox(0,0)[b]{\small{$-$}}}
%  \put(180,80){\vector(0,-1){50}}
  \put(180,20){\makebox(0,0)[c]{$x$}}
 \end{picture}
 \caption{a wild underground loop}
 \label{pic2}
\end{figure}
%\begin{figure}[h]%PICTURE5
% \setlength{\unitlength}{1pt}
% \begin{picture}(150,100)(0,0)
%  \put(0,0){\includegraphics{pic2.eps}}
%  \put(65,35){$\scriptstyle{y^p-y=x}$}
%  \put(57,-5){$\scriptstyle{x=0}$}
%  \put(57,90){$\scriptstyle{y=\zeta_{p-1}}$}
%  \put(37,57){$\scriptstyle{y=0}$}
% \end{picture}
% \caption{a wild underground roop}
% \label{pic2}
%\end{figure}

In our special case, the discussion of \ref{para-analytic-cont} tells us
that $e^{-\pi z}$ \emph{admits an extension to any disk
$\matheur{D}(a,1^-) \subset \matheur{D}(0,1^+)$, analytic in that disk,
and well-defined up to multiplication by some $p$th root of unity}. It
is easy to find a formula for this extension: it must be proportional to
$e^{\pi (a-z) }$, and by evaluation at $a$, we find that it is $e^{\pi
(a-z)}E_{\pi}(b)$, where $b$ is any solution of the equation $b^p-b=a$.
Its $p$th power is, as expected, $e^{-p\pi z}$ itself.

Therefore, this multivalued exponential $e^{-\pi z}$ on $\matheur{D}(0,1^+)$
appears as a multivalued section of the logarithmic \'etale covering
$\frac{-1}{\pi}\log : \matheur{D}\bigl(1,p^{-\frac{1}{p-1}+}\bigr)
\rightarrow \matheur{D}(0,1^+)$.  This is just the \emph{opposite} way
from the complex situation, where the logarithm is a multivalued
section of the \'etale covering of $\C^{\times}$ given by the exponential.
\end{exa}

\subsection{$p$-adic chiaroscuro: overconvergence.}
\label{sub-chiaroscuro}
\begin{para}
 In fact, Dwork's exponential $E_{\pi}$ is more
 than just an element of the $\pi$-adic completion of $\mathfrak{v}[z]$: it
 is \emph{overconvergent}, \ie extends to an analytic function on a disk
 of radius $\rho > 1$ ($\rho = p^{(p-1)/p^2}$).

 Overconvergence is a fundamental notion in $p$-adic analysis: dimming
 the contours of an affinoid turns out to be the key to the finiteness
 properties of $p$-adic cohomology, as was recognized by Dwork, and
 subsequently developed by P. Monsky, G.  Washnitzer, P. Berthelot
 (rigid cohomology, $\mathcal{D}^{\dagger}$-modules)...
  \begin{figure}[h]%PICTURE6
   \begin{picture}(80,80)(0,0)
    \put(0,0){\includegraphics{fig6.eps}}
    \put(35,35){$\bbp$}
   \end{picture}
   \caption{}
  \end{figure}
\end{para}

\begin{para}
 Let again $S$ be $\matheur{D}(0,1^+)
 \setminus \bigcup \matheur{D}(\zeta_i,1^-)$.  For any
 $\rho >1$, let us consider the bigger affinoids
 $S_{\rho}=\matheur{D}(0,{\rho}^+)
 \setminus \bigcup \matheur{D}(\zeta_i,\frac{1}{\rho}^-)$.  The ring of
 \emph{overconvergent analytic functions} on $S$ is
 
 \begin{equation*}
  \mathcal{H}^{\dagger}(S)= \bigcup_{\rho >1}\mathcal{O}(S_{\rho}).
 \end{equation*}
 Its relevance to the algebraic viewpoint on analytic continuation comes
 from the following result \cite{utdppdfa}:
\end{para}
\begin{pro}
  Let $f \in \mathcal{O}(S)$ satisfy a \emph{monic} polynomial equation
 with coefficients in $\mathcal{H}^{\dagger}(S)$. Then $f \in
 \mathcal{H}^{\dagger}(S)$.
\end{pro}
\begin{para}
 A $F$-crystal $(\mathcal{M},\nabla, F)$ is {\it overconvergent} if
 $(\mathcal{M},\nabla)$ as well as $F(\phi)$ extends over some
 $S_{\rho}$.  This is the case in example \examref{exa-dwork-exp}.  For
 unit-root $F$-crystals, Crew has given the following characterization:

\begin{pro}
 The unit-root $F$-crystal $(\mathcal{U}_V,\nabla_V, F_V )$
 attached to a $p$-adic representation $V$ is overconvergent if and only
 if the images in $GL(V)$ of the inertia groups at the missing points
 ${\ovl{\zeta}}_1\ldots,{\ovl{\zeta}}_{\nu},\infty$ are finite.
\end{pro}

 \end{para}

\subsection{(Overconvergence and Frobenius) Dwork's derivation of
the $p$-adic Gamma function and exponential sums.}
\label{subsection-deriv-of-gamma}

\begin{para}
 We come back to the situation of
 \examref{exa-dwork-exp}.  For any $\alpha \in {\Z}_p$, let us consider
 $M_{\alpha}^{\dagger}:= z^{\alpha}e^{\pi z}\mathcal{H}^{\dagger}(S)$
 endowed with the derivation $z\frac{d}{dz}$. A simple computation shows
 that
 \begin{equation*}
  z^{\alpha}e^{\pi z}z^{k+1}\equiv
   -\frac{\alpha + k}{\pi}z^{\alpha}e^{\pi z}z^k 
 \end{equation*}
 in $M_{\alpha}^{\dagger}/z\frac{d}{dz}M_{\alpha}^{\dagger}$, from which
 one deduces that this cokernel has dimension $1$ over $K$ and is
 generated by $[z^{\alpha}e^{\pi z}]$.

 Following Dwork, let us introduce the operator $\psi$ defined by
 \begin{equation*}
  \psi\Bigl(\sum a_nz^n\Bigr)=\sum a_{pn}z^n.
 \end{equation*}
 This is a left inverse of the Frobenius operator induced by
 $\phi: z \mapsto z^p$.  It acts on $\mathcal{H}^{\dagger}(S)$ and
 commutes with $z\frac{d}{dz}$ up to multiplication by
 $p$: $\psi \circ z\frac{d}{dz}= pz\frac{d}{dz} \circ \psi$.
 The operator $\psi$ can also be applied to any element of
 $M_{\alpha}^{\dagger}$ for $\alpha \in \N$: one finds that
 \begin{equation*}
  \psi(z^{\alpha}e^{\pi z}f)=z^{\beta}e^{\pi z}\psi\bigl(z^{\alpha -
   p\beta}E_{\pi}(z)f\bigr) 
 \end{equation*}
 where $\beta$ is the {\it successor} of $\alpha \in {\Z}_p$, \ie  the
 unique $p$-adic integer $\beta$ such that $p\beta - \alpha \in {\Z}\cap
 [0,p[$ (note that since $\alpha - p\beta >-p\;$, the terms containing a
 negative power of $z$ disappear when applying $\psi$). This formula
 makes sense for any $\alpha \in {\Z}_p$, and we can see that $\psi$
 applies $M_{\alpha}^{\dagger}$ into $M_{\beta}^{\dagger}$, and commutes
 with $z\frac{d}{dz}$. Let us write the induced map of one-dimensional
 cokernels $[{\psi}]:M_{\alpha}^{\dagger} /
 z\frac{d}{dz}M_{\alpha}^{\dagger}\rightarrow
 M_{\beta}^{\dagger}/z\frac{d}{dz}M_{\beta}^{\dagger}$ in the form
 \begin{equation*}
  [{\psi}]([z^{\alpha}e^{\pi z}]) = 
   {\pi}^{p\beta  - \alpha}\Gamma_p(\alpha)[z^{\beta}e^{\pi z}].
 \end{equation*}
\end{para}

\begin{para}
 This function $\Gamma_p$ is then nothing but Morita's $p$-adic
 Gamma function, characterized by its continuity and the functional equation
 \begin{equation*}
  \Gamma_p(0)=1, \;\;\Gamma_p(\alpha +1)/\Gamma_p(\alpha) =
   \begin{cases}
    -\alpha & \text{ if $\alpha$ is a unit,}\\
    -1 & \text{ if $|\alpha|_p < 1$.}
   \end{cases}
 \end{equation*}
 Let us check this functional equation:
 \begin{itemize}
  \item for $\alpha = 0$, one has
	$[z^{k}e^{\pi z}]=0$ if $k>0$; hence
	$\Gamma_p(0)[e^{\pi z}]=[e^{\pi z}][\psi(E_{\pi}(z))]$ $=[e^{\pi z}]$;
  \item if $|\alpha|_p=1$,
	$\beta$ is also the successor of $\alpha +1$, and
	$\Gamma_p(\alpha +1)[z^{\beta}e^{\pi z}] = {\pi}^{\alpha - p\beta
	+1}[{\psi}]([z^{\alpha +1}e^{\pi z}])= {\pi}^{\alpha - p\beta
	+1}[{\psi}](-\frac{\alpha}{\pi}[z^{\alpha}e^{\pi
	z}])=-a\Gamma_p(\alpha)[z^{\beta}e^{\pi z}]$;
  \item finally, if $|\alpha|_p<1$, one has
	$\alpha = p\beta$, and $\beta +1$ is the successor of $\alpha
	+1$; one gets $\Gamma_p(\alpha
	+1)[z^{\beta +1}e^{\pi z}] = {\pi}^{-p+1}[{\psi}]([z^{\alpha +1}
	e^{\pi z}]) = 
	-\frac{1}{p}[{\psi}](-\frac{\alpha}{\pi}[z^{\alpha} e^{\pi z}])
	$ $=
	\frac{\beta}{\pi}\Gamma_p(\alpha)[z^{\beta}e^{\pi z}] =
	- \Gamma_p(\alpha)[z^{\beta +1}e^{\pi z}]$.
 \end{itemize}
\end{para} 

\begin{para}\label{para-cont-of-gamma}
 Let us now check the continuity of $\Gamma_p$ on $\Z_p$, or
 better, its analyticity on any disk
 $\matheur{D}(-k,\vert p\vert^+),\; k=0,1,\ldots,p-1$.
 Let us write $E_{\pi}(z)=e^{\pi(z-z^p)}=\sum_0^{\infty} e_n z^n$.

 For any $\alpha \in \Z_p\cap \matheur{D}(-k,1^-)$, one has
 \begin{align*}
  & \;\psi(z^{\alpha} e^{\pi z})\equiv \pi^k\Gamma_p(\alpha)z^{\beta} e^{\pi z}
  = z^{\beta} e^{\pi z}\psi(z^{-k}E_\pi(z)) \\
  = \;& z^{\beta} e^{\pi z}\sum_{n=0}^{\infty} e_{pn+k} z^n
  \equiv z^{\beta} e^{\pi z}\sum_{n=0}^{\infty} e_{pn+k}(-\pi)^{-n} (\beta)_n,
 \end{align*} 
 where $(\beta)_n=(\frac{\alpha + k}{p})_n$ is the Pochhammer symbol.
 Easy estimates now show that $\Gamma_p(\alpha)=\sum_0^{\infty}
 e_{pn+k}(-\pi)^{-n-k} (\frac{\alpha + k}{p})_n$ is analytic on 
 $\matheur{D}(-k,\vert p\vert^+)$.
\end{para}
\begin{para}\label{para-gross-koblitz}
 {\itshape Gross-Koblitz formula:}
 for $k=0,1,\ldots,p-2$, one has
 \begin{equation*}
  \Gamma_p\biggl(\frac{k}{p-1}\biggr)=-\pi^{-k}\sum_{\zeta_{p-1}}\;
   \zeta_{p-1}^{-k}\zeta_p^{\zeta_{p-1}}\;\;\in \Q_p[\pi]\cap
   {\ovl{\Q}}
 \end{equation*}
 where $\zeta_{p-1}$ runs over the $(p-1)$th roots of unity,
 and $\zeta_p^{\zeta_{p-1}}$ is as before the unique $p$th root of 
 unity $\equiv 1+\zeta_{p-1}\pi \pmod{\pi}$.

 \medskip
 Let us sketch the proof.  We choose $\alpha=\frac{k}{p-1}$, so that
 $\alpha=\beta$ and $\psi$ acts on the ind-Banach-space
 $M_\alpha^{\dagger}$ via the formula:
 $\psi(z^{-k}E_{\pi}(z)f) =
 (z^{\alpha}e^{\pi z})^{-1}\psi(z^{\alpha}e^{\pi z}f)$.
 Coming back to the definition of $\psi$, one observes that for any
 $g\in \mathcal{H}^{\dagger}(S)$, 
 \begin{equation*}
  \psi(z^{-k}g)(z)= \frac{1}{p}\sum_{t\in \phi^{-1}(z)}\;t^{-k}g(t).
 \end{equation*}
 Using the fact that the domain of analyticity of this function is 
 bigger than the domain of analyticity of $g$ itself, one shows that
 $\psi$ is a nuclear operator of $M_\alpha^{\dagger}$. In particular, it
 has a trace, which is the trace of the ``composition operator''
 $\Psi: g\in \mathcal{H}^{\dagger}(S)\mapsto \frac{1}{p}
 \sum_{t\in \phi^{-1}(z)}\;t^{-k}E_\pi(t)g(t)$.
 The computation of this trace is done by approximating $E_\pi$ by
 polynomials and studying the resulting action on the subspace of
 polynomials.  One finds $\Tr\Psi = \frac{1}{p-1}\sum_{\zeta_{p-1}}
 \zeta_{p-1}^{-k}E_\pi(\zeta_{p-1})$.

 At last, because $\psi \circ z\frac{d}{dz}= pz\frac{d}{dz}\circ\psi$,
 one has
 \begin{align*}
  & \pi^k\Gamma_p\biggl(\frac{k}{p-1}\biggr) =
 \Tr\biggl([\psi] \Big|
  M_{\alpha}^{\dag}\big/z\frac{d}{dz}M_\alpha^{\dag}\biggr)
 =(1-p)\Tr\bigl(\psi\big|{M_\alpha^{\dagger}}\bigr)\\
 = & -\sum_{\zeta_{p-1}}\zeta_{p-1}^{-k}E_\pi(\zeta_{p-1})
 = -\sum_{\zeta_{p-1}}\;\zeta_{p-1}^{-k}\zeta_p^{\zeta_{p-1}}.
 \end{align*}

 We refer to \cite{edpaase} for a detailed account. A more general form of the Gross-Koblitz,
proved along the same lines, shows that for any $1\leq k < p^r$, the product 
\begin{equation*}
  \prod_0^{r-1} \Gamma_p\biggl(\frac{p^i k}{p^r-1}\biggr)
 \end{equation*}
belongs to the cyclotomic field $\Q(\zeta_{p^r})$.
  
\end{para}
\begin{para}
 A more straightforward and elementary proof has been discovered
 by A. Robert. It goes as follows.

 As we have just seen, the right hand side of the Gross-Koblitz 
 formula can be written
 \begin{align*}
  &\;-\pi^{-k}\sum_{\zeta_{p-1}}\zeta_{p-1}^{-k}E_\pi(\zeta_{p-1})
  =-\pi^{-k}\sum_{n=0}^{\infty}
  \biggl(\sum_{\zeta_{p-1}}\zeta_{p-1}^{n-k}\biggr) e_n\\
  =&\;(1-p)\pi^{-k}\sum_{m=0}^{\infty} e_{(p-1)m+k}.
 \end{align*}

 Using the expansion of $\Gamma(\alpha)$ given in
 \ref{para-cont-of-gamma}, in the case of $\alpha=\frac{k}{p-1}$, we
 thus have to show that 
 \begin{equation*}
  \sum_{n=0}^{\infty}
   e_{pn+k}(-\pi)^{-n} \biggl(\frac{k}{p-1}\biggr)_n 
   = (1-p)\sum_{m=0}^{\infty} e_{(p-1)m+k}.
 \end{equation*}
 Denote the left hand side by $G_k$ (for any $k\in \N$).
 Due to the overconvergence of $E_\pi$, it is not difficult to see that 
 $\displaystyle\lim_{k\rightarrow \infty}G_k=0$.  One has Robert's formula:
 \begin{equation*}
  G_k-G_{p-1+k}=(1-p)e_k.
 \end{equation*}
 Summing up consecutive expressions, one gets a telescoping sum which yields
 the desired equality $\displaystyle  G_k= G_k-\lim_{m\rightarrow
 \infty}G_{(p-1)m+k} = (1-p)\sum_{m=0}^{\infty} e_{(p-1)m+k}$.
 It remains to prove Robert's formula. One first observes that 
 $z\frac{d}{dz}
 E_\pi=(\pi z-p\pi z^p)E_\pi(z)$, which yields the relation 
 $ne_n=\pi(e_{n-1}-pe_{n-p})$ for $n\geq p$, hence
 $\pi e_{p-1+m}=p\pi e_m + (m+p)e_{m+p}$ for $m\geq 0$. Then 
\begin{align*}
  &\; G_k-G_{p-1+k}\\
 =\;& e_k+ \sum_0^{\infty} e_{p(n+1)+k}(-\pi)^{-n-1}
 \biggl(\frac{k}{p-1}\biggr)_{n+1}\\
 &\phantom{(1-p)e_k+p(1-p)}
 -\sum_0^{\infty} e_{p-1+pn+k}(-\pi)^{-n} \biggl(\frac{k}{p-1}+1\biggr)_n\\
 =\;& e_k+ \sum_0^{\infty}
 \biggl[\frac{k}{p-1}e_{p(n+1)+k}+\pi e_{p-1+pn+k}\biggr]
 (-\pi)^{-n-1} \biggl(\frac{k}{p-1}+1\biggr)_{n}\\
 =\;& e_k + \sum_0^{\infty}
 \biggl[p\frac{(k+(n+1)(p-1))}{p-1}e_{p(n+1)+k}+p\pi e_{pn+k}\biggr]
 (-\pi)^{-n-1} \biggl(\frac{k}{p-1}+1\biggr)_{n}\\
 =\;&(1-p)e_k+p\sum_0^{\infty} \frac{(k+(n+1)(p-1))}{p-1}e_{p(n+1)+k}
 (-\pi)^{-n-1}\biggl(\frac{k}{p-1}+1\biggr)_{n}\\
  &\phantom{(1-p)e_k+p(1-p)}
 - p\sum_1^{\infty} \frac{k+n(p-1)}{p-1} e_{pn+k}(-\pi)^{-n}
 \biggl(\frac{k}{p-1}+1\biggr)_{n-1}\\
 =\;& (1-p)e_k.
\end{align*}
\end{para}
\begin{para}
 The right hand side of the Gross-Koblitz formula is a special case of an
 {\it exponential sum}, \ie an expression of the form
 \begin{equation*}
  S_r(\ovl{f},\ovl{g},\ovl{h})
   = \sum_{x_1,\ldots,x_d \in \F_{p^r},\;
   \ovl{h}(\udl{x})\neq 0}\;\chi\bigl(N_{\F_{p^r}/\F_{p}}
   \ovl{g}(\udl{x})\bigr)\exp\biggl(\frac{2\pi i}{p}
   \Tr_{\F_{p^r}/\F_{p}}\ovl{f}(\udl{x})\biggr) 
 \end{equation*}
 where $\ovl{h}$ is a polynomial with coefficients in
 $\F_{p^r}$, $\ovl{f}$ and $\ovl{g}$ are rational functions with
 coefficients in $\F_{p^r}$ with no pole where $\ovl{h}$
 vanishes, and where $\chi$ is a character of
 $\F_{p}^{\times}$.  This includes Gauss, Jacobi and Kloosterman
 sums as special cases.  In fact, it is known, after the classical works
 of Gauss, Artin, Weil, that counting solutions of systems of polynomial 
 equations in finite fields amounts to the computation of exponential
 sums $S_r(\ovl{f},\ovl{g},\ovl{h})$.  They are studied via their
 generating series, the so-called $L$-series:
 \begin{equation*}
  L(\ovl{f},\ovl{g},\ovl{h};t)=
   \exp\biggl(\sum_{r\geq 1}S_r(\ovl{f},\ovl{g},\ovl{h})\frac{t^r}{r}\biggr).
 \end{equation*}
 Dwork's methods, refined by P. Robba and others, allow to tackle the
 questions of the rationality of $L$, its degree, and of the functional
 equation relating $L(\ovl{f},\ovl{g},\ovl{h};t)$ to
 $L(\ovl{f},\ovl{g},\ovl{h};\frac{1}{p^rt})$.
 In the one-dimensional case, the solution is elementary and follows the
 pattern sketched in \ref{para-gross-koblitz}. Namely \cite{edpaase}:
 \begin{enumerate}
  \item One considers liftings $f,g,h$ of $\ovl{f},\ovl{g},\ovl{h}$ in 
	characteristic zero, and one introduces the
	affinoid set $S$ defined by $\vert h\vert =1$. One sets
	$F=g(z)^{1/p-1}\exp(\pi f(z))$. Then $F'/F$ is a rational function, 
	and one shows that the differential operator
	$\displaystyle \frac{d}{dz} + \frac{F'}{F}$ has an index in
	$\mathcal{H}^{\dagger}(S)$ (in a slightly generalized sense, and
	which can be computed thanks to the work of Robba). This is the 
	crucial point. Hence the cohomology spaces of the de
	Rham complex: $\Omega^0=F\mathcal{H}^{\dagger}(S)\rightarrow 
	\Omega^1=F\mathcal{H}^{\dagger}(S)dz$ are finite-dimensional over
	$\C_p$.
  \item One observes that $E:= F^{\phi-\id}$ (\ie $E(z)= 
	F(z^p)/F(z)$) is an element of $\mathcal{H}^{\dagger}(S)$. This
	allows to define two endomorphisms of the de Rham complex by setting:
	\begin{equation*}
	 \phi^0(F.g)=FEg^{\phi},\;\;\phi^1(Fgdz)
	  =FE\phi'g^{\phi}dz \; \text{ (with $\phi'(z)=pz^{p-1}$)};
	\end{equation*} 
	\begin{equation*}
	 \psi^0(F.g)(z)=F(z)\sum_{\phi(t)=z} \frac{g(t)}{E(t)}\;,\;\;
	  \psi^1(Fgdz)=\sum_{\phi(t)=z} \frac{g(t)}{E(t)\phi'(t)}.
	\end{equation*}
	One has $\psi^{\ast} \circ \phi^{\ast}=p\cdot\id$.  In particular, the
	operators $\mathrm{H}^i(\psi^{\ast}), \;i=0,1,$ are invertible
	on the cohomology spaces. The same argument as in
	\ref{para-gross-koblitz} shows that $\psi^i$ is nuclear on
	$\Omega^i$.  In particular, it has a trace.
  \item Polynomial approximation allows to prove the trace formula:
	\begin{equation*}
	 \tr(\psi^1)-\tr(\psi^0)=\tr(\mathrm{H}^1(\psi^{\ast}))
	  -\tr(\mathrm{H}^0(\psi^{\ast}))=S_r(\ovl{f},\ovl{g},\ovl{h}).
	\end{equation*}
	It follows that
	\begin{equation*}
	 L(\ovl{f},\ovl{g},\ovl{h};t)
	  =\frac{\det(\id-t\mathrm{H}^1(\psi^{\ast}))}
	  {\det(\id-t\mathrm{H}^0(\psi^{\ast}))}.
	\end{equation*}
	This is a rational function since the cohomology is
	finite-dimensional.  Moreover, since the
	$\mathrm{H}^i(\psi^{\ast})$
	are invertible, the computation of its degree amounts to the
	computation of the index of
	$\displaystyle \frac{d}{dz} + \frac{F'}{F}$ (the computation of
	the dimension of $\mathrm{H}^0$, $0 $ or $1$, being essentially
	trivial).
  \item The functional equation follows from a topological ``dual theory''
	in which the transpose of $\phi$ and $\psi$ play the roles of $\psi$
	and $\phi$ respectively.
 \end{enumerate}


 Let us remark at last that this method has an \emph{archimedean
 analogue} in the so-called thermodynamic formalism ---in dimension
 one.  The analogy is especially striking in  
 the presentation of D. Mayer \cite{cfart}:
 exponential sums correspond to ``partition functions'', $L$ to the 
 Ruelle zeta-function, Frobenius to the
 ``shift'', Dwork's operator $\phi^{\ast}$ to the ``transfer operator''
 (or Ising-Perron-Frobenius-Ruelle operator) and is given by a 
 ``composition operator'' (loc. cit. 7.2.2), its nuclearity is
 established by the same argument, and the trace formula has the same
 form (loc. cit. 7.17).

 In view of this close analogy, one could dream of an archimedean 
 proof of the rationality of Weil zeta
 functions, parallel to Dwork's $p$-adic proof...
\end{para}

\newpage
\section{The tale of $F(\frac{1}{2},\frac{1}{2},1;z)$.}\label{sec-tale}

\begin{abst}
This is a detailed concrete illustration of the somewhat abstract
non-archimedean notions discussed in the previous section. We refer to
\cite{hfml} (and \cite[9]{husemoeller87:_ellip}) for an excellent account of
the archimedean tale of $F(\frac{1}{2},\frac{1}{2},1;z)$.
\end{abst}

\subsection{Dwork's hypergeometric function.}\label{sub-dwork-hyper}

Let us consider the Legendre pencil of elliptic
curves with parameter $z\neq 0,1,\infty$, given in inhomogeneous coordinates by
\begin{equation*}
 y^2=x(x-1)(x-z).
\end{equation*}
As a scheme over ${\Z}[z,\frac{1}{{2z(1-z)}}]$, we denote it by $X$. The
first de Rham cohomology module $\mathrm{H}^1_{\mathrm{dR}}(X)$ is free
of rank $2$ over ${\Z}[z,\frac{1}{{2z(1-z)}}]$, and endowed with the
Gauss-Manin connection $\nabla$ (derivation with respect to the
parameter $z$); it is generated by the class $\omega$ of $\frac{dx}{y}$,
and $\nabla(\frac{d}{dz})(\omega)$. The canonical symplectic form
(cup-product) satisfies
$\an{\omega, \nabla(\frac{d}{dz})(\omega)}=\frac{2}{z(z-1)}$. The
Gauss-Manin connection is given by the hypergeometric differential
equation with parameters $(\frac{1}{2},\frac{1}{2},1)$
\begin{equation*}
\nabla(L_{\frac{1}{2},\frac{1}{2},1})(\omega)=0,\text{ with }
L_{\frac{1}{2},\frac{1}{2},1} =
z(1-z)\frac{d^2}{dz^2}+(1-2z)\frac{d}{dz}-\frac{1}{4}.
\end{equation*}
Let $p$ be an odd prime. The {\it Hasse invariant} is the polynomial
$h_p(z) \in {\Z}[\frac{1}{2}][z]$ obtained by truncating the
hypergeometric series $(-1)^\frac{p-1}{2}F(\frac{1}{2},\frac{1}{2},1;z)$
at order $\frac{p-1}{2}$. It enjoys the following well-known properties:
\begin{itemize}
 \item functional equations: $h_p(z)\equiv
       (-1)^\frac{p-1}{2}h_p(1-z)\equiv
       z^\frac{p-1}{2}h_p(\frac{1}{z})\pmod{p}$,
 \item for any $\ovl{\zeta} \in \ovl{\F}_p \setminus \lbrace0,1\rbrace$,
       one has $h_p(\ovl{\zeta})=0$ if and only if the elliptic
       curve $X_{\ovl{\zeta}}$ over $\ovl{\F}_p$ is {\it supersingular}, \ie
       has no geometric point of order $p$,
       \cf \cite[13.3]{husemoeller87:_ellip}.
\end{itemize}
The roots of $h_p \pmod{p}$ are distinct and lie in ${\F}_{p^2}$: we
denote them by $\ovl{\zeta}_1,\ldots,\ovl{\zeta}_{(p-1)/2}$, and choose
liftings ${\zeta}_1,\ldots,{\zeta}_{(p-1)/2}$ in
$\widehat{{\Z}_p^{\mathrm{ur}}}$.

Non-supersingular elliptic curves are called {\it ordinary}; they have
exactly $p$ points of order $p$.

Dwork's hypergeometric function is

\begin{equation*}
f_p(z)=
(-1)^\frac{p-1}{2}\frac{F(\frac{1}{2},\frac{1}{2},1;z)}
{F(\frac{1}{2},\frac{1}{2},1;z^p)}\in {\Z}\Bigl[\frac{1}{2}\Bigr][[z]].
\end{equation*}

One has $f_p(z)\equiv h_p(z)\pmod{p}$. Dwork discovered that although
the $p$-adic radius of convergence of this series is exactly $1$, $f_p$
\emph{does extend to a $p$-adic analytic function on ${\A}^1$ deprived
from the supersingular disks} $\matheur{D}(\zeta_j,1^-)$; we denote this
extension by the same symbol $f_p$.

In terms of this function, he obtained in 1958 his famous $p$-adic 
formula for the number of
rational points of an ordinary elliptic curve defined over ${\F}_{p^n}$:

\emph{for any $s_0\in {\F}_{p^n}, s_0\neq 0,1,\ovl{\zeta}_j,$
 the number of ${\F}_{p^n}$-points of $X_{s_0}$ is}

\begin{equation*}
1 - \prod_{k=0}^{n-1} f_p(\omega^{p^k}) +
 p^n \left(1 - \prod_{k=0}^{n-1} f_p(\omega^{p^k})^{-1}\right),
\end{equation*}
where $\omega$ denotes the unique $(p^n-1)$th root of unity $\equiv s_0
\pmod{p}$.

\subsection{The ordinary unit-root $F$-crystal.}\label{sub-ordinary-unit-root}

On completing $\widehat{{\Z}_p^{\mathrm{ur}}}[z,\frac{1}{{2z(1-z)}}]$
$p$-adically, $(\mathrm{H}^1_{\mathrm{dR}}(X), \nabla)$ gives rise to an
overconvergent $F$-crystal $(\mathcal{H}, \nabla, F)$. We do not discuss
here the construction of the Frobenius structure, which is a general
$p$-adic feature of Gauss-Manin connections. In fact, various analytic
or geometric constructions are available, but in our present case, the
Frobenius structure can be made quite explicit, \cf \cite{pc}.

\medskip
Let us introduce a few notations:
\begin{itemize}
 \item $\mathcal{R}$ = the $p$-adic completion of
       $\widehat{{\Z}_p^{\mathrm{ur}}}\bigl[z,\frac{1}{{h_p(z)}}\bigr]$ 
 \item $\mathcal{R}_{\mathrm{ord}}$ = the $p$-adic completion of
       $\widehat{{\Z}_p^{\mathrm{ur}}}\bigl[z,\frac{1}{{z(1-z)h_p(z)}}\bigr]$
 \item $S=\matheur{M}(\mathcal{R}
       [\frac{1}{p}]) = \matheur{D}(0,1^+) \setminus
       \bigl(\bigcup \matheur{D}(\zeta_j,1^-)\bigr)$
 \item $S_{\mathrm{ord}}=\matheur{M}(\mathcal{R}_{\mathrm{ord}}
       [\frac{1}{p}]) = S\setminus
       \bigl(\matheur{D}(0,1^-)\cup \matheur{D}(1,1^-)\bigr)$
       is the {\it ordinary locus} 
 \item $S_{\mathrm{nss}}= {\A}^1\setminus
       \bigl(\bigcup \matheur{D}(\zeta_j,1^-)\bigr)$
       is the {\it non-supersingular locus}.
\end{itemize}
\noindent The restriction of $(\mathcal{H},\nabla, F)$ to the ordinary
locus possesses a unique non-zero unit-root sub-$F$-crystal
$(\mathcal{U}, \nabla, F)$. This unit-root $F$-crystal extends over $S$
(and even over $S_{\mathrm{nss}}$ as an $F$-isocrystal with logarithmic
singularity at $\infty$ in the sense of \cite{atfff}) \cf also
\cite{ogus00:_ellip}; we use the same
symbol for the extension. It can be described along the following lines:
\begin{itemize}
 \item $\mathcal{U}$ is the unique rank-one horizontal
       submodule of $\mathcal{H}$.
 \item $\mathcal{U}\otimes_{\mathcal{R}_{\mathrm{ord}}}
       \widehat{\mathcal{R}_{\mathrm{ord}}^{\mathrm{et}}} = (\mathcal{H}
       \otimes_{\mathcal{R}_{\mathrm{ord}}}
       \widehat{\mathcal{R}_{\mathrm{ord}}^{\mathrm{et}}})^{\nabla}
       \otimes_{\widehat{{\Z}_p^{\mathrm{ur}}}}
       \widehat{\mathcal{R}_{\mathrm{ord}}^{\mathrm{et}}}.$ (Here
       $\widehat{\mathcal{R}_{\mathrm{ord}}^{\mathrm{et}}}$ is defined
       as in \ref{para:r}.)
 \item for any ordinary $s_0 \in \ovl{\F}_p$, the
       associated $p$-adic representation of
       $\pi_1^{\mathrm{alg}}(S_0,s_0)$ is
       $(\mathcal{H}\otimes_{\mathcal{R}_{\mathrm{ord}}}
       \widehat{\mathcal{R}_{\mathrm{ord}}^{\mathrm{et}}})^{\nabla =0,
       F(\phi)=\id} \simeq \mathrm{H}^1_{\mathrm{et}}(X_{s_0},{\Z}_p).$
 \item For any $u \in (\mathcal{U}\otimes_\mathcal{R}
       \widehat{\mathcal{R}^{\mathrm{et}}})^{\nabla}$, ``the'' image of
       $u$ in $\mathcal{O}(\matheur{D}(s_0,1^-))$ is a bounded solution of the
       differential operator $L_{\frac{1}{2},\frac{1}{2},1}$ on
       $\matheur{D}_{\C_p}(s_0,1^-)$; this also characterizes $\mathcal{U}$.
 \item $\mathcal{U}|_{\matheur{D}(0,1^-)}^{\nabla}$ has a
       canonical $\Z$-submodule, which can be identified with
       $\mathrm{H}^1((X_z)^{\mathrm{an}},{\Z})$ for any $z\in
       \matheur{D}(0,1^-)\setminus \lbrace 0\rbrace$; for a generator
       $u$, one has
       \begin{equation*}
	\an{\omega,u}=\sqrt{-1}F\Bigl(\frac{1}{2},\frac{1}{2},1;z\Bigr)
	 \text{ in } \mathcal{O}(\matheur{D}(0,1^-))
       \end{equation*}
       ($\sqrt{-1}$ appears as residue of $\frac{dx}{y}\big|_{z=0}$ at
       $x=0$), \cf \cite{pbl}.  Similarly,
       $\mathcal{U}|_{\matheur{D}(1,1^-)}^{\nabla}$ has a canonical
       $\Z$-submodule; for a generator $u$, one has
       \begin{equation*}
	\an{\omega,u} = \sqrt{-1}F\Bigl(\frac{1}{2},\frac{1}{2},1;1-z\Bigr)
	 \text{ in } \mathcal{O}(\matheur{D}(1,1^-)).
       \end{equation*}
\end{itemize}

\noindent The existence of the unit-root $F$-crystal $(\mathcal{U},
\nabla, F)$ over $\mathcal{R}$ then amounts to the two
function-theoretic facts:
\begin{equation*}
 \dlog F\Bigl(\frac{1}{2},\frac{1}{2},1;z\Bigr)\;\text{ and }\;
 f_p(z) =
 \Bigl(\sqrt{-1}F\Bigl(\frac{1}{2},\frac{1}{2},1;z\Bigr)\Bigr)^{1-\sigma}
\end{equation*}
both extend to units in $\mathcal{R}$. 

The inertia at any supersingular point maps \emph{onto} $\Aut
\mathrm{H}^1_{\mathrm{et}}(X_{s_0},{\Z}_p)\simeq {\Z}_p^{\times}$ (J.-I. 
Igusa, \cf \eg \cite{anopu}).  The inertia at $\infty$ acts as $\pm1$.
According to Crew's criterion, the $F$-crystal $(\mathcal{U},\nabla,F)$
is \emph{not} overconvergent.

We refer to \cite{opm} for a study in the same spirit of more general
hypergeometric equations and other differential equations ``coming from
geometry''.

\subsection{Analytic continuation: a pr\'ecis.}\label{sub-analytic-cont}

\begin{para}
 We have seen that Dwork's hypergeometric function extends to an analytic
 function on the whole of $S_{\mathrm{nss}}$. Some of its special values
 have been computed.  Let us mention \cite{thfwpp}, \cite{anjsasvogphf}
 \begin{enumerate}
 \item $f_p(1)=1$ (Koblitz),
 \item if $p \equiv 1\pmod{4}$,
       \[
       f_p(-1) =
       (-1)^\frac{p-1}{4}\frac{\Gamma_p(1/4)^2}{\Gamma_p(1/2)} =
       \frac{\Gamma_p(1/4)}{\Gamma_p(1/2)\Gamma_p(3/4)} \quad\text{ (Young).}
       \]
 \end{enumerate}
 (the condition $p \equiv 1\pmod{4}$ ensures that $-1$ is an
 ordinary modulus); $f_p(-1)$ is a Gauss integer (Van Hamme).
\end{para}

\begin{para}
 We have seen that the logarithmic derivative
 $\dlog F(\frac{1}{2},\frac{1}{2},1;z) $ also extends to an analytic
 function on $S_{\mathrm{nss}}$.  It satisfies the functional equations
 \begin{align*}
  \dlog F\Bigl(\frac{1}{2},\frac{1}{2},1;z\Bigr)
   & = -\dlog F\Bigl(\frac{1}{2},\frac{1}{2},1;1-z\Bigr)\\
   & = -\frac{1}{2z} - \frac{1}{z^2}
   \dlog F\Bigl(\frac{1}{2},\frac{1}{2},1;\frac{1}{z}\Bigr).
 \end{align*}
\end{para}

\begin{para}
 Let us turn to the more subtle case of $F(\frac{1}{2},\frac{1}{2},1;z)$
 itself. The general discussion of
 \ref{para-analytic-cont}/\ref{para-unit-root} applies to the ordinary
 unit-root $F$-crystal and tells us that $F(\frac{1}{2},\frac{1}{2},1;z)$
 \emph{admits an extension to any unit disk $\matheur{D}(s,1^-) \subset S$,
 analytic in that disk, and well-defined up to multiplication by a unit
 in ${\Z}_p$}. In other words, $\;F(\frac{1}{2},\frac{1}{2},1;z)$ is the
 specialization of an element (in fact, a unit) of
 $\widehat{\mathcal{R}^{\mathrm{et}}_{\mathrm{nss}}}$.
 It may be suggestive to denote
 such an element by $\;F(\frac{1}{2},\frac{1}{2},1;\eta)$, where $\eta$
 stands for the Berkovich generic point of $S$ corresponding to the
 sup-norm).

 \noindent In $\matheur{D}(1,1^-)$, this extension is
 $F(\frac{1}{2},\frac{1}{2},1;1-z)$ (up to $\Z_p^{\times}$). The
 refined structure of $F$-isocrystal with logarithmic singularity at
 $\infty$ allows to construct an extension of
 $F(\frac{1}{2},\frac{1}{2},1;z)^2$ to $\matheur{D}(\infty,1^-)$,
 analytic in that disk: it is $\frac{1}{z} F(\frac{1}{2},\frac{1}{2},1;\frac{1}{z})^2$ (up to $\Z_p^{\times}$); notice the
 square, which comes from the exponent $\pm 1/2$ at $\infty$.
\end{para}

\begin{para}
 The Krasner representation of $f_p(z)$ and $\dlog
 F(\frac{1}{2},\frac{1}{2},1;z)$ as uniform limits of rational functions
 without pole (nor zero) on $S$ can be made explicit: for $n>0$, let
 $g_n$ be the polynomial obtained by truncating
 $F(\frac{1}{2},\frac{1}{2},1;z)$ at order $p^n-1$; then $f_p(z)= \lim
 (-1)^{\frac{p-1}{2}}\frac{g_{n+1}(z)}{g_n(z^p)}$ and $\dlog
 F(\frac{1}{2},\frac{1}{2},1;z)=\lim\dlog g_n(z)$ \cite[3.4]{pc}.  On
 the other hand, $F(\frac{1}{2},\frac{1}{2},1;z)$ cannot be approximated
 by rational functions on $S$, but $L_{\frac{1}{2},\frac{1}{2},1}(g_n)$
 tends uniformly to $0$ on $\matheur{D}(0,1^+)$ \cite{sbdsdlaafh}.  One
 has

\begin{equation*}
\sqrt{-1}F\Bigl(\frac{1}{2},\frac{1}{2},1;z\Bigr) \equiv
 \bigl(h_p(z)\bigr)^{-\frac{1}{p-1}}
\pmod{p}
\end{equation*}
but we do not know any explicit representation of 
$F(\frac{1}{2},\frac{1}{2},1;z)$ as a $p$-adic limit of algebraic
functions on $S$.
\end{para}

\begin{para}
 Any solution of $L_{\frac{1}{2},\frac{1}{2},1}$ on an ordinary disk
 $\matheur{D}(s,1^-)$, which is not proportional to (the extended)
 $F(\frac{1}{2},\frac{1}{2},1;z)$, is not bounded in
 $\matheur{D}_{\C_p}(s,1^-)$,
 hence does not extend to neighboring disks in any reasonable
 \emph{analytic} sense. In spite of this obstruction to continuation,
 one can nevertheless ``jump from disk to disk'' (see figure \ref{fig7})
 \begin{figure}[h]%PICTURE7
  \begin{picture}(300,100)(0,-10)
   \put(0,0){\includegraphics{fig7.eps}}
   \put(10,0){{\scriptsize $\matheur{D}(\zeta_{\mathrm{can}},1^-)$}}
  \end{picture}
  \caption{}
  \label{fig7}
 \end{figure}
 and extend them as locally analytic functions, in the following way.

 For each $\ovl{\zeta} \in S_0 \setminus \lbrace 0,1\rbrace$,
 there is a {\it canonical lifting} $\;\zeta_{\mathrm{can}} \in
 {\Z}_p^{\mathrm{ur}}$: namely, the modulus $\zeta_{\mathrm{can}}$ for which
 $X_{\zeta_{\mathrm{can}}}$ has complex multiplication by the quadratic
 order $\End(X_{\ovl{\zeta}})$.  Let us fix a branch of
 $F(\frac{1}{2},\frac{1}{2},1;z)$ in
 $\matheur{D}(\zeta_{\mathrm{can}},1^-)$, \ie an analytic solution
 $F_{\ovl{\zeta}}\in \widehat{{\Z}_p^{\mathrm{ur}}}[[z-{\zeta_{\mathrm{can}}}]]$
 of $L_{\frac{1}{2},\frac{1}{2},1}$ in
 $\matheur{D}(\zeta_{\mathrm{can}},1^-)$ which is a specialization of
 $F(\frac{1}{2},\frac{1}{2},1;\eta)\in
 (\widehat{\mathcal{R}^{\mathrm{et}}})^{\times}$ (any other branch is of
 the form $c.F_{\ovl{\zeta}}$ with $c \in {\Z}_p^{\times}$. Let
 $u_{\ovl{\zeta}} \in
 \mathcal{U}|_{\matheur{D}(\zeta_{\mathrm{can}},1^-)}^{\nabla}$ be
 defined by

 \begin{equation*}
 \an{\omega,u_{\ovl{\zeta}}}=\sqrt{-1}F_{\ovl{\zeta}}\;\text{ in }\;
 \mathcal{O}(\matheur{D}({\zeta_{\mathrm{can}}},1^-)).
 \end{equation*}
 On the other hand, let $v_{\ovl{\zeta}} \in
 \mathcal{M}|_{\matheur{D}(\zeta_{\mathrm{can}},1^-)}^{\nabla}$ be defined by
 \begin{equation*}
 \sqrt{-1} F_{\ovl{\zeta}}({\zeta_{\mathrm{can}}}).v_{\ovl{\zeta}}
  = \omega({\zeta_{\mathrm{can}}}),
  \;\text{ so that }\; \an{v_{\ovl{\zeta}},u_{\ovl{\zeta}}}=1.
 \end{equation*}
 Then $\an{\omega,v_{\ovl{\zeta}}}$ defines an
 unbounded solution of
 $L_{\frac{1}{2},\frac{1}{2},1}$ in
 $\matheur{D}_{\C_p}(\zeta_{\mathrm{can}},1^-)$.
\end{para}

\begin{para}
 Finally, in any supersingular disk $\matheur{D}_{{\C}_p}(\zeta_{i},1^-)$,
 there is no non-zero bounded solution of $L_{\frac{1}{2},\frac{1}{2},1}$.
\end{para}

\subsection{The Tate and Dwork-Serre-Tate parameters.}
\label{sub-serre-tate-para}

\begin{para}\label{para-tate-para}
 It turns out that the formal group
 ${\widehat{X}}_{(\widehat{\mathcal{R}_{\mathrm{ord}}^{\mathrm{et}}})}$
 of $X$ over $\widehat{\mathcal{R}_{\mathrm{ord}}^{\mathrm{et}}}$ is
 isomorphic to ${\widehat{\G}}_m$ (\cf~\cite[I]{fgade}). Any such
 isomorphism transforms the canonical differential
 $\omega_{\mathrm{can}}$ on ${\widehat{\G}}_m$ into
 $\Theta \frac{dx}{y}$, for a suitable element $\Theta \in
 \widehat{\mathcal{R}_{\mathrm{ord}}^{\mathrm{et}}}$ well-defined up to
 multiplication by an element of ${\Z}_p^{\times}$. It is possible to
 arrange normalizations so that the following relation hold (loc. cit.):
 
 \begin{equation*}
  \Theta
   = \biggl(\sqrt{-1}F\Bigl(\frac{1}{2},\frac{1}{2},1;\eta\Bigr)\biggr)^{-1}.
 \end{equation*}
 Let $\zeta\in \widehat{{\Z}_p^{\mathrm{ur}}}, \zeta \neq 0,1,$ be a point of
 $S$. By specialization, the element $\Theta(\zeta) \in
 \widehat{{\Z}_p^{\mathrm{ur}}}^{\times}$ (well-defined up to
 ${\Z}_p^{\times}$) links up $\omega_{\mathrm{can}}$ and $\frac{dx}{y}$ on
 $\widehat{X}_{\zeta}\simeq{\widehat{\G}}_m$; it is called the {\it Tate
 parameter} or {\it Tate constant} of $X_{\zeta}$.  It is well-defined
 up to multiplication by an element of ${\Z}_p^{\times}$.

 This is the $p$-adic analogue of the following familiar situation: let
 us consider the two-step uniformization of a complex elliptic curve
 \begin{equation*}
  \C \xrightarrow{\exp\bigl(\frac{2i\pi}{\omega_1}\cdot \;\,\bigr)}
   {\C}^{\times}
   \longrightarrow \C^{\times}/q^{\Z} \simeq
   {\C}/(\omega_1{\Z}+\omega_2{\Z})
 \end{equation*}
 with $q=\exp(2i\pi\frac{\omega_2}{\omega_1})$.  If the elliptic curve
 is $X_{\zeta}$ with $\zeta\in \matheur{D}(0,1^-)\setminus \lbrace
 0\rbrace$, and if $\omega_1$ and $\omega_2$ are fundamental periods of
 $\frac{dx}{y}$ ($\omega_1$ being the period attached to the vanishing
 cycle), then $\frac{2i\pi}{\omega_1}$ appears as the analogue of
 $\Theta(\zeta)$, and it is well-known that $\omega_1/2i\pi =
 iF(\frac{1}{2},\frac{1}{2},1;\zeta)$ (up to sign).
\end{para}

\begin{para}\label{para-dwork-para}
 We come back to the symplectic basis $v_{\ovl{\zeta}}, u_{\ovl{\zeta}}$
 and write $\frac{1}{{\sqrt{-1} F_{\ovl{\zeta}}}} \omega$ in the form
 $v_{\ovl{\zeta}} + \tau u_{\ovl{\zeta}}$, where

 \begin{equation*}
 \tau = -\frac{\an{\omega,v_{\ovl{\zeta}}}}{\an{\omega,u_{\ovl{\zeta}}}}
  \in
  (z-{\zeta_{\mathrm{can}}})\widehat{{\Q}_p^{\mathrm{ur}}}[[z-{\zeta_{\mathrm{can}}}]].
 \end{equation*}
 This defines an unbounded element of
 $\mathcal{O}(\matheur{D}_{{\C}_p}(\zeta_{\mathrm{can}},1^-))$ (notice that
 another choice of $u_{\ovl{\zeta}}$ multiplies $\tau$ by an element of
 $\Z_p^{\times}$).

 \noindent Applying $\nabla(d/dz)$ to
 $\frac{1}{{\sqrt{-1}F_{\ovl{\zeta}}}} \omega = v_{\ovl{\zeta}} + \tau 
 u_{\ovl{\zeta}}$ and using $\lan{\omega,\nabla(\frac{d}{dz})(\omega)} =
 \frac{2}{z(z-1)}$, one derives:
 \begin{equation*}
 \frac{d\tau}{dz} = \frac{2}{z(1-z) F_{\ovl{\zeta}}^2}.
 \end{equation*}
 The exponential of $\tau$ is called the {\it Dwork-Serre-Tate
 parameter}.  We shall recall later its meaning as a parameter for
 $p$-divisible groups. It satisfies
 the following remarkable integrality property \cite[Th.4]{pc}, \cite{tdd}
 \begin{equation*}
 q = e^{\tau} \in 1 + (z-{\zeta_{\mathrm{can}}})
 \widehat{{\Z}_p^{\mathrm{ur}}}[[z-{\zeta_{\mathrm{can}}}]].
 \end{equation*}
 This suggests the following question: does $q$ arise from an element of
 $\widehat{\mathcal{R}_{\mathrm{ord}}^{\mathrm{et}}}\;$? The answer is
 \emph{no}. This may be seen by considering $q\pmod{p}$: the 
 formula for $d\tau/dz = \dlog(q)$ shows that $q\pmod{p}$ is non-constant
 ($\dlog(q) \equiv \frac{2(h_p)^{2/p-1}}{z(z-1)}\pmod{p}$).  If $q$ comes
 from an element of $\widehat{\mathcal{R}_{\mathrm{ord}}^{\mathrm{et}}}$ which
 specializes to $1$ at every canonical modulus $\zeta_{\mathrm{can}}$, one has 
 $q\equiv 1\pmod{p}$, a contradiction.
\end{para}

\subsection{Complex counterpart: the supersingular locus.}
\label{sub-complex-supersing}

 Let us now revert things and try to understand the complex situation
 from the $p$-adic viewpoint!

 Let $\matheur{D}^+$ and $\matheur{D}^-$ be the (complex-conjugate)
 connected components of $\matheur{D}(\frac{1}{2},\frac{3}{2}^-)\setminus
 \bigl(\matheur{D}(\frac{-1}{4},\frac{3}{4}^+) \cup
 \matheur{D}(\frac{5}{4},\frac{3}{4}^+)\bigr)$. Let $S$ 
 be the complement of $\matheur{D}^+ \cup \matheur{D}^-$ in the complex plane.
 Note that $S$ is closed and arcwise connected, and that
 $\pi_1^{\mathrm{top}}(S)$ 
 is a free group with two generators $\gamma^+,\gamma^-$ (see figure
 \ref{pic3}). Its 
 interior $S^{\circ}$ has three connected (simply-connected) components
 $\matheur{D}_0= \matheur{D}(\frac{-1}{4},\frac{3}{4}^-),\; \matheur{D}_1=
 \matheur{D}(\frac{5}{4},\frac{3}{4}^-)\;$ and $\matheur{D}_{\infty}$.

\begin{figure}[h]
 \begin{picture}(150,180)(0,0)
  \put(0,0){\includegraphics[clip]{fig8.eps}}
  \put(69,105){$\matheur{D}^{+}$}
  \put(69,30){$\matheur{D}^{-}$}
  \put(175,90){$\matheur{D}_{\infty}$}
  \put(42,82){$\matheur{D}_0$}
  \put(98,82){$\matheur{D}_1$}
  \put(0,58){$-1$}
  \put(55,58){$0$}
  \put(55,66.7){\makebox(0,0)[c]{$\bp$}}
  \put(80,58){$\FRAC{1}{2}$}
  \put(95,58){$1$}
  \put(95,66.7){\makebox(0,0)[c]{$\bp$}}
  \put(138,58){$2$}
  \put(145,108){$\gamma^{+}$}
  \put(150,40){$\gamma^{-}$}
 \end{picture}
 \caption{}
 \label{pic3}
\end{figure}

 Let $\mathcal{O}(S)$ denote the ring of continuous functions on $S$
 analytic in $S^{\circ}$. According to Mergelyan's theorem
 \cite{uatfoacv}, they are uniform limits of rational functions on every
 compact $K\subset S$ such that $\pi_0(\C\setminus K,0)$ is finite.
 This ring is however not stable under differentiation, and we consider
 its differential closure $\mathcal{R}$ in $\mathcal{O}(S^{\circ})$. We
 also consider the integral closure $\mathcal{R}^{\mathrm{et}}$ of
 $\mathcal{R}$ in $\mathcal{O}(S^{\circ})$. Every element of
 $\mathcal{R}^{\mathrm{et}}$ defines a multivalued locally analytic
 function on $S$, \ie an analytic germ which may be analytically
 continued along any path of $S$ not ending at $-1,\frac{1}{2}, 2$, in
 such a way that the germ at the other extremity is analytic.

 For instance, $F(\frac{1}{2},\frac{1}{2},1;z)$ may be viewed as an element
 of $\mathcal{R}^{\mathrm{et}}$: in fact,
 $F(\frac{1}{2},\frac{1}{2},1;z)^4 \in \mathcal{O}(S)$. An explicit
 continuous extension of $F(\frac{1}{2},\frac{1}{2},1;z)^4$ from
 $\matheur{D}_0$ to $S$ is given by $F(\frac{1}{2},\frac{1}{2},1;1-z)^4$
 in $\matheur{D}_1$, $\frac{1}{z^2 }
 F(\frac{1}{2},\frac{1}{2},1;\frac{1}{z})^4$ in $\matheur{D}_{\infty}$,
 and the values
 $\frac{{\Gamma(1/4)}^8}{64\pi^6},\;\frac{{\Gamma(1/4)}^8}{16\pi^6} ,
 \;\frac{{\Gamma(1/4)}^8}{64\pi^6}$ at $-1,\;\frac{1}{2},\;2$
 respectively \cite[p.189]{f2}.  On the other hand, note that $\dlog
 F(\frac{1}{2},\frac{1}{2},1;z)$ is in $\mathcal{R}$ but not
 in $\mathcal{O}(S)$.

%Let us now revert things and try to understand the complex situation
%from the $p$-adic viewpoint!
%
%Let $\matheur{D}^+$ and $\matheur{D}^-$ be the (complex-conjugate)
%connected components of $\matheur{D}(\frac{1}{2},\frac{1}{2}^-)\setminus
%\bigl(\matheur{D}(\frac{1}{4},\frac{1}{4}^+) \cup
%\matheur{D}(\frac{3}{4},\frac{1}{4}^+)\bigr)$. Let $S$ be the complement of
%$\matheur{D}^+ \cup \matheur{D}^-$ in the complex plane. Note that $S$
%is closed and arcwise connected, and that $\pi_1^{\mathrm{top}}(S)$ is a
%free group with two generators $\gamma^+,\gamma^-$ (see figure
%\ref{pic3}).  Its interior $S^{\circ}$ has three connected
%(simply-connected) components $\matheur{D}_0=
%\matheur{D}(\frac{1}{4},\frac{1}{4}^-)$, $\matheur{D}_1=
%\matheur{D}(\frac{3}{4},\frac{1}{4}^-)$ and $\matheur{D}_{\infty}$
%(see figure \ref{pic3}).
%
%\begin{figure}[h]
% \begin{picture}(150,180)(0,0)
%  \put(0,0){\includegraphics[clip]{fig8.eps}}
%  \put(69,105){$\matheur{D}^{+}$}
%  \put(69,30){$\matheur{D}^{-}$}
%  \put(175,90){$\matheur{D}_{\infty}$}
%  \put(42,60){$\matheur{D}_0$}
%  \put(100,60){$\matheur{D}_1$}
%  \put(26,62){$0$}
%  \put(82,62){$\FRAC{1}{2}$}
%  \put(133,58){$1$}
%  \put(145,108){$\gamma^{+}$}
%  \put(160,40){$\gamma^{-}$}
% \end{picture}
% \caption{}
% \label{pic3}
%\end{figure}
%\begin{figure}[h]%PICTURE8
%\setlength{\unitlength}{1pt}
%\begin{picture}(150,150)(0,0)
%\put(0,0){\includegraphics{pic3.eps}}
%\put(67,100){$\matheur{D}^{+}$}
%\put(67,30){$\matheur{D}^{-}$}
%\put(40,80){$\matheur{D}_0$}
%\put(95,80){$\matheur{D}_1$}
%\put(0,65){$-1$}
%\put(50,54){$0$}
%\put(88,54){$\displaystyle\frac{1}{2}$}
%\put(133,65){$1$}
%\put(140,103){$\gamma^{+}$}
%\put(140,12){$\gamma^{-}$}
%\put(0,130){$\matheur{D}_{\infty}$}
%\end{picture}
%\caption{}
%\label{pic3}
%\end{figure}
%
%Let $\mathcal{O}(S)$ denote the ring of continuous functions on $S$
%analytic in $S^{\circ}$. According to Mergelyan's theorem
%\cite{uatfoacv}, they are uniform limits of rational functions on every
%compact $K\subset S$ such that $\pi_0(\C\setminus K,0)$ is finite.  This
%ring is however not stable under differentiation, and we consider its
%differential closure $\mathcal{R}$ in $\mathcal{O}(S^{\circ})$. We also
%consider the integral closure $\mathcal{R}^{\mathrm{et}}$ of
%$\mathcal{R}$ in $\mathcal{O}(S^{\circ})$. Every element of
%$\mathcal{R}^{\mathrm{et}}$ defines a multivalued locally analytic
%function on $S$, \ie an analytic germ which may be analytically
%continued along any path of $S$ not ending at $0,\frac{1}{2}, 1$, in
%such a way that the germ at the other extremity is analytic.  \par For
%instance, $F(\frac{1}{2},\frac{1}{2},1;z)$ may be viewed as an element
%of $\mathcal{R}^{\mathrm{et}}$: in fact,
%$F(\frac{1}{2},\frac{1}{2},1;z)^4 \in \mathcal{O}(S)$. An explicit
%continuous extension of $F(\frac{1}{2},\frac{1}{2},1;z)^4$ from
%$\matheur{D}_0$ by
%$S$ is given by $F(\frac{1}{2},\frac{1}{2},1;1-z)^4$ in $\matheur{D}_1$,
%$\frac{\pi^4}{z^2 \log^4 16}F(\frac{1}{2},\frac{1}{2},1;\frac{1}{z})^4$ in
%$\matheur{D}_{\infty}$, and the values $1,\frac{{\Gamma(1/4)}^8}{16\pi^6},1$ at
%$0, \frac{1}{2}, 1$ respectively. On the other hand, note that
%$\dlog\;F(\frac{1}{2},\frac{1}{2},1;z)$ is in $\mathcal{R}$
%but not in $\mathcal{O}(S)$.

The de Rham cohomology of the Legendre elliptic pencil gives
rise to an $\mathcal{R}[1/z(1-z)]$-module with connection
$(\mathcal{H},\nabla)$.  It admits a \emph{unique} non-zero horizontal
submodule $\mathcal{U}$ which extends to $S$; we use the same
symbol for the extension. It can be described along the following lines:

\begin{itemize}
 \item For any $u \in (\mathcal{U}\otimes_\mathcal{R}
       {\mathcal{R}^{\mathrm{et}}})^{\nabla}$, ``the'' image of $u$ in
       $\mathcal{O}(\matheur{D}_j) (j=0,1, \text{ or } \infty),$ is a
       bounded solution of the differential operator $L_{\frac{1}{2},\frac{1}{2},1}$.
 \item $\mathcal{U}\otimes_\mathcal{R}
       {\mathcal{R}^{\mathrm{et}}[\frac{1}{z(1-z)}]} =
       \bigl(\mathcal{H}\otimes_{\mathcal{R}[\frac{1}{z(1-z)}]}
       {\mathcal{R}^{\mathrm{et}}[\frac{1}{z(1-z)}]}\bigr)^{\nabla}\otimes_\C
       {\mathcal{R}^{\mathrm{et}}[\frac{1}{z(1-z)}]}.$
 \item $\bigl(\mathcal{H}\otimes_{\mathcal{R}[\frac{1}{z(1-z)}]}
       {\mathcal{R}^{\mathrm{et}}[\frac{1}{z(1-z)}]}\bigr)^{\nabla}$ has a
       canonical ${\Z}[i]$-submodule ($i=\sqrt{-1}$) which can be
       locally identified with the part
       $\mathrm{H}^1(X_z^{\mathrm{an}},{\Z}[i])_{\mathrm{isotriv}}$ of
       $\mathrm{H}^1(X_z^{\mathrm{an}},{\Z}[i])$ where
       $\pi_1^{\mathrm{top}}(S,z)$ acts through a finite group.
 \item $\mathcal{U}|_{\matheur{D}_0}^{\nabla}$ has a
       canonical $\Z$-submodule, which can be identified with the part
       of $\mathrm{H}^1(X_z^{\mathrm{an}},{\Z}[i])_{\mathrm{isotriv}}$ invariant
       under complex conjugation. One of the two generators $u$ satisfies
       \begin{equation*}
	\Lan{\omega,\frac{u}{2i\pi}}=
	 i F\Bigl(\frac{1}{2},\frac{1}{2},1;z\Bigr)
	 \;\text{ in }\; \mathcal{O}(\matheur{D}_0).
       \end{equation*}
       Similarly, $\mathcal{U}|_{\matheur{D}_1}^{\nabla}$ has a canonical
       $\Z$-submodule, and for one of the two generators $u$, one has
       \begin{equation*}
	\Lan{\omega,\frac{u}{2i\pi}} =
	 i F\Bigl(\frac{1}{2},\frac{1}{2},1;1-z\Bigr)\;\text{ in }\;
	 \mathcal{O}(\matheur{D}_1).
       \end{equation*}
       Note the occurrence of $2i\pi$ in these formulas. 
 \item The local monodromy $\gamma^+$ around $\matheur{D}^+$
       maps \emph{onto}
       \[\Aut \mathrm{H}^1(X_{s},{\Z}[i])_{\mathrm{isotriv}}\simeq
       {\Z}/4{\Z}.\]
       Same for $\gamma^-$.
       The local monodromy at $\infty$ acts as $\pm1$.
 \item In $\matheur{D}^+$ and in $\matheur{D}^-$, there is no non-zero
       bounded analytic solution
       of $L_{\frac{1}{2},\frac{1}{2},1}$.
\end{itemize}


This is quite similar to \subsecref{sub-ordinary-unit-root},
\subsecref{sub-analytic-cont}, $\matheur{D}^+\cup \matheur{D}^-$ playing
the role of the \emph{supersingular locus}. This picture is assuredly
very different from the traditional view of monodromy for
$F(\frac{1}{2},\frac{1}{2},1;z)$, and may shed some light upon the
divergences between the topological and algebraic approaches to analytic
continuation in the $p$-adic case (here of course, the coexistence of
the two pictures is explained by the fact that the principle of unique
continuation (\defref{dfn-unique-cont}) fails for the sheaf
of germs of continuous functions on $S$, analytic on $S^{\circ}$).

The tale of $F(\frac{1}{2},\frac{1}{2},1;z)$ is not finished: we have
not yet explored the islands of supersingularity. We shall reach them in
\subsecref{sub-p-adic-betti-lattice}.

\newpage
\section{Abelian periods as algebraic integrals.}

\begin{abst}
 We discuss periods of abelian varieties. Their $p$-adic counterparts 
 live naturally in Fontaine's ring $\matheur{B}_{\mathrm{dR}}$. We present 
 Colmez' construction of abelian $p$-adic periods, which relies on 
 $p$-adic integration and reflects as closely as possible the complex
 picture.
 %In the relative case, the relation with the Picard-Fuchs
% differential equation is delicate, except in the special case of
% multiplicative reduction, where complex and $p$-adic periods can 
% be computed essentially as evaluations of formal solutions
% (\eg hypergeometric functions).
 We also deal with the concrete evaluation of elliptic $p$-adic periods.
\end{abst}

\subsection{Over $\C$.}\label{sub-complex-integ}

Let $A$ be an abelian variety over $\C$ of dimension $g$, 
and $\omega_1,\ldots,\omega_g$ a basis of invariant differential forms.
%Set $\Lambda=\mathrm{Im}(i\colon\mathrm{H}_1(A(\C),\Z)\rightarrow
%\C^g)$ with $i$ sending
%\begin{equation*}
%\gamma\mapsto\left(
%\int_{\gamma}\omega_i\right).
%\end{equation*}
Let $\Lambda$ be the image of the map
\begin{equation*}
 \iota:\mathrm{H}_1(A(\C),\Z)\rightarrow \C^g\; ;\;\; \gamma\mapsto
 \left(\int_{\gamma}\omega_i\right)_i.
\end{equation*}
One has a canonical isomorphism 
\begin{equation*}
\C^g/\Lambda\stackrel{\sim}{\longrightarrow}A(\C).
\end{equation*}
Let us denote the projection $\C^g\rightarrow A(\C)$ by $\mathrm{pr}$.
For any differential one-form $\omega$ of the second kind, one can 
consider a primitive function $f_{\omega}$
of $\mathrm{pr}^{\ast}\omega$; it is meromorphic (univalued) on 
$\C^g$, because all residues of $\omega$ are $0$ by definition, and
it is unique up to addition of an arbitrary constant.
Then $f_{\omega}(z_1+z_2+z_3)-f_{\omega}(z_1+z_2)-f_{\omega}(z_1+z_3)
+f_{\omega}(z_1)$ defines a periodic function on $(\C^g)^3$, hence
induces a meromorphic function on $A(\C)^3$, which we denote by
$F^3_{\omega}$.
Note that the function $F^3_{\omega}$ can also be defined purely
algebraically by the following conditions:
\begin{itemize}
\item $F^3_{\omega}(z_1,0,z_3)=F^3_{\omega}(z_1,z_2,0)=0$.
\item $dF^3_{\omega}=m^{\ast}_{123}\omega-m^{\ast}_{12}\omega
-m^{\ast}_{13}\omega+m^{\ast}_1\omega,$
\end{itemize}
where $m_{123}$ is the addition $A^3\rightarrow A$ sending 
$(z_1,z_2,z_3)$ to $z_1+z_2+z_3$, etc.
For $\gamma\in\mathrm{H}_1(A(\C),\Z)$ the value
$f_{\omega}\bigl(i(\gamma)+a\bigr)-f_{\omega}(a)$ 
does not depend on $a\in\C^g$ ($a$ being chosen so that
$i(\gamma)+a$ and $a$ are not poles of $f_{\omega}$),
and one has the equality
\begin{equation*}
f_{\omega}\bigl(i(\gamma)+a\bigr)-f_{\omega}(a)=
\int_{\gamma}\omega,
\end{equation*}
which describes the period pairing 
\begin{equation*}
\mathrm{H}^1_{\mathrm{dR}}(A)
\times\mathrm{H}_1(A(\C),\Z)\longrightarrow\C
\end{equation*}
(or equivalently, the isomorphism
$\mathrm{H}^1_{\mathrm{dR}}(A)\otimes\C
\stackrel{\sim}{\rightarrow}
\mathrm{H}^1_{\mathrm{B}}(A(\C),\Z)\otimes\C$.)
Here $\mathrm{H}^1_{\mathrm{dR}}(A)$ denotes the first \emph{algebraic}
de Rham cohomology group of $A$, which coincides with the group of 
differential forms of the second kind modulo exact forms.

\subsection{Over the $p$-adics; prolegomena}\label{sub-prolegomena}

If we try to translate this into the $p$-adic setting, one has to face
at once the problem: what is \emph{integration over a loop}?

\begin{para}\label{para-padic-int-1}
 One tentative way is via Berkovich's theory, where loops do exist.  Let us for
 instance consider the Legendre elliptic curve $X_z$. Viewed as a
 $p$-adic space, for $z \in \matheur{D}(0,1^-)\setminus \lbrace 0 \rbrace$
 ($p\neq 2$), this is a Tate curve: $X_z^{\mathrm{an}} \simeq
 {\C}_p^{\times}/q^{\Z}$. The canonical differential
 $\omega_{\mathrm{can}}$ inherited from ${\C}_p^{\times}$ and
 $\frac{dx}{y}$ are proportional:

 \begin{equation*}
 \omega_{\mathrm{can}}=\Theta(z).\frac{dx}{y} \;\text{ with }\;
 1/\Theta(z)=\sqrt{-1}F\Bigl(\frac{1}{2},\frac{1}{2},1;z\Bigr).
 \end{equation*}

 The basic element $\gamma$ of $\mathrm{H}_1(X_z^{\mathrm{an}},{\Z})$
 can be identified with the generator $q$ of $q^{\Z}$, and it is natural
 to set $\int_{\gamma}\omega_{\mathrm{can}}=\int_{0}^q \frac{dt}{t}=
 \log(q)$ (choosing a branch of the $p$-adic logarithm). One then has

 \begin{equation*}
 \omega_2^{(p)}:=\int_{\gamma}^{(p)}\frac{dx}{y} =
 \sqrt{-1}F\Bigl(\frac{1}{2},\frac{1}{2},1;z\Bigr)\log(q).
 \end{equation*}
 This formula, as well as
 \begin{align*}
 & \sqrt{q} =
 \frac{z}{16}
 e^{\frac{F^{\ast}(\frac{1}{2},\frac{1}{2},1;z)}
         {F(\frac{1}{2},\frac{1}{2},1;z)}}\\
 \biggl( & \text{with }F^{\ast}\Bigl(\frac{1}{2},\frac{1}{2},1;z\Bigr)
  = 4\sum_{n>0}\binom{2n}{n}^2 \biggl(\sum_1^{2n} \frac{(-1)^{m-1}}{m}\biggr) 
 \Bigl(\frac{z}{16}\Bigr)^n\biggr)
 \end{align*}
 are the same as those encountered in the complex situation. This leaves
 the open problem: how can one construct ``the other period''
 $\omega_1^{(p)}$, \ie the $p$-adic analogue of
 $\omega_1=2\pi\sqrt{-1}F(\frac{1}{2},\frac{1}{2},1;z)$, since the
 corresponding loop is missing?
\end{para}

\begin{para}\label{para-padic-int-2}
 One simple tentative answer would be to replace the topological
 covering ${\C}_p^{\times}\rightarrow X_z^{\mathrm{an}}$ by other \'etale
 coverings. Natural candidates for this purpose are \'etale coverings of
 order $p^n$, especially those corresponding to torsion points of order
 $p^n$ which are close to the origin.  Let $\zeta_{p^n}$ be a $p^n$th
 root of unity in ${\C}_p^{\times}$ and let $x_n$ be its image in
 $X_z^{\mathrm{an}}$. In the corresponding complex case, taking
 $\zeta_{p^n}=e^{2i\pi/{p^n}}$ would give the right answer
 $p^n\int_{0}^{x_n}\frac{dx}{y}=\omega_1$.  In the $p$-adic case, we get
 instead $p^n\int_{0}^{x_n}\frac{dx}{y}=0$.  Indeed, already in the case
 of the multiplicative group, we have
 $p^n\int_{1}^{\zeta_{p^n}}\frac{dt}{t}=p^n \log(1+(\zeta_{p^n}-1))=0\,$  
 $p$-adically: in other words, $2i\pi$ ``\emph{is missing}''.
\end{para}

\begin{para}[{\itshape Riemann-Shnirelman sums.}]
 It is appropriate to evoke here Shnirel\-man's approach to integration over loops,
 indeed one of the earliest works in $p$-adic analysis. This
 is an adaptation to the $p$-adic case of the computation of integrals by
 Riemann sums
 \begin{align*}
  \frac{1}{2i\pi}\int_{C(a,r)} f(z)dz & = \int_0^1 f(a+re^{2i\pi\theta})
   re^{2i\pi \theta}d\theta
   & = \,\lim_{m\to \infty}\frac{1}{m}\sum_{\zeta^m=1}f(a+r\zeta) r\zeta.
 \end{align*}
 Now, let $a, r \in \C_p$ and let $f$ be a $\C_p$-valued function
 on the circumference
 $C(a,|r|)$.  Let us denote the limit
 \begin{equation*}
  \lim_{m\to \infty}\frac{1}{m}\sum_{\zeta^m=1} f(a+r\zeta).r\zeta 
 \end{equation*}
 symbolically
 by $\int_{C(a,r)} f(z)dz$. For better convergence, Shnirelman actually
 restricted the
 limit to those integers $m$ prime to $p$. In the case of an analytic
 function $f$ on $C(a,|r|)$, this restriction is unnecessary (one can even
 take $m=p^k$) and one has the following analogue
 of Cauchy's theorem of residues (\cf \cite[app.]{paascorw})

 \begin{lem}
  Assume that $f$ is a meromorphic function on $\matheur{D}(a,|r|^+)$,
  and that its poles $z_1,\ldots, z_{\nu}$ all lie in
  $\matheur{D}(a,|r|^-)$. Then $\int_{C(a,r)} f(z)dz$ exists and equals
  $\sum \Res_{z_i} f$.
 \end{lem}

 In the (trivial) special case $f = 1/z$, we get
 $\int_{C(0,1)}dz/z =1$ (not $0$ as
 in \ref{para-padic-int-2}!), but $2i\pi$ is still missing.
\end{para}

\begin{para}
 Actually, there is no way to remedy this if one remains in $\C_p$.  A
 deeper reason for that, due to Tate, is that while
 $\Gal({\ovl{\Q}}_p/{\Q}_p)$ acts on the inverse system of $p^n$th
 roots of unity ($n \geq 0$) through the cyclotomic character $\chi$,
 there is no element $(2i\pi)_p$ in $\C_p$ such that $g((2i\pi)_p)=\chi
 (g) (2i\pi)_p$ for every $g \in \Gal({\ovl{\Q}}_p/{\Q}_p)$.

 \noindent It turns out that $(2i\pi)_p$ is in a sense the only
 ``missing piece'': there is a good theory of $p$-adic periods (due to
 J.M. Fontaine, W. Messing \cite{ppapec}, G. Faltings) which lives in
 some ${\Q}_p$-algebra 
 isomorphic to $\C_p[[(2i\pi)_p]]$, and which we shall now describe in the
 case of abelian varieties.
\end{para}

\subsection{The Fontaine ring $\matheur{B}_{\mathrm{dR}}$.}
\begin{para}
 {\itshape The ring $\matheur{R}$.\ }\index{000R@$\matheur{R}$}
 Let $\mathcal{O}_{\C_p}$ be the ring of
 integers of $\C_p$, \ie $\{x \in \C_p$, $|x|_p \leq 1\}$.  Let us set
 \begin{equation*}
  \matheur{R}:=\underset{x\mapsto x^p}{\limproj}
   \OO_{\C_p},
 \end{equation*}
 \ie $\matheur{R}$ is the set of all series
 $\left(x^{(n)}\right)_{n\in\N}$
 such that $\left(x^{(n+1)}\right)^p=x^{(n)}$.
 This is in fact a ring of characteristic $p$ with 
 \begin{equation*}
  \left(\bigl(x^{(n)}\bigr)+\bigl(y^{(n)}\bigr)\right)^{(n)}=
   \lim_{m\rightarrow\infty}\left(x^{(n+m)}+y^{(n+m)}\right)^{p^m}
 \end{equation*}
 and
 \begin{equation*}
  \left(\bigl(x^{(n)}\bigr)\cdot\bigl(y^{(n)}\bigr)\right)^{(n)}=
   x^{(n)}\cdot y^{(n)}.
 \end{equation*}
 Let $\mathrm{W}(\matheur{R})$ be the Witt ring with coefficients in
 $\matheur{R}$. for $x\in \matheur{R}$, let $[x]$ denote the Teichm\"uller
 representative in $\mathrm{W}(\matheur{R})$.
\end{para}

\begin{para}
 {\itshape The ring $\matheur{B}^+_{\mathrm{dR}}$.\ }
 \index{000bdr+@$\matheur{B}^+_{\mathrm{dR}}$}
 To any element 
 $(x_0,x_1,\ldots,x_n,\ldots)=\sum p^n[x^{p^{-n}}_n]$ in
 $\mathrm{W}(\matheur{R})$, we associate
 \begin{equation*}
 \theta((x_n))=\sum^{\infty}_{n=0}p^nx^{(n)}_n.
 \end{equation*}
 This defines a surjective homomorphism
 $\theta\colon\mathrm{W}(\matheur{R})\rightarrow \OO_{\C_p}$,
 whose kernel is principal.
 It extends to a homomorphism $\theta:
 \mathrm{W}(\matheur{R})\left[\frac{1}{p}\right]\rightarrow\C_p$,
 and $\matheur{B}^+_{\mathrm{dR}}$ is defined as the $(\Ker\theta)$-adic 
 completion
 \begin{equation*}
 \matheur{B}^+_{\mathrm{dR}} :=\lim_{\longleftarrow}
 \mathrm{W}(\matheur{R})\textstyle{\left[\frac{1}{p}\right]}
 \Big/(\Ker\theta)^n.
 \end{equation*}
 By continuity, $\theta$ further extends to a homomorphism $\theta\colon
 \matheur{B}^+_{\mathrm{dR}}\rightarrow\C_p$.
 Then $\matheur{B}^+_{\mathrm{dR}}$ is a complete discrete valuation
 ring with maximal ideal $\Ker\theta$ and residue field $\C_p$.
 Moreover, the Galois group $\Gal(\ovl{\Q_p}/\Q_p)$ acts on 
 $\matheur{B}^+_{\mathrm{dR}}$ in such a way that $\theta$ is equivariant 
 with respect to this Galois action,
 \begin{equation*}
  \Gr^{\bp}\matheur{B}^+_{\mathrm{dR}}\simeq \bigoplus_{r\in \N}\C_p(r),
 \end{equation*}
 where $\Gr^{\bp}$ refers to the filtration by the powers of
 $\Ker\theta$, and where the ``twist'' $(r)$ indicates that the 
 $\Gal(\ovl{\Q_p}/\Q_p)$-action is twisted by the $r$th power of the
 cyclotomic character.

 It turns out that $\matheur{B}^+_{\mathrm{dR}}$ contains naturally a copy of
 ${\ovl{\Q}}_p$ (which $\theta$ maps isomorphically to ${\ovl{\Q}}_p \subset
 \C_p$). More precisely, P. Colmez \cite{lnasddb} has shown that
 $\matheur{B}^+_{\mathrm{dR}}$ is the separated \emph{completion} of
 ${\ovl{\Q}}_p$ with respect to the topology defined by taking
 $\Bigl(p^n\mathcal{O}^{(k)}_{{\ovl{\Q}}_p}\Bigr)_{n,k}$ as a basis of
 neighborhoods of $0$, where $\mathcal{O}^{(k)}_{{\ovl{\Q}}_p}$ denotes
 the subring of ${\ovl{\Z}}_p$ defined inductively as
 \begin{equation*}
 \mathcal{O}^{(0)}_{{\ovl{\Q}}_p}
  ={\ovl{\Z}}_p,\;\mathcal{O}^{(k)}_{{\ovl{\Q}}_p}
  =\Ker\Bigl(d: \mathcal{O}^{(k-1)}_{{\ovl{\Q}}_p}
  \rightarrow
  \Omega^1_{\mathcal{O}^{(k-1)}_{{\ovl{\Q}}_p}/{\Z}_p}\otimes
  {\ovl{\Z}}_p\Bigr).
 \end{equation*}
 It is easy to deduce from this description that
 $\matheur{B}^+_{\mathrm{dR}}$ contains $\widehat{{\Z}_p^{\mathrm{ur}}}$.
\end{para}

\begin{para}
 {\itshape Some remarkable elements of $\matheur{B}^+_{\mathrm{dR}}$.}
 For any $x \in \matheur{B}^+_{\mathrm{dR}}$ such that $|\theta(x)-1|_p< 1$,
 the series $\log(x) = - \sum_{n>0} \frac{(1-x)^n}{n}$ converges in
 $\matheur{B}^+_{\mathrm{dR}}$.  In particular, let $\underline {z} =
 (\ldots,z^{(1)},z^{(0)})$ be an element of $\matheur{R}$ such that $z^{(0)}\in
 {\ovl{\Z}}_p$ or $\widehat{{\Z}_p^{\mathrm{ur}}}$.  Then one can define the
 element \index{000Log@Log}
 $\Log \udl{z} = \log
 \bigl(\frac{z^{(0)}}{[\udl{z}]}\bigr)$\footnote{we follow Colmez' sign 
 convention; Fontaine's $\LOG$ is $\log - \Log$.%\index{000LOG@LOG}
 }.
 Note that $\theta(\Log \underline{z})=0$. 

 In the special case where each
 $z^{(n)}=\zeta_{p^n}$ is a primitive $p^n$th root of unity,
 $\Log(\ldots, \zeta_{p^n}, \ldots,1)$ is the element $(2i\pi)_p$ 
 \index{0002ipp@$(2i\pi)_p$}
 we were looking for. This is a generator of $\Ker \theta$, and 
 $\Gal({\ovl{\Q}}_p/{\Q}_p)$ acts on it through the cyclotomic character;
 in other words, we have a $\Gal({\ovl{\Q}}_p/{\Q}_p)$-equivariant
 isomorphism

 \begin{equation*}
 \Gr^{\bp}\matheur{B}^+_{\mathrm{dR}}\;\simeq\; {\C}_p[(2i\pi)_p].
 \end{equation*}
 
 Note that another choice of $(\ldots, \zeta_{p^n},
 \ldots,1)$ changes $(2i\pi)_p$ by multiplication by a unit in ${\Z}_p$.
 Similarly,
 up to addition of an element of $(2i\pi)_p{\Z}_p$, $\Log\;\underline
 {z}\;$ depends only on
 $z^{(0)}$; it is sometimes simply denoted by $\Log\;z^{(0)}$.

 By choosing a double embedding of ${\ovl{\Q}}$ into $\C$ and
 ${\ovl{\Q}}_p$,
 one obtains a canonical element $(2i\pi)_p$, attached to the sequence
 $(\ldots, e^{2i\pi/{p^n}},\ldots,1)$. Similarly, if $z^{(0)}\in
 {\ovl{\Q}}$, then $\Log z^{(0)}$ is 
 well-defined up to addition of an element of $(2i\pi)_p{\Z}$.
\end{para}

\subsection{Colmez' construction of abelian $p$-adic periods.}

Let $A$ be an abelian variety defined over a $p$-adic local field
$K$. The Fontaine-Messing $p$-adic period pairing is a pairing

\begin{equation*}
\int^{(p)}\;:\;\mathrm{H}^1_{\mathrm{dR}}(A) \otimes T_p(A_{\ovl{K}})
 \longrightarrow \matheur{B}^+_{\mathrm{dR}}
\end{equation*}
where $T_p(A_{\ovl{K}}) = \varprojlim \Ker([p^n]: A_{\ovl{K}}\rightarrow
A_{\ovl{K}})$ is the Tate module (a ${\Z}_p$-module of rank $2\dim A$
with $\Gal({\ovl{K}}/K)$-action).

We present Colmez' construction, which is parallel to \ref{sub-complex-integ}.

Let $\omega$ be a differential form on $A$ of second kind, and
$F_{\omega}^3$ the function on $A^3$ determined as in
\ref{sub-complex-integ}.  Then:
\begin{pro}[{\cite[4.1]{ppdva}}]
There exists a locally meromorphic function $F_{\omega}$ on 
$A(\matheur{B}_{\mathrm{dR}})$, unique up to constant, such that:

\begin{enumerate}
 \renewcommand{\theenumi}{\arabic{enumi}}
 \item $dF_{\omega}=\omega$.
 \item $F_{\omega}(z_1+z_2+z_3)-F_{\omega}(z_1+z_2)-
       F_{\omega}(z_1+z_3)+F_{\omega}(z_1)=F^3_{\omega}(z_1,z_2,z_3)$.
 \item If $\omega=dF$, then $F_{\omega}=F$.
\end{enumerate}

\noindent
Moreover, if $\alpha\colon A_1\rightarrow A_2$ is a morphism of abelian 
varieties and $\omega$ is a differential form of second kind on
$A_2$, then $F_{\alpha^{\ast}\omega}=\alpha^{\ast}F_{\omega}$.
\end{pro}

Note that the function $F_{\omega}$ is \emph{not} multivalued; this
fact comes from the following lemma specific to the $p$-adic case:

\begin{lem}[{\cite[4.3]{ufdpplpdvaamc}}]
For any neighborhood $V$ of $0$ in $A(\matheur{B}_{\mathrm{dR}}^+)$, there
exists an open subgroup $U$ of $A(\matheur{B}_{\mathrm{dR}}^+)$ contained in 
$V$ such that $A(\matheur{B}_{\mathrm{dR}}^+)/U$ is a torsion group.
\end{lem}

Take a proper model $\mathscr{A}$ of $A$ over $\OO_{K}$.
Let $\gamma=(\cdots,u_n,\cdots,u_2,u_1=0)\in\mathrm{T}_p(A_{\ovl{K}})$ 
with each $u_n\in\mathscr{A}(\OO_{\C_p})$, and choose
$a_n\in\mathscr{A}(\matheur{B}^+_{\mathrm{dR}})$ so that neither
$a_n$ nor $a_n+_A\widehat{u}_n$ is close to a pole of $\omega$.
For suitable liftings 
$\widehat{u}_n\in\mathscr{A}(\matheur{B}^+_{\mathrm{dR}})$ of 
$u_n$ (\ie $\theta(\widehat{u}_n)=u_n$), the following holds. 

\begin{thm}[{\cite[5.2]{ufdpplpdvaamc}}]
The limit 
\begin{equation*}
 \oint^{(p)}_{\gamma}\omega \;:=\;
 \lim_{n\rightarrow\infty}p^n\bigl(F_{\omega}(a_n)-F_{\omega}
 (a_n+_A\widehat{u}_n)\bigr)
\end{equation*}
converges to an element in $\matheur{B}^+_{\mathrm{dR}}$, and it defines
a non degenerate bilinear pairing
\begin{equation*}
 \oint^{(p)}\;:\;
 \mathrm{H}^{1}_{\mathrm{dR}}(A)\otimes
 \mathrm{T}_p(A_{\ovl{K}})\longrightarrow\matheur{B}_{\mathrm{dR}}^+,
\end{equation*}
compatible with the Galois action and the filtrations.
\end{thm}

\begin{rem}
 On composing the period pairing $\int^{(p)}$ with $\theta$, one gets a
 bilinear map $\mathrm{H}^1_{\mathrm{dR}}(A) \otimes T_p(A_{\ovl{K}})
 \longrightarrow \C_p$ which sends $\Omega^1(A)\subset
 \mathrm{H}^1_{\mathrm{dR}}(A)$ to $0$.  This degenerate pairing
 describes ``half'' of the (Hodge)-Tate-Raynaud decomposition
 \begin{equation*}
  \mathrm{H}^1_{\mathrm{et}}(A_{\ovl{K}},{\Q}_p)\otimes {\C}_p \simeq
   \Omega^1(A)\otimes
   {\C}_p(-1) \oplus \mathrm{H}^1(\mathcal{O}(A))\otimes {\C}_p
 \end{equation*}
 of $\Gal({\ovl{K}}/K)$-modules.
\end{rem}


\subsection{Some computations of elliptic $p$-adic periods.}
\label{sub-computation}

\begin{para}{\itshape Tate elliptic curves.\ }
In this case, $T_p(A_{\ovl{K}})$ sits in an exact sequence 
\begin{equation*}
 0\rightarrow {\Z}_p(1)\rightarrow T_p(A_{\ovl{K}}) \rightarrow
  {\Z}_p \rightarrow 0.
\end{equation*}
Let us take for $\gamma_1$ the image of $(2i\pi)_p \in {\Z}_p(1)$ in
$T_p(A_{\ovl{K}})$, and for $\gamma_2$ any lifting of $1 \in {\Z}_p$.
For concreteness, let $A=X_z$ be in Legendre form as in
\ref{para-padic-int-1}.  Let $\eta \in \mathrm{H}^1_{\mathrm{dR}}$
satisfy $\an{\omega,\eta}=1$ (\eg $\frac{xdx}{4y}$).  We set
$\omega^{(p)}_i=\int^{(p)}_{\gamma_i}\omega,
\eta^{(p)}_i=\int^{(p)}_{\gamma_i}\eta$.  Then one has the ``Legendre
relation''
\begin{equation*}
\omega^{(p)}_1\eta^{(p)}_2-\eta^{(p)}_1\omega^{(p)}_2=(2i\pi)_p
\end{equation*}
and the formulas (\cf \cite{pbl}, \cite{get})
\begin{equation*}
\frac{\omega^{(p)}_1}{(2i\pi)_p} = \sqrt{-1}F\Bigl(\frac{1}{2},
\frac{1}{2},1;z\Bigr) = 1/\Theta(z), \;\; \omega^{(p)}_2 =
\sqrt{-1}F\Bigl(\frac{1}{2},\frac{1}{2},1;z\Bigr) \Log q
\end{equation*}
(compare with the complex case in \ref{para-tate-para}, and also
\cite[p.406]{f2}).  Note that the ``period'' we found in
\ref{para-padic-int-1} by integration along the Berkovich loop is
 \emph{not} a Fontaine-Messing period (there is a $\log$ instead of
 $\Log$).
\end{para}

\begin{para}{\itshape Elliptic curves with ordinary reduction.\ }
We now assume that $z\in {\Z}_p^{\mathrm{ur}}$ and that $A=X_z$ has good
ordinary reduction (with the notation of \subsecref{sub-dwork-hyper},
this means that $z(1-z)h_p(z)\neq 0$).  Again, $T_p(A_{{\ovl{\Q}}_p})$
sits in an exact sequence (\cite[A.2.4]{alraec})
\begin{equation*}
0\rightarrow {\Z}_p \gamma_1\rightarrow
T_p(A_{{\ovl{\Q}}_p}) \rightarrow T_p(A\;\mathrm{mod}\, p) \rightarrow 0.
\end{equation*}
 Then one still has the formula $\Theta(z)=
\frac{(2i\pi)_p}{\omega^{(p)}_1}$, where $\Theta(z)\in
\widehat{{\Z}_p^{\mathrm{ur}}}$ is the Tate constant discussed in
\ref{para-tate-para} (compare \cite{pbl}, \cite[4.3]{itoecwcm}).  In
particular, $\frac{\omega^{(p)}_1}{(2i\pi)_p}$ is given by the
evaluation at $z$ of ``the'' extension of
$F(\frac{1}{2},\frac{1}{2},1;?)$ discussed in \subsecref{sub-serre-tate-para}.

\noindent On the other hand, let $\gamma_2 \in T_p(A_{{\ovl{\Q}}_p})$ be
such that the Weil pairing $\an{\gamma_1,\gamma_2}\in {\Z}_p(1)={\Z}_p
(2i\pi)_p$ is the chosen generator $(2i\pi)_p$. Then one has the
``Legendre relation'', and $\omega^{(p)}_2 = \omega^{(p)}_1\cdot\frac{{\Log
\;q}}{{(2i\pi)_p}}$, where $q\in \widehat{{\Z}_p^{\mathrm{ur}}}$ denotes
here the so-called {\it Dwork(-Serre-Tate) parameter} of $A$ introduced
in \ref{para-dwork-para}.  If $q=1$, \ie when $A$ is the \emph{canonical
lifting} of $A\pmod{p}$, then one can choose $\gamma_2\;$ such the
corresponding value of ${\Log 1}$ is $0$, and then $\omega^{(p)}_2 =0.$
\end{para}

\subsection{Periods of CM elliptic curves and values of the Gamma function.}

\begin{para}
 The case of an elliptic curve $A$ with \emph{supersingular reduction} is
 more delicate.  If $A$ does not have complex multiplication (CM), one
 can show, using \cite[A.2.2]{alraec}, that the four basic periods
 $\omega^{(p)}_1,\omega^{(p)}_2,\eta^{(p)}_1,\eta^{(p)}_2$ are
 algebraically independent over ${\ovl{\Q}}_p$, so that one cannot
 expect ``formulas'' for the periods. 

  On the other hand, it is well-known that the complex
 periods of an elliptic curve with complex multiplication can be
 expressed in terms of special values of the $\Gamma$ function. So one
 may ask for a $p$-adic analogue involving $\Gamma_p$
 (\cf~\ref{subsection-deriv-of-gamma}).  Such a formula has been found by A.
 Ogus, actually not for the periods, but for the action of Frobenius.
\end{para}

\begin{para} {\itshape The Lerch-Chowla-Selberg-Ogus formulas.}
 \label{para-lerch-chowla}
 Let $A$ be an elliptic curve with complex multiplication by
 ${\Q}(\sqrt{-d})$ over ${\ovl{\Q}}\,$ ($-d$ denotes a fundamental
 discriminant).  Let $\epsilon = (\frac{-d}{}) $ be the quadratic
 character $({\Z}/d)^\times\to {\Z}/2$ induced by the embedding
${\Q}(\sqrt{-d})\subset {\Q}(\zeta_d) $, $h=$ the class number of
 ${\Q}(\sqrt{-d})$, and $w=$ the 
 number of roots of unity in ${\Q}(\sqrt{-d})$.  For any $u\in
 ({\Z}/d)^{\times}$, we denote by $\an{\frac{u}{d}}$ the unique rational
 number in $]0,1]$ such that $d\an{\frac{u}{d}}\equiv u$.

 Let $v$ be a place of ${\ovl{\Q}}$ of residue characteristic $p$, with
 associated embedding ${\ovl{\Q}}\subset {\ovl{\Q}}_p$.  To any lifting
 $\psi_v$ to $\Gal({\ovl{\Q}}_p/{\Q}_p)$ of the Frobenius element in
 $\Gal({\Q}_p^{\mathrm{ur}}/{\Q}_p)$, one attaches a $\psi_v$-linear
 endomorphism $\Psi_v$ of
 $\mathrm{H}^1_{\mathrm{dR}}(A/{\ovl{\Q}})\otimes {\ovl{\Q}}_p$
 \footnote{using its crystalline interpretation, \cf \cite[4]{fadrc1};
 we refer to \cite{chambert98:_cohom} for a nice recent survey of crystalline
 cohomology.}.
 On the other hand, it is well-known that $A$ has supersingular
 reduction at $v$ if and only if $\epsilon(p)=-1$ or $0$.
\end{para} 

\begin{thm}\label{thm-ogus-formula}
 There exists a basis $(\omega, \eta)$ of
 $\mathrm{H}^1_{\mathrm{dR}}(A/{\ovl{\Q}})$ of eigenvectors under the
 action of $\sqrt{-d}$ ($\omega$ being the class of a regular
 differential), and an element $\gamma \in \mathrm{H}_1(A(\C),{\Q})$,
 such that
 \begin{align*}
  \int_{\gamma} \omega
  & = \sqrt{2i\pi}\prod_{u\in ({\Z}/d)^{\times}}
  \Bigl(\Gamma\Lan{\frac{u}{d}}\Bigr)^{\epsilon(u)w/4h} \\
  \int_{\gamma} \eta 
  & = \sqrt{2i\pi}\prod_{u\in({\Z}/d)^{\times}}
  \Bigl(\Gamma\Lan{-\frac{u}{d}}\Bigr)^{\epsilon(u)w/4h}
 \end{align*}
 and such that for every place of ${\ovl{\Q}}$
 of residue characteristic $p$ satisfying $\epsilon (p)=-1$,
 \begin{align*}
  \Psi_v^{\ast} (\omega)
  & = p\prod_{u\in ({\Z}/d)^{\times}}
  \Bigl(\Gamma_p\Lan{p\frac{u}{d}}\Bigr)^{-\epsilon(u)w/4h}\eta \\
  \Psi_v^{\ast} (\eta)
  & = \prod_{u\in ({\Z}/d)^{\times}}
  \Bigl(\Gamma_p\Lan{-p\frac{u}{d}}\Bigr)^{-\epsilon(u)w/4h}\omega
 \end{align*}
 up to multiplication by some root of unity.
\end{thm}

\noindent This is a concatenation of \cite[3.15, 3.9]{ogus90:_chowl_selber},
taking into account the formula $2hd/w=- \sum_1^d \epsilon(u) u$,
\cf \cite[VII]{vudtdaz}).

\medskip\noindent {\it Remark.} The last two formulas extend to the case when
$\epsilon (p)=+1$ if one interchanges $\omega$ and $\eta$ in the right
hand sides (for $d=4$ and $p\equiv 1\pmod{4}$, this
is compatible with Young's formula in \subsecref{sub-analytic-cont}, since
$\Gamma_p(1/2)^4=1$, and with the formula
$F(\frac{1}{2},\frac{1}{2},1;-1)=\frac{1}{2}
\frac{\Gamma(1/4)}{\Gamma(1/2)\Gamma(3/4)}$.). 

 In that case, the expressions 
 \begin{equation*}
  \prod\Bigl(\Gamma_p\Lan{p\frac{u}{d}}\Bigr)^{-\epsilon(u)w/4h},\;
   \prod\Bigl(\Gamma_p\Lan{p\frac{u}{d}}\Bigr)^{-\epsilon(u)w/4h}
 \end{equation*}
 are algebraic numbers: indeed, let $r$ be the order of the subgroup
 $\an{p}$ of $({\Z}/d)^{\times}$ 
 generated by $p$, and $k= p^r-1/d$; then
\begin{equation*}
 \prod_{u\in
  ({\Z}/d)^{\times}}\Bigl(\Gamma_p\Lan{p\frac{u}{d}}\Bigr)^{-\epsilon(u)}=
  \prod_{w\in ({\Z}/d)^{\times}/\an{p}}\prod_1^{r}
  \Gamma_p\biggl(\frac{p^i wk}{p^r-1}\biggr)^{-\epsilon(w)}
\end{equation*}
and one concludes by the Gross-Koblitz formula.  
 
 \begin{para} {\itshape The ramified case.}\label{para-ramif} In the
  ramified case, \ie when
  $\epsilon (p)=0$, there is again an analogue of the
  last two formulas for
  $p\neq 2$. This relies on
  Coleman's computation of the Frobenius matrix of Fermat curves of degree
  divisible by $p$ --- which have arboreal reduction modulo $p$. The result
  takes the same form as in the case $\epsilon (p)=-1$: $\Psi$ is now
  attached to an element $\psi$ of degree one in the Weil group (\ie
  a lifting of the Frobenius element in
  $\Gal({\ovl{\Q}}_p/{\Q}_p)$), the expressions $\,{p\frac{u}{d}}\,$ have
  to be replaced by $\,{\psi\frac{u}{d}}$, and $\Gamma_p$ must be extended to 
  ${\Q}_p$ 
  \cite[6.5]{coleman90:_froben_fermat}. 
  
  Coleman's complicated expressions have been simplified by F.~Urfels in
  his thesis (Strasbourg, 1998; unpublished). If one passes to the
  Iwasawa logarithm $\log_p$ \index{Iwasawa logarithm@Iwasawa logarithm}
  \index{000logp@$\log_p$}
  (at the cost of rational powers of $p$), the result is that one should
  replace $\log_p \Gamma_p \lan{\psi\frac{u}{d}}$ in Ogus' formula by
  $G_p\bigl(\lan{\psi\frac{u}{d}}\bigr)-G_p\bigl(\lan{\frac{u}{d}}\bigr)$, 
  where $G_p$ denotes Diamond's LogGamma function
  \index{Diamond's LogGamma function@Diamond's LogGamma function}
  \index{000Gp@$G_p$}
  \begin{equation*}
   G_p(x)= \lim_{m\to \infty}\frac{1}{p^m}\sum_{n=0,\ldots, p^m-1}
    (x+n)\log_p(x+n) -(x+n) .
  \end{equation*}
 
\noindent {\it Example.} We take $p=3$. Let $\Q[\sqrt{-3n}]$ be an
 imaginary quadratic field with fundamental discriminant $-3n$, $n$
 prime to $3$. Let  $A$  be an elliptic curve with 
complex multiplication by some order in $\Q[\sqrt{-3n}]$.

Let $\psi$ be an element of the Weil group which acts by $-1$ on
$\Q_3/\Z_3\cong \mu_{3^\infty}$. One has a formula $\Psi_v^{\ast}
(\omega) = \kappa\eta$ with

\begin{equation*}
 \log_3 \kappa =-(\epsilon(u)w/4h)\sum_{u\in ({\Z}/3n)^{\times}}
 \Bigl(G_3\Bigl(\Lan{\psi\frac{u}{3n}}\Bigr) - 
 G_3\Bigl(\Lan{\frac{u}{3n}}\Bigr)\Bigr)
\end{equation*}
 We notice that for $u\in
 ({\Z}/3n)^{\times}$,
 $\epsilon\Bigl(3n\Lan{\psi\frac{u}{3n}}\Bigr)=-\epsilon (u)$, 
 whence
\begin{equation*}
  \log_3 \kappa = (w/2h) \sum_{u\in ({\Z}/3n)^{\times}}
 \epsilon(u)G_3\Bigl(\Lan{\frac{u}{3n}}\Bigr).
\end{equation*}

 On the other hand, for $p=3$,
the Teichm\"uller character is nothing but the Legendre symbol
$(\frac{-3}{{\,}})$. Using this, the latter expression can be rewritten,
by the Ferrero-Greenberg formula (\cf \cite[chap. 17]{cf1a2}), as
\begin{align*}
 \log_3\, \kappa = (w/2h) . (L'_3(0,\epsilon) +
 L_3(0,\epsilon)\log_3n)\\ =  (w/2h ) .\sum_{v\in ({\Z}/n)^{\times}} 
 \Bigl(\frac{n}{{v}}\Bigr)\log_3\Gamma_3 \Lan{\frac{v}{n}}.  
\end{align*} 
 When $(\frac{n}{{3}})=1$, then
$L_3(0,\epsilon)= 0$ and the Gross-Koblitz formula shows that
$L'_3(0,\epsilon)= \sum 
  \bigl(\frac{n}{{v}}\bigr)\log_3\Gamma_3 \lan{ \frac{v}{n}} $ is the
 Iwasawa logarithm of an algebraic number (a Gauss sum, {\it \cf
 loc.~cit.}), hence $\kappa\in \ovl{\Q}$.

 \noindent  A particularly simple case is $n=1$: one finds $\log_3
 \kappa=0$, thus $\kappa
 \in 3^\Q.\mu_\infty \subset \ovl{\Q}$.

 \medskip
 \noindent For $n=8$, a contrario, one has $(\frac{n}{{3}})=-11$: one finds  
 \begin{equation*} \log_3 \kappa = 
  (1/2).\Bigl[\log_3\Gamma_3 \Lan{\frac{1}{8}}-\log_3\Gamma_3
   \Lan{\frac{3}{8}}-\log_3\Gamma_3
   \Lan{\frac{5}{8}}+\log_3\Gamma_3 \Lan{\frac{7}{8}}\Bigr]=0
 \end{equation*}
 by the functional equation of
 $\Gamma_3$, so that again $\kappa
 \in 3^\Q.\mu_\infty \subset \ovl{\Q}$.
\end{para}


\begin{para} {\itshape Colmez' product formula.} There is a natural
 extension of $|\quad|_p$ on ${\ovl{\Q}}_p$ to
 $\matheur{B}^+_{\mathrm{dR}}$ (however not as an absolute value), such
 that $\vert (2i\pi)_p \vert_p = p^{-\frac{1}{p-1}}$
 \cf \cite{ufdpplpdvaamc}. 
 
 Colmez has remarked that the logarithm of the
 product $|2i\pi| \prod_p |(2i\pi)_p|_p $ $=$ $ 2\pi\prod_p
 p^{-\frac{1}{p-1}}$ 
 is formally equal to $\log (2\pi) + \zeta^{\prime}(1)/\zeta(1)$, a
 divergent sum which can be renormalized using the functional equation of
 $\zeta$: setting $\zeta^{\prime}(1)/\zeta(1)$ $=$ 
 $-\zeta^{\prime}(0)/\zeta(0)=-\log(2\pi)$, the renormalized product is
 \begin{equation*}
  |2i\pi| \sideset{}{'}\prod_p |(2i\pi)_p|_p = 1.
 \end{equation*}
 He has given an amazing generalization of this product formula to
 periods of CM elliptic curves and many other CM abelian varieties
 (loc.~cit.).
 \end{para}

\newpage
\section{Periods as solutions of the Gauss-Manin connection.}

\begin{abst}
 The variation of the periods in a family of complex algebraic varieties
 is controlled by the Picard-Fuchs differential equation. We discuss the
 p-adic situation, present Dwork's general viewpoint on p-adic
 periods. We then discuss the question of the existence of an arithmetic
 structure on the space of solutions, which would give an intrinsic
 meaning to the period modulo $\ovl{\Q}^\times$. We present such an
 arithmetic structure (analogous to the Betti lattices in the complex
 situation) in the case of abelian varieties with either multiplicative
 reduction or supersingular reduction. In the latter case, we relate the
 periods, in the presence of complex multiplication, to special values
 of the $p$-adic Gamma function.
 \end{abst}

\bigskip
\subsection{Stokes.}

We come back to the basic question of the meaning of \emph{integration
over a loop}, already discussed in \subsecref{sub-prolegomena}.  A
different approach is based on the Stokes lemma: integrating \emph{exact}
differentials on a loop gives zero. This is the approach favored by
K. Aomoto in the complex case and by B. Dwork in the $p$-adic case.

\begin{para}
 In order to see how this idea can be implemented, let us consider the
 Hankel expression for the gamma function
 \begin{equation*}
  \frac{1}{\Gamma(1-\alpha)} = \frac{1}{2i\pi}
   \int_{\gamma} x^{\alpha}e^x\frac{dx}{x}
 \end{equation*}
 where $\gamma$ is the following loop based at $-\infty$.

 \begin{figure}[h]%PICTURE9
  \begin{picture}(200,40)(0,0)
   \put(0,0){\includegraphics{fig9.eps}}
   \put(108,12.5){$\bbp$}
   \put(117,10){$\scriptstyle{0}$}
   \put(135,15){$\gamma$}
  \end{picture}
  \caption{}
 \end{figure}

 \noindent Stokes' lemma suggests to attach to this integral the
 following complex
 \begin{equation*}
 x^{\alpha}e^x\mathcal{O} \xrightarrow{xd/dx}  x^{\alpha}e^x\mathcal{O}
 \end{equation*}
 where $\mathcal{O}$ is a suitable ring of analytic functions.  From the
 point of view of index theory (Malgrange-Ramis), a natural choice is
 \begin{equation*}
 \mathcal{O}= {\C}[[x]]_{-1,1^-}=\biggl\{
 \sum_{n\geq 0} a_n x^n \biggm| \exists \kappa >0, \exists r \in ]0,1[,\;
 |a_n| \leq \kappa r^n/n!\biggr\}
 \end{equation*}
 identified with the ring of entire functions of {\it exponential order}
 $O(e^{r|x|})$ for some $r<1$\footnote{we use the traditional
 notation for rings of Gevrey series; the value $(-1,1^-)$ of the index
 is ``characteristic'', \ie an extremal value for which $\mathrm{H}^1$
 is of dimension one.}.  Using the formula $x^{\alpha}e^x x^k=
 \frac{-1}{\alpha + k}x^{\alpha}e^x x^{k+1} + xd/dx(\frac{1}{\alpha +
 k}x^{\alpha}e^x x^k)$, a simple computation then shows that
 $\mathrm{H}^0=0$ and that $\mathrm{H}^1$ is of dimension $1$; it is
 generated by $[\frac{1}{2i\pi}x^{\alpha}e^x]$, whose integration along
 $\gamma$ gives $\frac{1}{\Gamma(1-\alpha)}$.

 \par \noindent The $p$-adic analogue
 \begin{equation*}
 \C_p[[x]]_{-1,1^-}=\biggl\{
 \sum_{n\geq 0} a_n x^n \biggm|\exists \kappa >0, \exists r \in ]0,1[,\;
 |a_n|\leq \kappa r^n/| n!|_p \biggr\}
 \end{equation*}
 is nothing but the ring of \emph{overconvergent} analytic functions on
 $\matheur{D}(0,|\pi|_p)$ (here $\pi$ is Dwork's constant).  By change
 of variable $z=\pi x$, we thus recover the complex
 $M_{\alpha}^{\dagger} \xrightarrow{zd/dz} M_{\alpha}^{\dagger}$ of
 \ref{subsection-deriv-of-gamma}, where $\Gamma_p(\alpha)$ $\bigl(=\pm
 \frac{1}{{\Gamma_p(1-\alpha)}}\bigr)$ appeared.
\end{para}

\begin{para}
 This approach inspired by Stokes' lemma is well-suited for studying
 many kinds of hypergeometric functions (confluent or not). Basically,
 the integral will satisfy difference equations with respect to the
 exponents ($\alpha$ in the previous example), and differential
 algebraically on parameters. In the complex situation, one of the main
 problems is the construction of nice loops (such as $\gamma$); this is
 the object of a so-called ``topological intersection theory''
 generalizing the usual Betti homology and period pairing. In the
 $p$-adic case, the focus has been more on the construction and
 properties of the Frobenius structure, in particular its analyticity
 with respect to the exponents (``Boyarsky principle'' \cite{otbp}).
\end{para}

\begin{para}
 When no exponential is involved and when the exponents are rational,
 the integrand is algebraic; if it depends algebraically on parameters,
  the period integral is a solution of a {\it Gauss-Manin connection}. 

 Let us say a few words about the classical case of
 $F(a,b,c;z)$.  We assume that $a,b,c \in {\Z}_p$, and that $c-a,c-b,b,a$
 all lie outside $\Z$.  We set
 \begin{equation*}
 f_{a,b,c;z}=
 x^b(1-x)^{c-b}(1-zx)^{-a}
 \end{equation*}
 where, for simplicity, $z$ is limited to $\vert z(1-z)\vert = 1$.  Let
 $Z$ be the affinoid $\matheur{D}(0,1^+)\setminus
 (\matheur{D}(0,1^-)\cup \matheur{D}(1,1^-)\cup \matheur{D}(1/z,1^-))$.
 The relevant complex is
 \begin{equation*}
 f_{a,b,c;z}\mathcal{H}^{\dagger}(Z)
 \xrightarrow{xd/dx} f_{a,b,c;z}.\mathcal{H}^{\dagger}(Z).
 \end{equation*}
 Its $\mathrm{H}^1=\mathrm{H}^1_{a,b,c}$ is of dimension $2$; it is
 generated by $[f_{a,b,c;z}]$ and $[\frac{f_{a,b,c;z}}{1-x}]$. The
 Gauss-Manin connection is
 \begin{equation*}
 \nabla\Bigl(\frac{d}{dz}\Bigr)
  \begin{pmatrix}
   [f_{a,b,c;z}] \\ {}
   [\frac{f_{a,b,c;z}}{1-x}]
  \end{pmatrix}
 =
 \begin{pmatrix}
  -\frac{c}{z} & \frac{c-b}{z} \\ 
  \frac{c-a}{1-z} & \frac{a+b-c}{1-z}
 \end{pmatrix}
 \begin{pmatrix}
  [f_{a,b,c;z}] \\ {}
  [\frac{f_{a,b,c;z}}{1-x}]
 \end{pmatrix}
 \end{equation*}
  \cf \cite[3.1, 1.2]{lopde}.  The Frobenius structure relates
 $\mathrm{H}^1_{a,b,c}$ to
 $\mathrm{H}^1_{a^{\prime},b^{\prime},c^{\prime}}$, where
 $a^{\prime},b^{\prime},c^{\prime}$ are the successors of $a,b,c$
 respectively (\cf~\ref{subsection-deriv-of-gamma}).  For all these
 hypergeometric series, there is a story similar to that of
 $F(\frac{1}{2},\frac{1}{2},1;z)$.  In this context, it is fruitful to
 combine the Boyarsky principle and the contiguity relations.
\end{para}

\begin{para}
 In this sketch of Dwork's approach, ``periods''
 are just analytic solutions of the Gauss-Manin connection. One can ask
 more: namely, one can ask for a dual theory of $p$-adic cycles and a
 period pairing as in the complex case\footnote{Dwork has developed a
 ``dual theory'' and a pairing given by residues, which play somehow the
 role of ``$p$-adic cycles'' (\cf \cite[5]{uinacdd}); however, this dual
 theory is an avatar of de Rham cohomology with proper supports and
 does not enjoy the properties of discreteness and horizontality that we
 are looking for.}.  At best, we can expect these ``{\it $p$-adic Betti
 lattices}'' to be locally horizontal, functorial with respect to the
 endomorphisms of the geometric fibres, and defined over $\Z$ or at least
 over some number field.

 It turns out that this can be done in some
 cases, \eg in the cases of abelian varieties with multiplicative
 reduction, and of abelian
 varieties with supersingular reduction (\cite{pbl}, \cite{tdmeigdvpdg}). Let
us outline the results.
\end{para}
  

\subsection{$p$-adic Betti lattices for abelian varieties with 
multiplicative reduction.}\label{sub-p-adic-betti-lattice}

\begin{para}\label{para-multi-uniformize}
 {\itshape Over $\C$.}
 We first recall the ``multiplicative uniformization'' of complex
 abelian varieties.  Let $A$ be an abelian variety of dimension $g$ over
 $\C$.  One has the analytic representation
 \begin{equation*}
  A(\C)=T(\C)/M,
 \end{equation*}
 where $T$ is a torus of dimension $g$, and $M$ is a lattice of rank
 $g$.  Set $M' :=\mathrm{Hom}(T,\G_m)\,(=M_{A^{\vee}})$.  Then we
 have the exact sequence
 \begin{equation*}
  0\longrightarrow 2i\pi {M'}^{\vee}\longrightarrow\Lambda\longrightarrow
   M\longrightarrow 0,
 \end{equation*}
 where $\Lambda$ is the period lattice (of rank $2g$), and this sequence
 splits by choosing a branch of $\log$.  The inclusion $M\hookrightarrow
 T$ defines a pairing $M\times M' \rightarrow\C^{\times}$.  A
 polarization, on the other hand, induces $M\rightarrow M'$.  These
 data give rise to the pairing ({\it ``multiplicative period''})
\begin{equation*}
 q=(q_{ij})\colon M\otimes M\longrightarrow\C^{\times}
\end{equation*}
and $-\log|q|$ is a scalar product on $M_{\R}$.
\end{para}

\begin{para}\label{para-degenerate}
 {\itshape Over the $p$-adics.}
 Let $K$ be a $p$-adic local field
 and $A$ an abelian variety over $K$ having (split) multiplicative
 reduction.  One has a similar analytic representation \cite{aacodavocr}
 \begin{equation*}
  A(K)=T(K)/M
 \end{equation*}
 Set again $M' :=\Hom(T,\G_m)$.  We have the following identifications
 \begin{equation*}
  M\simeq \mathrm{H}_1(A^{\mathrm{rig}},\Z)\quad\mathrm{and}\quad
   M'\simeq \mathrm{H}_1(A^{\vee\,\mathrm{rig}},\Z),
 \end{equation*}
 where $A^{\mathrm{rig}}$ denotes the associated rigid analytic variety
 over $K$ to $A$.  Then, just as in the complex case, one can 
 construct a natural exact sequence:
 \begin{equation*}
  0\longrightarrow (2i\pi)_p{M'}^{\vee}\longrightarrow\Lambda
   \longrightarrow M\longrightarrow 0,
 \end{equation*}
 split by the choice of a branch of the $p$-adic logarithm $\log_p$, 
 and a non-degenerate pairing
 \begin{equation*}
 \int^p\colon\mathrm{H}^1_{\mathrm{dR}}(A)\times\Lambda\longrightarrow
 K[(2i\pi)_p].
 \end{equation*}
 (The construction involves one-motives and Frobenius \cite{pbl}).
 Restricted to $(2i\pi)_p{M'}^{\vee}$ which is in a 
 canonical way a subgroup of $T_p(A_{\ovl{K}})$, this pairing coincides
 with the Fontaine-Messing pairing.
\end{para}

\begin{para}
 {\itshape Degenerating abelian pencils.}  Let $\mathcal{A}\rightarrow
 {\A}^1\setminus \{0,\zeta_1,\ldots,\zeta_r\}$ be a pencil of polarized
 abelian varieties defined over a number field $k$. We consider an
 element $\omega$ of the relative $\mathrm{H}^1_{\mathrm{dR}}$ and its
 associated Picard-Fuchs differential equation $\mathcal{L}\omega=0$. We
 assume that the connected component of identity of the special fibre at
 $0$ of the Neron model of $\mathcal{A}$ is a split torus $T$ over $k$.
 The dual abelian pencil then has the same property, with a torus
 $T^{\prime}$.  We set $M=\Hom(T',{\G}_m)$, $M'=\Hom(T,{\G}_m)$.

 Let us fix $k\hookrightarrow \C$.  The constant subsheaf of
 $R_1(f^{\mathrm{an}}_{\C})_{\ast}{\Z}$ in the neighborhood of $0$
 identifies with $2i\pi M^{\prime \vee}$: its fibre at any $z (\neq
 0,\zeta_i)$ is the lattice $2i\pi M_z^{\prime \vee}$ attached to
 $\mathcal{A}_z$ (\ref{para-multi-uniformize}).  Similarly for $M$.  On
 the other hand, the choice of a branch of $\log$ identifies
 $\mathrm{H}_{1\mathrm{B}}(\mathcal{A}_z, {\Z})\simeq
 (R_1(f^{\mathrm{an}}_\C)_{\ast}{\Z})_z$ with $\Lambda = 2i\pi
 M_z^{\prime \vee}\oplus M_z$.  Then for any $\gamma \in 2i\pi M^{\prime
 \vee}$, $y(z)= \frac{1}{2i\pi}\int_{\gamma_z}\omega_z$ is a solution of
 $\mathcal{L}$ in $k[[z]]$.

 More generally, for any $\gamma^{\ast} \in \Lambda$, $\frac{1}{2i\pi}
 \int_{\gamma^{\ast}_z} \omega_z $ is a solution of the Picard-Fuchs
 differential equation of the form $y(z) \frac{\log az^n}{2i\pi} +
 y^{\ast}(z)$, with $y^{\ast}(z)\in k[[z]]$, $a\in k$, $n\in \Z$.  (In
 the case of the Legendre pencil, we have already met this situation
 with $y(z)=iF(\frac{1}{2},\frac{1}{2},1;z), y^{\ast}(z)= 2
 F^{\ast}(\frac{1}{2},\frac{1}{2},1;z), a= 2^{-8}, n = 2$).

 Let now fix an embedding $k\hookrightarrow K \subset \C_p$. For any $z
 \neq 0$ close enough to $0$, $M$ and $M^{\prime}$ identify respectively
 with the lattices $M_z$, $M_z^{\prime}$ attached to $\mathcal{A}_z$
 (\ref{para-degenerate}), and the choice of a determination of the
 $p$-adic logarithm identifies $\Lambda$ with $\Lambda_z$. It turns out
 that the $p$-adic pairing $\int^p$ of \ref{para-degenerate} extends to
 a horizontal pairing over a punctured disk around $0$, which is given
 by the ``same'' formula as in the complex case
 \begin{equation*}
 \frac{1}{(2i\pi)_p} \int^p_{\gamma^{\ast}_z}
 \omega_z = y(z)\frac{\log_p az^n}{(2i\pi)_p} + y^{\ast}(z),
 \end{equation*}
 where $y(z), y^{\ast}(z)$ are the \emph{same formal series as above,
 evaluated $p$-adically} (loc. cit.). The computation of
 \ref{para-padic-int-1} can be considered as a special case.

 \noindent This is also closely related to the work of T. Ichikawa on
 ``universal periods'' for Mumford curves
 \cite{sutorsamcoig}.
\end{para}

\begin{para}
 {\itshape Relation to the Fontaine-Messing periods.} Since the restriction
 of $\int^p$ to $(2i\pi)_p M_z^{\prime \vee}\subset
 T_p(\mathcal{A}_{z,\ovl{K}})$ is the Fontaine-Messing pairing, we
 obtain $(2i\pi)_p y(z)$ as a Fontaine-Messing period of
 $\mathcal{A}_z$. More generally, $y(z)\Log_p az^n +(2i\pi)_p
 y^{\ast}(z)$ appears as a period ($\Log$ instead of $\log$). We conclude
 that in this degenerating case, the Fontaine-Messing periods behave
 relatively well with respect to the Gauss-Manin connection. A similar
 observation applies to the pairing composed with $\theta$ (Hodge-Tate
 periods).
\end{para}

\subsection{$p$-adic Betti lattices for abelian varieties with
supersingular reduction.}\label{sub-p-adic-betti-lattice-supersing}

\begin{para}\label{para-ss-abelian}
 Any supersingular abelian variety over $\ovl{\F}_p$ is isogenous to
 a power of a supersingular elliptic curve. We refer to
 \cite[4]{waterhouse69:_abelian} for a detailed study of supersingular
 elliptic curves over finite fields (including the classification of
 isogeny classes and isomorphism classes, refining Deuring's classical
 work), and to \cite{li_o98:_modul} for the higher dimensional case. Let
 us simply recall a few basic facts:
 \begin{itemize}
  \item there is exactly one ${\ovl{\F}_p}$-isogeny class of
	supersingular elliptic curves.
  \item If $A_0$ is any supersingular elliptic curve over
	${\F}_{p^n}$, then $\bcD:= \End(A_0)_{\ovl{\F}_p}$ 
	\index{000D@$\bcD$}
	is a maximal order in a quaternion algebra over $\Q$, ramified
	exactly at $p,\infty$; we set $\bD=\bcD_{\Q}$.\index{000D@$\bD$}
  \item $A_0$ is ${\ovl{\F}_p}$-isomorphic (but not necessarily
	${\F}_{p^n}$-isomorphic) to an elliptic curve defined over
	${\F}_{p^2}$.
  \item If $n=1$, $\End(A_0)\otimes {\Q}={\Q}[Fr_{p}]$, an
	imaginary quadratic field; if $p\neq 2$, this is
	${\Q}[\sqrt{-p}]$. More precisely, the order $\End(A_0)$ can be either 
	${\Z}[Fr_{p}]$ or the maximal order in $\End(A_0)\otimes
	{\Q}$ (which coincide if $p\equiv 3\pmod{4}$). If $p=2$, there is
	another possibility, namely $\End(A_0)={\Z}[Fr_2]\simeq
	{\Z}[\sqrt{-1}]$.
  \item There exists a supersingular elliptic
	curve over ${\F}_p$ whose Frobenius endomorphism satisfies
	$(Fr_p)^2=-p.\id\,$ and whose endomorphism algebra is the maximal
	order in ${\Q}[Fr_{p}]$. Such elliptic curves belong to a single
	${\F}_p$-isogeny class.
  \item If $p\neq 2,3$, there is exactly one ${\F}_p$-isogeny class of
	supersingular elliptic curves over ${\F}_p$ (but several
	${\ovl{\F}_p}$-isomorphism classes in general).
 \end{itemize}
\end{para}

\begin{para}\label{para-supersing}
 Let $A_0$ be a supersingular abelian variety of 
 dimension $g$ over $\F_{p^n}$. By functoriality, the elements of 
 $\End\bigl((A_0)_{\ovl{\F}_p}\bigr)\otimes \Q\simeq M_g(\bD)$ act
 linearly on the crystalline cohomology 
 ${\mathrm{H}^{\ast}_{\mathrm{crys}}\bigl((A_0)_{\ovl{\F}_p}/\ovl{\Q}_p\bigr)}=
 \mathrm{H}^{\ast}_{\mathrm{crys}}\bigl((A_0)_{\ovl{\F}_p}/
 \widehat{\Z_p^{\mathrm{ur}}}\bigr)\otimes_{\widehat{\Z_p^{\mathrm{ur}}}}
 \ovl{\Q}_p$. 
 In degree one, this provides an embedding of $M_g(\bD)$ into the
 endomorphism ring of 
 $\mathrm{H}^1_{\mathrm{crys}}\bigl((A_0)_{\ovl{\F}_p}/\ovl{\Q}_p\bigr)$.

 On the other hand, one has a canonical isomorphism
 \begin{equation*}
  \bigwedge^2
  \mathrm{H}^1_{\mathrm{crys}}\bigl((A_0)_{\ovl{\F}_p}/\ovl{\Q}_p\bigr) \simeq
  \mathrm{H}^2_{\mathrm{crys}}\bigl((A_0)_{\ovl{\F}_p}/\ovl{\Q}_p\bigr) 
 \end{equation*} 
 and a canonical $\Q$-structure in $
  \mathrm{H}^2_{\mathrm{crys}}\bigl((A_0)_{\ovl{\F}_p}/\ovl{\Q}_p\bigr)
  \simeq 
  \ovl{\Q}_p^{g(2g-1)}$ coming from the fact that the whole cohomology
 in degree $2$ is generated by algebraic cycles (since
 $(A_0)_{\ovl{\F}_p}$ is isogenous to the $g$th power of a
 supersingular elliptic curve, one reduces easily to the case $g=2$, in
 which case this fact is well-known).

\medskip Let $F$ be either $\ovl{\Q}$ or a splitting number field for $\bD$: 
 $\bD\otimes_{\Q} F\simeq
 M_2(F)$.  In the latter case, this amounts to saying that $F$ is 
 totally imaginary, and for any place $v$ above $p$,
 $[F_v:\Q_p]$ is even.  We fix an embedding of $F$ in $\ovl{\Q}_p$.  Via 
 this embedding, the $M_g(\bD)$-action extends to a
 $F$-linear $M_{2g}(F)$-action on
 $\mathrm{H}^1_{\mathrm{crys}}\bigl((A_0)_{\ovl{\F}_p}/\ovl{\Q}_p\bigr)$.

 \noindent We consider an irreducible cyclic sub-$M_{2g}(F)$-module 
 $M_{2g}(F).u\subset \mathrm{H}^1_{\mathrm{crys}}\bigl((A_0)_{\ovl{\F}_p}/\ovl{\Q}_p\bigr)$.
 Obviously, such submodules exist, and have $F$-dimension $2g$.  We may and shall choose $u$ in such a 
 way that its exterior square of this $F$-space coincides with the canonical $F$-structure on $
\mathrm{H}^2_{\mathrm{crys}}\bigl((A_0)_{\ovl{\F}_p}/\ovl{\Q}_p\bigr) $ (this can be achieved by replacing $u$ by a
 suitable multiple).

 \begin{pro}\label{betti-la} Up to a homothety by a factor in 
 $\sqrt{F^\times}$, the normalized
 $M_{2g}(F)$-submodule $\;M_{2g}(F)\cdot u\subset
 \mathrm{H}^1_{\mathrm{crys}}\bigl((A_0)_{\ovl{\F}_p}/\ovl{\Q}_p\bigr)$
 depends only on $F\subset \ovl{\Q}_p$, not on the choice
 of $u$. 
  In particular, for $F=\ovl{\Q}$, this defines a canonical $\ovl{\Q}$-structure in $
\mathrm{H}^1_{\mathrm{crys}}\bigl((A_0)_{\ovl{\F}_p}/\ovl{\Q}_p\bigr)$, stable under 
$\End\bigl((A_0)_{\ovl{\F}_p}\bigr)$.
\end{pro}
 
 \begin{proof} Indeed, two such $M_{2g}(F)$-submodules are related 
 by some $h\in GL_{2g}(\ovl{\Q}_p)$, such that
 $\bigwedge^2(h)\in GL_{g(2g-1)}(F)$. Now 
 $M_{2g}(F)\cdot hu = h(M_{2g}(F)\cdot u)$ implies
 that $h$ normalizes
 $GL_{2g}(F)$. It follows that the image of $h$ in 
 $PGL_{2g}(\ovl{\Q}_p)$ lies in $PGL_{2g}(F)$.

  \noindent (When $g>1$, the proposition is not surprising, since
  $\bigwedge^2 V$ is a faithful representation of $PSL(V)$
  for any space $V$ of dimension $>1$).
 \end{proof}

\end{para}
\begin{rems}
 \begin{enumerate}
  \renewcommand{\theenumi}{\alph{enumi}}
  \item A minimal choice for $F$ is a splitting
	quadratic field $\Q(\sqrt{-d})\,$ (the splitting property amounts to
	saying that $d>0$ and $p$ ramifies  or remains prime in
	$\Q(\sqrt{-d})$. A natural choice is $d=p$).
  \item Let $\mathcal{O}_F$ be the ring of integers of $F$. Then
	$\End((A_0)_{\ovl{\F}_p})\otimes \mathcal{O}_F$ is an order in
	$M_{2g}(F)$, and there is a full
	$(\End((A_0)_{\ovl{\F}_p})\otimes \mathcal{O}_F)$-lattice in 
	$(M_{2g}(F)\cdot u)$. As an $\mathcal{O}_F$-module, it is
	projective of rank $2g$.  If $F=\ovl{\Q}$, \ie
	$\mathcal{O}_F=\ovl{\Z}$, then this module is 
	{\it canonically determined} by the
	normalization of its top exterior power, taking into account the 
	canonical isomorphism
	\begin{equation*}
	 \bigwedge^{2g}
	  \mathrm{H}^1_{\mathrm{crys}}\bigl((A_0)_{\ovl{\F}_p}/\ovl{\Z}_p\bigr)
	  \simeq
	  \mathrm{H}^{2g}_{\mathrm{crys}}
	  \bigl((A_0)_{\ovl{\F}_p}/\ovl{\Z}_p\bigr)
	  \simeq \ovl{\Z}_p,
	\end{equation*} 
	the latter isomorphism coming from the trace map. 
 
	Hence there is a 
	{\it canonical $\ovl{\Z}$-structure} in the
	crystalline cohomology of supersingular abelian 
	varieties (not used in the sequel).
 \end{enumerate}
\end{rems}

\begin{para}
 Let $A$ be an abelian variety defined over a 
 subfield $k\subset \C_p$, which has good
 {\it supersingular reduction} $A_0$ over the residue field of $k$. 
 There is a canonical embedding
 $\End(A)\hookrightarrow  \End((A_0)_{\ovl{\F}_p})$. 

The canonical $\ovl{\Q}$-subspace of $\; \mathrm{H}^1_{\mathrm{crys}}((A_0)_{\ovl{\F}_p}/\ovl{\Q}_p)\;$ defined in
\ref{betti-la} will be denoted by 
\begin{equation*}
  \mathrm{H}^1_{\mathrm{B}}(A_{\C_p}, \ovl{\Q}) \subset 
  \mathrm{H}^1_{\mathrm{crys}}\bigl((A_0)_{\ovl{\F}_p}/\ovl{\Q}_p\bigr) 
 \end{equation*}
because, as we shall see, it shares many properties with the Betti space
 $\mathrm{H}^1_{\mathrm{B}}(A_\C, \ovl{\Q})= \mathrm{H}^1_{\mathrm{B}}(A_\C,\Z)\otimes \ovl{\Q}$ of a
complex
 abelian variety $A_\C$). Its dual will be denoted by $ \mathrm{H}_{1,\mathrm{B}}(A_{\C_p}, \ovl{\Q}) $.

Similarly, we shall write $\mathrm{H}^1_{\mathrm{B}}(A_{\C_p}, F)$ for the $F$-subspace described \ref{betti-la}
(well-defined up to a homothety in $\sqrt{F^\times}$), and $ \mathrm{H}_{1,\mathrm{B}}(A_{\C_p},F) $ for its $F$-dual.
 
\medskip  Assume that $k\subset \ovl{\Q}_p$.  There is a functorial 
 isomorphism \cite{fadrc1}
 \begin{equation*}
  \mathrm{H}^1_{\mathrm{dR}}(A)\otimes_k\ovl{\Q}_p \simeq
 \mathrm{H}^1_{\mathrm{crys}}\bigl((A_0)_{\ovl{\F}_p}/\ovl{\Q}_p\bigr),
 \end{equation*} 
 from which one derives a canonical isomorphism
 \begin{equation*}
  \mathrm{H}^1_{\mathrm{dR}}(A)\otimes_k\ovl{\Q}_p \simeq 
  \mathrm{H}^1_{\mathrm{B}}(A_{\C_p},\ovl{\Q})\otimes_{\ovl{\Q}}\ovl{\Q}_p,
 \end{equation*} 
 which is functorial with respect to homomorphisms 
 of abelian varieties with supersingular reduction

 \noindent This isomorphism, which is far from being tautological, as we shall see, 
 can be translated into a $\ovl{\Q}_p$-valued pairing between 
 $\mathrm{H}_{1\mathrm{B}}(A_{\C_p}, F)$ and $\mathrm{H}^1_{\mathrm{dR}}(A)$. 
This pairing is conveniently expressed in the form of a {\it ``period 
 matrix''} $\Omega=(\omega_{ij})$ with
 entries in $\ovl{\Q}_p$, depending on the choice of a basis 
 $(\gamma_i)_{i=1,\ldots ,2g}$ of $\mathrm{H}_{1\mathrm{B}}(A_{\C_p}, F)$
 and a basis $(\omega_j)_{j=1,\ldots ,2g}$ of 
 $\mathrm{H}^1_{\mathrm{dR}}(A)$ (if $A$ is principally polarized, we
 choose symplectic bases). 

\noindent Warning: these $p$-adic periods attached to abelian varieties with supersingular reduction have little to
do with Fontaine-Messing periods  (a motivic interpretation of these periods is proposed in
\cite{tdmeigdvpdg}). 

\end{para}

\begin{para}
 {\itshape Horizontality.} If we let $A$ move in a family $A_z$,
 $\mathrm{H}^1_{\mathrm{crys}}\bigl((A_0)_{\ovl{\F}_p}/\ovl{\Q}_p\bigr)$ may be
 identified with a space of horizontal sections of the de Rham
 cohomology localized in the disk parameterizing the liftings $A_z$ of
 $(A_0)_{\ovl{\F}_p}$ \cite{fadrc1}. Thus the period matrix $\Omega(z)$ is a fundamental
 matrix of solutions of the Gauss-Manin connection.
\end{para}

In the sequel of this subsection, we assume for simplicity that 
 $g=1$, \ie $A$ is an elliptic curve.


\begin{para} {\itshape CM periods, $\,\Gamma_p$-values, and transcendence.}

 \medskip\noindent $\bullet$  Let $A_0$ be as before a
 supersingular elliptic curve over ${\F}_{p^n}$. Let $A$ be a CM-lifting of
 $A_0$, \ie an elliptic curve defined over some number field $k \subset
 \ovl{\Q}\subset \ovl{\Q}_p$ with residue field ${\F}_{p^n}$, such that
 $E= \End(A) \otimes \Q$ is
 quadratic. Note that $E$ is ipso facto a subfield of $\bD$. We denote by
 $c$ the conductor of the
 order $\End(A)$ (in particular, $ce \in \End(A)$).

 \medskip\noindent $\bullet$ Let $(\gamma_1,\gamma_2)$ be any symplectic
 basis of $\mathrm{H}_{1\mathrm{B}}(A,F)$. We can write
 $ce.\gamma_1=a\gamma_1+b\gamma_2$ with $a,b \in F$, so that
 \begin{equation*}
 \tau := \frac{\omega_{12}}{\omega_{11}}=
 \frac{-a+c\sqrt{-d}}{b}\in {\P}^1(FE).
 \end{equation*}
 In particular, $\tau \in {\P}^1({\ovl{\Q}})$.
 We conjecture the following converse, which is a $p$-adic analogue of
 T.~Schneider's theorem \cite{auei} (in the complex case):

 \begin{coj}\label{coj-complex-mult}
  Assume that $A$ is a lifting of $A_0$ defined over $\ovl{\Q}$, and
  that $\tau = \frac{\omega_{12}}{\omega_{11}} \in {\P}^1({\ovl{\Q}})$.
  Then $A$ has complex multiplication.
 \end{coj} 

 
\medskip\noindent $\bullet$ 
 Let $E= {\Q}[\sqrt{-d}]$ and 
 $E'= {\Q}[\sqrt{-d'}]$ be two distinct subfields of
 $\bD\;$ ($d, d'$
 squarefree; we do not exclude the case $d=d'$). According to
 \cite[II 1.4]{adadq}, $E\cap \bcD=\OO_E,\;E'\cap \bcD=\OO_{E'}$. The
 elements  $e = \sqrt{-d}
 \in E$ and $e'=\sqrt{-d'} \in  E'$ generate the
 algebra $\bD$. We remark
 that
 $ee'+e'e$ commutes to $e$ and $e'$, hence is an
 integer $m \in \Z$ (an
 even integer if $d$ or $d'\equiv 3\pmod{4}$, which is divisible by
 $4$ if both $d,
 d'\equiv 3\pmod{4}$). We have
 $(ee'-e'e)^2=m^2-4dd'$. Since
 $ee'-e'e$ anticommutes with $e$, it is not an integer,
 hence $\vert m\vert <
 2\sqrt{dd'}$. On the other hand, the images of $e$ and $e'$
 in $\bcD\otimes
 {\F}_p$ commute ($ee'-e'e$ acts trivially on regular
 differentials in
 characteristic $p$), therefore $p$ divides $4dd'-m^2$. 
 
 \medskip\noindent $\bullet$ On the other hand, there exists a unique symplectic basis
 of eigenvectors for $e$ in $\mathrm{H}^1_{\mathrm{dR}}(A)$ of the form
 $(\omega_1,\omega_2+\sigma \omega_1)$, with
 $\sigma \in k$. We get the relation
 \begin{equation*}
  \omega_{22}+\sigma\omega_{12} = {\til{\tau}}
   (\omega_{21} + \sigma\omega_{11}), \text{ with } {\til{\tau}} =
   \frac{-a-c\sqrt{-d}}{b}.
 \end{equation*} 
 Let now $E'={\Q}[\sqrt{-d'}]\neq E$ be another subfield of $\bD$,
 and assume that $(\gamma_1,\gamma_2)$ is a basis of eigenvectors for
 $E'$: $e'.\gamma_1=\sqrt{-d'}\gamma_1$. We get
 $(ee^{\prime}+e^{\prime}e)\gamma_1=\frac{2a\sqrt{-d^{\prime}}}{c}\gamma_1
 = m\gamma_1$, so that
 \begin{equation*}
 \frac{\til{\tau}}{\tau} =
 \frac{m+2 \sqrt{dd^{\prime}}}{m-2\sqrt{dd^{\prime}}}
 \in {\Q}[\sqrt{dd^{\prime}}]
 \end{equation*}
  the sign of the square root being chosen in such a way that $p \mid
 {2\sqrt{dd'} -m}$ (in the unramified case).

 \medskip\noindent $\bullet$ The previous relations, together with $\det\Omega =
 1$, show that the transcendence degree over $\ovl{\Q}$ of the entries
 of $\Omega$ is at most one in the CM case. 

 \medskip\noindent $\bullet$ We now describe the
 $p$-adic number $\omega_{11}$ modulo $\ovl{\Q}^{\times}$, keeping the
 notation of \ref{para-lerch-chowla} and \thmref{thm-ogus-formula}
 ($\epsilon$ is the Dirichlet 
 character, $w$ the number of roots of unity, $h$ the class number).

 \begin{thm}\label{thm-cm-period}
  In case $p$ does not ramify in the field of
  complex multiplications ${\Q}[\sqrt{-d}] = \End(A)\otimes{\Q}$, one has
  \begin{equation*}
   \omega_{11}\; \sim \prod_{u\in ({\Z}/d)^{\times}}
    \Bigl(\Gamma_p\Lan{p\frac{u}{d}}\Bigr)^{-\epsilon(u)w/8h} \;\text{ in }\;
    \ovl{\Q}_p^{\times}/\ovl{\Q}^{\times}.
  \end{equation*}
   \end{thm}

\medskip \noindent {\it Examples.} If $A$ is the Legendre curve $X_{1/2}\;$
  ($y^2=x(x-1)(x-1/2)$)
  and $p\equiv 3\pmod{4}$, we find $\omega_{11} \sim \Gamma_p(1/4) .$
  Similarly, if $\End(A)
  \otimes {\Q} = {\Q}(\sqrt{-3})$ and $p\equiv 2\pmod{3}$, then
  $\omega_{11} \sim
  \Gamma_p(1/3)^{3/2}.$

 \begin{proof}
  We may assume that the number field $k$ ($\subset \ovl{\Q}_p$) is
  Galois over $\Q$. We denote by $k'=k\cap {\Q}_p$ the subfield of $k$
  fixed by the local Galois group $\Gal({\ovl{\Q}}_p/{\Q}_p)$. Let us
  consider the Weil restriction $B:= \mathcal{R}_{k/ k'}(A)$
  (\cf \cite[7.6]{bosch90:_nm}). This is an abelian variety
  over $k' \subset {\Q}_p$, with good, supersingular reduction at $p$
  (like $A$), and we have $B_k\simeq \prod_{\tau \in \Gal(k/k')}A^{\tau}$.

  \noindent On the other hand, let $\psi$ be any element of degree one
  in the Weil group (\ie any lifting of the Frobenius element in
  $\Gal({\ovl{\Q}}_p/{\Q}_p)$), and let 
  ${\Psi}_{A^{\tau}}$ (\resp  ${\Psi}_B$) be the corresponding
  semilinear endomorphism of $\mathrm{H}^1_{\mathrm{dR}}(A^{\tau}/k)\otimes
  {\ovl{\Q}}_p$ (\resp $\mathrm{H}^1_{\mathrm{dR}}(B/k')\otimes {\ovl{\Q}}_p$).
  If ${\Phi}_B$ denotes the canonical linear action induced by
  ${Fr_{{\F}_p}}$ on
  $\mathrm{H}^1_{\mathrm{crys}}(B_0/{\Q}_p)\otimes{\ovl{\Q}_p}\simeq
  \mathrm{H}^1_{\mathrm{dR}}(B)_{\ovl{\Q}_p}$, one has the formula ${\Psi}_B =
  {\Phi}_B \circ (\id \otimes \psi)$ (\cf \cite[4]{fadrc1}).  Moreover, ${\Psi}_B=\bigoplus {\Psi}_{A^{\tau}}$ with
  respect to the decomposition
  $\mathrm{H}^1_{\mathrm{dR}}(B,{\ovl{\Q}_p})\simeq
  \bigoplus_{\tau} \mathrm{H}^1_{\mathrm{dR}}(A^{\tau},{\ovl{\Q}_p})$. 
 
  \noindent If $(\omega_1,\omega_2)$ is any symplectic basis of
  eigenvectors for ${\Q}[\sqrt{-d}]$ in
  $\mathrm{H}^1_{\mathrm{dR}}(A)_{\ovl{\Q}}$,
  with $\omega_1 \in \Omega^1(A)_{\ovl{\Q}}$, we
  can write  ${\Psi}_A(\omega_1)= \kappa.\omega_2$,
  where $\kappa$ modulo $\ovl{\Q}^{\times}$ is given by Ogus' formula in
  \thmref{thm-ogus-formula}
  \begin{equation*}
   \kappa\; \sim \prod_{u\in ({\Z}/d)^{\times}}
    \Bigl(\Gamma_p\Lan{p\frac{u}{d}}\Bigr)^{-\epsilon(u)w/4h} \;\text{ in }\;
    \ovl{\Q}_p^{\times}/\ovl{\Q}^{\times}.
  \end{equation*}
   Because the conjugates $A^{\tau}$ are isogenous to each other over $k$ (and
  since isogenies preserve
  the regular differential forms), we derive that the constant $\kappa$, resp.
  the ``period''
  $\omega_{11}$, is the same modulo $\ovl{\Q}^{\times}$ for each of them.
  Therefore, for any
  $\omega \in \Omega^1(B) \subset
  \mathrm{H}^1_{\mathrm{crys}}(B_0/{\Q}_p)$, we can write
  ${\Phi}_B(\omega)={\Psi}_B(\omega)= \kappa.\eta$ for some $\eta \in
  \mathrm{H}^1_{\mathrm{dR}}(B)_{\ovl{\Q}}$ whose pairing with any
  element of $\mathrm{H}_{1\mathrm{B}}(B,{\ovl{\Q}})$ belongs to
  ${\omega_{21}}.{\ovl{\Q}}$. Since $\omega_{11}^{-1}\omega\in
  \mathrm{H}^1_{\mathrm{B}}(B,{\ovl{\Q}}),\,\omega_{21}^{-1}
  \eta\in \mathrm{H}^1_{\mathrm{B}}(B,{\ovl{\Q}})$, and since
  $\Phi_B$ respects
  $\mathrm{H}_{1\mathrm{B}}(B,{\ovl{\Q}})$, we get
  $\kappa \sim \frac{\omega_{11}}{\omega_{21}} \sim \omega_{11}^2$,
  whence the result. 
 \end{proof}
 
\begin{coj}\label{coj-no-complex-multi}
Assume that $A$ is a lifting of $A_0$ defined over $\ovl{\Q}$, and that
$p$ does not ramify $\End(A)\otimes \Q$ (which is obviously the case if
$A$ does not have complex multiplication). Then
$\mathrm{H}^1_{\mathrm{dR}}(B)_{\ovl{\Q}}\neq
\mathrm{H}^1_{1\mathrm{B}}(B,{\ovl{\Q}})$.
\end{coj} 

\noindent In the absence of complex multiplication, this follows from the previous conjecture.
In the presence of complex multiplication, it amounts to the transcendence of $\omega_{11}$. 

 \noindent Taking into account \ref{thm-cm-period}, the conjecture would
 imply, for instance, the transcendence of the
 adic numbers $\Gamma_3(1/4)$, $\Gamma_7(1/4)$, $\Gamma_2(1/3)$,
 $\Gamma_5(1/3)$ (note, in contrast,
 that $\Gamma_5(1/4)$ and $\Gamma_7(1/3)$ are algebraic numbers,
 according to the Gross-Koblitz formula (\ref{para-gross-koblitz})).

 
 \noindent $\bullet$ Actually, it may happen, in the ramified CM case, that
 $\mathrm{H}^1_{\mathrm{dR}}(A)_{\ovl{\Q}}=
 \mathrm{H}^1_{1\mathrm{B}}(A,{\ovl{\Q}})$. 

 \noindent Let us first notice that the argument given in the proof of theorem
 \ref{thm-cm-period} still works in the ramified case, and allows to
 conclude that $\kappa \sim  \omega_{11}^2$. The computation of $\kappa
 \pmod{{\ovl{\Q}}^\times}$ (for $p$ odd) according to
 Ogus-Coleman-Urfels was explained in \ref{para-ramif}. In the example
 of an elliptic curve with complex multiplication by $\Q[\sqrt{-3}]$, we
 have seen that $\kappa\in \ovl{\Q}$, hence $\omega_{11}\in \ovl{\Q}$
 and $\mathrm{H}^1_{\mathrm{dR}}(A)_{\ovl{\Q}}= 
 \mathrm{H}^1_{1\mathrm{B}}(A,{\ovl{\Q}})$. 

This is also the case for an elliptic curve with complex multiplication
 by $\Q[\sqrt{-3n}]$ 
if the Legendre symbol $(\frac{n}{3})$ is $1$, but probably not in general if
$(\frac{n}{3})=-1$ (although this happens for $n=8$, after the last example in
\ref{para-ramif}).
\end{para}

\begin{para}\label{para-supersingular-disk}
 {\it $L_{\frac{1}{2},\frac{1}{2},1}$ in a supersingular disk.} For
 concreteness, we consider the special case of the supersingular
 Legendre elliptic curve with parameter $z = 1/2$ in characteristic
 $p\equiv 3\pmod{4}$. A basis of solutions of the hypergeometric
 differential operator 
 $L_{\frac{1}{2},\frac{1}{2},1}= HGDO(\frac{1}{2},\frac{1}{2},1)$ in the
 supersingular disk $\matheur{D}(\FRAC{1}{2}, 1^-)$ is given by
 \begin{equation*}
 F\Bigl(\frac{1}{4},\frac{1}{4},\frac{1}{2};(1-2z)^2\Bigr),\;\;
 (1-2z)F\Bigl(\frac{3}{4},\frac{3}{4},\frac{3}{2};(1-2z)^2\Bigr).
 \end{equation*}

\begin{itemize}
 \item Let us consider the symplectic basis
       $\omega_1=[\frac{dx}{2y}],\omega_2=[\frac{(2x-1)dx}{4y}]$ of
       $\mathcal{M}$ (de Rham cohomology). At $z=1/2$, this is a basis of
       eigenvectors for the action of $E':=\End(X_{1/2}) \otimes {\Q}
       ={\Q}(\sqrt{-1})$. The Gauss-Manin connection satisfies
       \begin{equation*}
	\omega_2=2z(z-1)\nabla\Bigl(\frac{d}{dz}\Bigr)\omega_1
	 + \frac{4z-5}{6}\omega_1.
       \end{equation*}
       The fundamental solution matrix $Y$ of
       the Gauss-Manin connection (expressed in the basis
       $(\omega_1,\omega_2)$), normalized by $Y(1/2)=\id$, is then given by
       \begin{align*}
	& y_{11} =
	F\Bigl(\frac{1}{4},\frac{1}{4},\frac{1}{2};(1-2z)^2\Bigr) +
	\frac{1}{2}(1-2z)F\Bigl(\frac{3}{4},\frac{3}{4},\frac{3}{2};
	(1-2z)^2\Bigr),\\
	& y_{12} = (1-2z)F\Bigl(\frac{3}{4},\frac{3}{4},\frac{3}{2};
	(1-2z)^2\Bigr),\\
	& y_{2i} = 2z(z-1) \frac{dy_{1i}}{dz}+\frac{4z-5}{6} y_{1i},
	\quad i=1,2.
       \end{align*}

 \item We choose a symplectic basis $(\gamma_1,\gamma_2)$ of
       eigenvectors for $E'$ in $\mathrm{H}^1_{1\mathrm{B}}(X_{1/2},F)$
       (note that ${\Q}(\sqrt{-1})$ is always contained in $F$). We have
       $\mathrm{H}^1_{1\mathrm{B}}(X_{1/2},F) =
       \mathrm{H}^1_{1\mathrm{B}}(X_z,F)$ for any point $z$ in the
       supersingular disk $\matheur{D}(\FRAC{1}{2}, 1^-)$.
       Because the period matrix $\Omega(z)$ is another solution matrix
       of the Gauss-Manin connection, we have the relation
       $Y(z)=\Omega(z).\Omega(1/2)^{-1}$.
       
       Let us consider a CM-point $\zeta$ in the supersingular
       disk $\matheur{D}(\FRAC{1}{2}, 1^-)$: $\End(X_{\zeta}) =
       E ={\Q}(\sqrt{-d})$ (and $E\neq E'$ in $\bD$).

       \begin{figure}[h]%PICTURE10
	\begin{picture}(150,120)(0,10)
	 \put(50,0){\includegraphics{fig10.eps}}
	 \put(55,118){$X_{\FRAC{1}{2}}$}
	 \put(120,118){$X_{\zeta}$}
	 \put(0,40){$\matheur{D}(\FRAC{1}{2}, 1^-)$}
	 \put(22,95){$\scriptstyle\Q(\sqrt{-1})$}
	 \put(142,95){$\scriptstyle\Q(\sqrt{-d})$}
	 \put(75,50){$\scriptstyle\FRAC{1}{2}$}
	 \put(114,51){$\scriptstyle\zeta$}
	\end{picture}
	\caption{}
       \end{figure} 
              
       \noindent Let $\sigma$ and $m$ be as before (in the present case, the
       integer $m$ is even:
       $m=2n$). Our period relations can be easily expressed in terms of
       $Y(\zeta)$. Apart from the
       obvious relation $\det Y(\zeta)=1$, we find a relation
       \footnote{F. Beukers has also found these
       relations independently, by a different
       method.}
       \begin{equation*}
	\left.\frac{{y_{11}(y_{22}+\sigma y_{12})}}{{y_{12}(y_{21}+\sigma
	 y_{11})}}\right|_{z=\zeta} = \frac{\til{\tau}}{\tau} =
	\frac{n+\sqrt{d}}{n-\sqrt{d}}.
       \end{equation*}

       \noindent Changing the viewpoint, one can consider $\zeta$ as
       fixed, and vary the prime $p\neq 2$ (or more precisely the place
       of ${\Q}(\zeta,\sqrt{-1},\sqrt{-d}$)).  This relation
       between values of $p$-adic hypergeometric functions at
       $\zeta$ holds whenever $\vert \zeta - 1/2 \vert_p < 1$; it
       depends on $p$ via the integer $n$ ($\vert n\vert < \sqrt{d}, \;
       p \mid d - n^2$). A relation of the same kind holds at the places
       at infinity (derived along similar lines, using the usual Betti
       lattices).  As a specific example, we can take $p=3, \zeta=$ a
       primitive $6$th root of unity.  Then $X_{\zeta}$ has complex
       multiplication by ${\Z}[\zeta]$, $\;\sigma =
       \frac{-\zeta^2}{2(1+\zeta)}\;, \;n=0$ and
       $\frac{\til{\tau}}{\tau}=-1$.
 \item On the other hand, one can combine \thmref{thm-cm-period} and
       the formula $Y(\zeta)=\Omega(\zeta).\Omega(1/2)^{-1}$ in order to
       express the value $\bmod{\,\ovl{\Q}^{\times}}$ of the $p$-adic
       hypergeometric functions $y_{ij}$ in terms of $\Gamma_p$ ({\it
       supersingular avatar of Young's formulas}).
       
       \noindent For instance, if $p=7, \zeta = 2(\sqrt{2} -1)$ (complex
       multiplication by ${\Z}[\sqrt{-2}]$), one has $\omega_{11}(\zeta) \sim
       (\Gamma_7(1/8)\Gamma_7(3/8))^{1/2},\,\omega_{11}(1/2) \sim
       \Gamma_7(1/4) $, from which
       one derives the 7-adic evaluation
       \begin{align*}
	F\left(\frac{3}{4},\frac{3}{4},\frac{3}{2};5-4\sqrt{2}\right) & \sim
	 \bigl(\Gamma_7(1/8)\Gamma_7(3/8)\bigr)^{1/2}\Gamma_7(1/4)^{-1}.
       \end{align*}
\end{itemize}
\end{para}

\subsection{Conclusion and vista.}
We have seen many instances of the following scenario: some familiar
notion or phenomenon from the complex realm shows itself in two
different aspects; each of these aspects has a natural $p$-adic
counterpart; these counterparts are not complementary aspects of the
same $p$-adic entity, but belong to totally different theories.  

In II, we shall see that the theory of $p$-adic period mappings is a privileged
field where this semantic splitting process integrates into an
harmonious picture, where the complementarity between the $p$-adic
counterparts is restored at a deeper level.

The $p$-adic period mappings which we shall deal with relate some
deformation spaces of $p$-divisible groups to certain grassmannians. It
will turn out that they can be described by quotients of analytic
solutions of the Gauss-Manin connection, just as in the complex case. In
particular, at the {\it modular level}, the Fontaine-Messing periods
(which links up directly the \'etale and crystalline representations of
$p$-divisible groups) go offstage or remain at the background.

On the other hand, we shall see that the differential equations
studied in section 1 (\resp 2) arise in the context of deformations of
supersingular (resp. ordinary) $p$-divisible groups.
